% Amélioration de la santé du système immunitaire — SVT 1ère A — Cours signé OnBuch+
% Programme officiel MINESEC. Compiler avec : tectonic NA04-amelioration-sante-systeme-immunitaire.tex
\newcommand{\DOCMATIERE}{SVT}
\newcommand{\DOCNIVEAU}{1ère A}
\newcommand{\DOCMODULE}{Module 2 : L'Éducation à la Santé}
\newcommand{\DOCLECON}{NA04}
\newcommand{\DOCTITRE}{Amélioration de la santé du système immunitaire}
% OnBuch+ — préambule « cours signés » ENRICHI (compilé avec tectonic / XeLaTeX)
% Système visuel « Pop » d'OnBuch+ : fond crème (jamais blanc), bordures encre
% épaisses, ombres « dures » décalées, titres Archivo Black en capitales.
% Version MATHÉMATIQUES : graphes (pgfplots/tikz), boîtes pédagogiques
% (définition, propriété, méthode, attention, le savais-tu, exercice résolu,
% exercice, TP, bilan), page de garde et signature OnBuch+ ancrée en bas.
%
% Macros attendues dans main.tex AVANT \input{preamble} :
%   \DOCMATIERE  \DOCNIVEAU  \DOCMODULE  \DOCLECON (numéro)  \DOCTITRE
\documentclass[11pt]{article}

\usepackage{fontspec}
\usepackage[french]{babel}
\usepackage[a4paper,margin=2cm,top=3.4cm,bottom=2.8cm,headheight=1.4cm,footskip=1.3cm]{geometry}
\usepackage{amsmath,amssymb,amsfonts}
\usepackage{mathtools} % \xmapsto, \coloneqq... souvent utilisés par les modèles
\usepackage{cancel} % simplifications barrées (bilans de forces)
% \abs{z} (module, valeur absolue) et \norm{u} (norme), utilisés par les modèles
\providecommand{\abs}{}\let\abs\relax\DeclarePairedDelimiter{\abs}{\lvert}{\rvert}
\providecommand{\norm}{}\let\norm\relax\DeclarePairedDelimiter{\norm}{\lVert}{\rVert}
\usepackage[table]{xcolor} % [table] : fournit aussi \rowcolors (alternance de lignes)
\usepackage{pagecolor}
\usepackage{tikz}
\usetikzlibrary{arrows.meta,positioning,calc,shapes.geometric,shapes.misc,shapes.arrows,decorations.pathmorphing,decorations.pathreplacing,decorations.markings,patterns,shadows,mindmap,trees,fit,backgrounds,matrix,intersections,angles,quotes}
\usepackage{pgfplots}
\usepackage[version=4]{mhchem} % équations nucléaires \ce{^{235}_{92}U}
\usepackage[european,siunitx]{circuitikz} % schémas électriques
\pgfplotsset{compat=1.18}
\usepgfplotslibrary{fillbetween,groupplots} % aires entre courbes (calcul intégral)
\usepackage{enumitem}
\usepackage{fancyhdr}
% [most] charge toutes les bibliothèques tcolorbox (listings, minted, raster,
% documentation...) et gonfle inutilement la mémoire TeX sur un document long
% (~300 boîtes) : on ne charge que ce qui est réellement utilisé.
\usepackage[skins,breakable]{tcolorbox}
\usepackage{graphicx}
\usepackage{colortbl}
\usepackage{booktabs}
\usepackage{tabularx}
\usepackage{diagbox} % case d'en-tête diagonale des tableaux à double entrée
% Colonnes extensibles centrée / gauche / droite (C, L, R), souvent utilisées par les modèles
\newcolumntype{C}{>{\centering\arraybackslash}X}
\newcolumntype{L}{>{\raggedright\arraybackslash}X}
\newcolumntype{R}{>{\raggedleft\arraybackslash}X}
\usepackage{array}
\usepackage{multicol}
\usepackage{siunitx}
% parse-numbers=false : accepte aussi \SI{2,5 \times 10^{-3}}{...}
\sisetup{locale=FR,output-decimal-marker={,},group-minimum-digits=4,parse-numbers=false}
\DeclareSIUnit\ohm{\ensuremath{\symup{\Omega}}} % U+2126 absent de la police
\DeclareSIUnit\division{div} % réglages d'oscilloscope (ms/div, V/div)
\DeclareSIUnit\tr{tr} % tours (vitesse de rotation, ex. \SI{1500}{\tr\per\minute})
\DeclareSIUnit\ch{ch} % cheval-vapeur (puissance moteur, unité usuelle hors SI)
\DeclareSIUnit\franccfa{FCFA} % franc CFA (coûts, contexte camerounais)
\DeclareSIUnit\dioptre{\ensuremath{\delta}} % vergence d'une lentille (dioptrie)
\usepackage{qrcode}
% \degree : commande fréquemment utilisée par les modèles pour un angle ou une
% température en dehors de \SI{}{} (non fournie par les paquets chargés).
\providecommand{\degree}{^{\circ}}
% \qed (fin de démonstration), utilisé par les modèles sans amsthm
\providecommand{\qed}{\hfill$\square$}
% \barre : double barre des valeurs interdites dans un tableau de variations (tkz-tab)
\providecommand{\barre}{\ensuremath{\|}}
% environnement proof (amsthm non chargé) utilisé par les modèles
\ifdefined\proof\else\newenvironment{proof}[1][Démonstration]{\par\noindent\textit{#1.}\ }{\qed\par}\fi
\usepackage{lastpage}
\usepackage{hyperref}
\hypersetup{hidelinks}

% ============================================================
% Palette « Pop » OnBuch+ — mêmes hex que la classe P (plus_ui.dart)
% ============================================================
\definecolor{poporange}{HTML}{F59321}
\definecolor{popdark}{HTML}{E07A0C}
\definecolor{popcream}{HTML}{FFF6E8}
\definecolor{popink}{HTML}{1C1714}
\definecolor{poppurple}{HTML}{7A5AE0}
\definecolor{popgreen}{HTML}{2E9E6A}
\definecolor{poppink}{HTML}{E0457B}
\definecolor{popblue}{HTML}{2F7DE1}
\definecolor{popgold}{HTML}{C98A12}
\definecolor{popmuted}{HTML}{8A7F76}
\definecolor{popline}{HTML}{EADFD2}
\definecolor{popgray}{HTML}{8A7F76}
\colorlet{popgrey}{popgray}
% Alias tolérants (le modèle nomme parfois les couleurs différemment)
\colorlet{popwhite}{white}
\colorlet{popblack}{popink}
\colorlet{popred}{poppink}
\colorlet{popyellow}{popgold}
% Teintes claires pour fonds de tableaux / zones
\colorlet{popblueL}{popblue!10!white}
\colorlet{poporangeL}{poporange!14!white}
\colorlet{popgreenL}{popgreen!12!white}
\colorlet{poppurpleL}{poppurple!12!white}
\colorlet{poppinkL}{poppink!12!white}
\colorlet{popgoldL}{popgold!14!white}

\pagecolor{popcream}

% ============================================================
% Typographie
% ============================================================
\defaultfontfeatures{Ligatures=TeX}
\setmainfont{PlusJakartaSans-Medium.ttf}[Path=../../../fonts/,BoldFont=PlusJakartaSans-Bold.ttf]
% Maths en sans-serif (Fira Math) avec chiffres et lettres droites de Plus
% Jakarta, pour que les formules se fondent dans le texte.
\usepackage[math-style=ISO,bold-style=ISO]{unicode-math}
\setmathfont{FiraMath-Regular.otf}[Path=../../../fonts/]
\setmathfont{PlusJakartaSans-Medium.ttf}[Path=../../../fonts/,range={up/num,up/{Latin,latin},\mathup}]
\usepackage{icomma} % virgule décimale sans espace en mode math
% Caractères mathématiques Unicode absents des polices de texte (Plus Jakarta,
% Archivo Black) : rendus par leur équivalent mathématique, en texte comme en
% formule — sinon ils s'affichent en carré vide (ex. « Arithmétique dans ℤ »).
\usepackage{newunicodechar}
\newunicodechar{ℝ}{\ensuremath{\mathbb{R}}}
\newunicodechar{ℤ}{\ensuremath{\mathbb{Z}}}
\newunicodechar{ℂ}{\ensuremath{\mathbb{C}}}
\newunicodechar{ℕ}{\ensuremath{\mathbb{N}}}
\newunicodechar{ℚ}{\ensuremath{\mathbb{Q}}}
\newunicodechar{𝔻}{\ensuremath{\mathbb{D}}}
\newunicodechar{α}{\ensuremath{\alpha}}
\newunicodechar{β}{\ensuremath{\beta}}
\newunicodechar{γ}{\ensuremath{\gamma}}
\newunicodechar{δ}{\ensuremath{\delta}}
\newunicodechar{ε}{\ensuremath{\varepsilon}}
\newunicodechar{θ}{\ensuremath{\theta}}
\newunicodechar{λ}{\ensuremath{\lambda}}
\newunicodechar{μ}{\ensuremath{\mu}}
\newunicodechar{σ}{\ensuremath{\sigma}}
\newunicodechar{τ}{\ensuremath{\tau}}
\newunicodechar{φ}{\ensuremath{\varphi}}
\newunicodechar{ψ}{\ensuremath{\psi}}
\newunicodechar{ω}{\ensuremath{\omega}}
\newunicodechar{Σ}{\ensuremath{\Sigma}}
% majuscules produites par \MakeUppercase dans les titres (α-aminés -> Α)
\newunicodechar{Α}{\ensuremath{\alpha}}
\newunicodechar{Β}{\ensuremath{\beta}}
\newunicodechar{Γ}{\ensuremath{\Gamma}}
\newunicodechar{Δ}{\ensuremath{\Delta}}
\newunicodechar{Ε}{\ensuremath{\varepsilon}}
\newunicodechar{Θ}{\ensuremath{\Theta}}
\newunicodechar{Λ}{\ensuremath{\Lambda}}
\newunicodechar{Μ}{\ensuremath{\mu}}
\newunicodechar{Τ}{\ensuremath{\tau}}
\newunicodechar{Φ}{\ensuremath{\Phi}}
\newunicodechar{Ψ}{\ensuremath{\Psi}}
\newunicodechar{Ω}{\ensuremath{\Omega}}
\newunicodechar{↦}{\ensuremath{\mapsto}}
\newunicodechar{∈}{\ensuremath{\in}}
\newunicodechar{∉}{\ensuremath{\notin}}
\newunicodechar{∧}{\ensuremath{\wedge}}
\newunicodechar{⇔}{\ensuremath{\Leftrightarrow}}
\newunicodechar{⟺}{\ensuremath{\Longleftrightarrow}}
\newunicodechar{⇒}{\ensuremath{\Rightarrow}}
\newunicodechar{⟹}{\ensuremath{\Longrightarrow}}
\newunicodechar{∘}{\ensuremath{\circ}}
\newunicodechar{∪}{\ensuremath{\cup}}
\newunicodechar{∩}{\ensuremath{\cap}}
\newunicodechar{≡}{\ensuremath{\equiv}}
\newunicodechar{⊂}{\ensuremath{\subset}}
\newunicodechar{⊥}{\ensuremath{\perp}}
\newunicodechar{∃}{\ensuremath{\exists}}
\newunicodechar{∀}{\ensuremath{\forall}}
\newunicodechar{∛}{\ensuremath{\sqrt[3]{\,}}}
\newunicodechar{∜}{\ensuremath{\sqrt[4]{\,}}}
\newunicodechar{⃗}{\ensuremath{{}^{\to}}}
\newunicodechar{ⁿ}{\ifmmode^{n}\else\textsuperscript{\ensuremath{n}}\fi}
\newunicodechar{ˣ}{\ifmmode^{x}\else\textsuperscript{\ensuremath{x}}\fi}
\newunicodechar{ᵇ}{\ifmmode^{b}\else\textsuperscript{\ensuremath{b}}\fi}
\newunicodechar{ʳ}{\ifmmode^{r}\else\textsuperscript{\ensuremath{r}}\fi}
\newunicodechar{ᵘ}{\ifmmode^{u}\else\textsuperscript{\ensuremath{u}}\fi}
\newunicodechar{ʸ}{\ifmmode^{y}\else\textsuperscript{\ensuremath{y}}\fi}
\newunicodechar{ᵃ}{\ifmmode^{a}\else\textsuperscript{\ensuremath{a}}\fi}
\newunicodechar{ᵉ}{\ifmmode^{e}\else\textsuperscript{\ensuremath{e}}\fi}
\newunicodechar{ᵏ}{\ifmmode^{k}\else\textsuperscript{\ensuremath{k}}\fi}
\newunicodechar{ᵖ}{\ifmmode^{p}\else\textsuperscript{\ensuremath{p}}\fi}
\newunicodechar{ᴺ}{\ifmmode^{N}\else\textsuperscript{\ensuremath{N}}\fi}
\newunicodechar{ᶜ}{\ifmmode^{c}\else\textsuperscript{\ensuremath{c}}\fi}
\newunicodechar{ʰ}{\ifmmode^{h}\else\textsuperscript{\ensuremath{h}}\fi}
\newunicodechar{ᵠ}{\ifmmode^{q}\else\textsuperscript{\ensuremath{q}}\fi}
\newunicodechar{ₙ}{\ifmmode_{n}\else\textsubscript{\ensuremath{n}}\fi}
\newunicodechar{ₐ}{\ifmmode_{a}\else\textsubscript{\ensuremath{a}}\fi}
\newunicodechar{ᵢ}{\ifmmode_{i}\else\textsubscript{\ensuremath{i}}\fi}
\newunicodechar{ᵣ}{\ifmmode_{r}\else\textsubscript{\ensuremath{r}}\fi}
\newunicodechar{ₖ}{\ifmmode_{k}\else\textsubscript{\ensuremath{k}}\fi}
\newunicodechar{ᵦ}{\ifmmode_{\beta}\else\textsubscript{\ensuremath{\beta}}\fi}
\newunicodechar{ₓ}{\ifmmode_{x}\else\textsubscript{\ensuremath{x}}\fi}
\newfontfamily\popdisplay{ArchivoBlack-Regular.ttf}[Path=../../../fonts/]
\newfontfamily\poplogo{SpaceGrotesk-Bold.ttf}[Path=../../../fonts/]
\newfontfamily\popsemi{PlusJakartaSans-SemiBold.ttf}[Path=../../../fonts/]
\newfontfamily\popxbold{PlusJakartaSans-ExtraBold.ttf}[Path=../../../fonts/]
\newfontfamily\popxboldls{PlusJakartaSans-ExtraBold.ttf}[Path=../../../fonts/,LetterSpace=3.0]

\setlist[enumerate]{leftmargin=1.7em,itemsep=5pt,topsep=4pt}
\setlist[itemize]{leftmargin=1.7em,itemsep=4pt,topsep=4pt,label={\color{poporange}\textbullet}}
\setlength{\parindent}{0pt}
\setlength{\parskip}{5pt}
\renewcommand{\arraystretch}{1.25}

% pgfplots : style Pop par défaut
\pgfplotsset{
  every axis/.append style={axis line style={popink,line width=1pt},tick style={popink},
    grid style={popline,line width=0.5pt},label style={font=\small},tick label style={font=\footnotesize}},
  popcurve/.style={poporange,line width=1.8pt,no markers,smooth},
}
% Même style disponible dans un \draw TikZ simple (hors pgfplots)
\tikzset{popcurve/.style={poporange,line width=1.8pt}}
\tikzset{popgrid/.style={popline,very thin}}
\colorlet{popcurve}{poporange} % le nom du style est parfois utilisé comme couleur

% ============================================================
% Pill / squiggle
% ============================================================
\newcommand{\poppill}[3]{%
  \tikz[baseline=-0.55ex]\node[draw=popink,line width=1.2pt,rounded corners=2.6pt,%
    fill=#2,inner xsep=6.5pt,inner ysep=3pt]{%
    \popxboldls\fontsize{7}{7}\selectfont\color{#3}\MakeUppercase{#1}};%
}
\newcommand{\popsquiggle}[1]{%
  \tikz[baseline=-0.4ex]{
    \draw[#1,line width=2.6pt,line cap=round]
      plot [smooth] coordinates {(0,0.09) (0.34,0.01) (0.68,0.21) (1.02,0.01) (1.36,0.10)};
  }%
}
% Mot-clé surligné (marqueur orange sous le texte)
\newcommand{\cle}[1]{{\popxbold\color{popdark}#1}}

% ============================================================
% En-tête / pied
% ============================================================
\pagestyle{fancy}
\fancyhf{}
\renewcommand{\headrulewidth}{0pt}
\renewcommand{\footrulewidth}{0pt}
\lhead{\raisebox{-0.15ex}{\poplogo\fontsize{13}{13}\selectfont\color{popink}OnBuch\color{poporange}+}}
\rhead{\poppill{\DOCMATIERE\ \textbullet\ \DOCNIVEAU\ \textbullet\ Leçon \DOCLECON}{white}{popink}}
\lfoot{\popxbold\fontsize{7.6}{7.6}\selectfont\color{popmuted}\MakeUppercase{Cours signé OnBuch+}\ \textbullet\ {\popsemi\DOCTITRE}}
\rfoot{\tikz[baseline=-0.6ex]\node[rounded corners=8.5pt,fill=popink,minimum height=17pt,minimum width=17pt,inner xsep=5pt,inner sep=0]{\popxbold\fontsize{8}{8}\selectfont\color{popcream}\thepage\,/\,\pageref*{LastPage}};}

% ============================================================
% Titres
% ============================================================
\newcommand{\chaptitre}[2]{%
  \begin{tcolorbox}[enhanced,colback=poporange,colframe=popink,boxrule=2.6pt,arc=11pt,
      drop shadow=popink,left=15pt,right=15pt,top=12pt,bottom=11pt,before skip=3pt,after skip=18pt]
    \poppill{#2}{popink}{popcream}\par\vspace{9pt}
    {\raggedright\hyphenpenalty=10000\exhyphenpenalty=10000\popdisplay\fontsize{22}{24}\selectfont\color{popcream}\MakeUppercase{#1}\par}
  \end{tcolorbox}
}
% Section numérotée : pastille orange avec le numéro + titre Archivo
\newcounter{popsec}
\newcommand{\coursec}[1]{%
  \stepcounter{popsec}\setcounter{popsub}{0}%
  \par\vspace{16pt}\noindent
  \begin{minipage}{\linewidth}
  \tikz[baseline=-0.7ex]\node[draw=popink,line width=1.6pt,fill=poporange,rounded corners=4pt,minimum size=22pt,inner sep=0]{\popdisplay\fontsize{12}{12}\selectfont\color{popcream}\thepopsec};%
  \hspace{8pt}{\popdisplay\fontsize{15}{17}\selectfont\color{popink}\MakeUppercase{#1}}\par
  \vspace{4pt}\noindent{\color{poporange}\rule{36pt}{3.6pt}}
  \end{minipage}\par\nopagebreak\vspace{8pt}
}
\newcounter{popsub}
\newcommand{\courssub}[1]{%
  \stepcounter{popsub}%
  \par\vspace{9pt}\noindent{\popxbold\fontsize{11.5}{13}\selectfont\color{popink}{\color{poporange}\thepopsec.\thepopsub}\hspace{6pt}#1}\par\nopagebreak\vspace{4pt}
}

% ============================================================
% Boîtes pédagogiques — même recette : fond blanc, bordure épaisse colorée,
% ombre dure encre, badge plein. Titre optionnel [#1].
% ============================================================
\tcbset{popbase/.style={breakable,enhanced,colback=white,boxrule=2.3pt,arc=9pt,
  shadow={3pt}{-3pt}{0pt}{popink},left=14pt,right=14pt,top=11pt,bottom=11pt,before skip=11pt,after skip=15pt,
  fontupper=\popsemi\fontsize{10.6}{14.5}\selectfont\color{popink},
  fontlower=\popsemi\fontsize{10.6}{14.5}\selectfont\color{popink}}}
% Coche dessinée (le glyphe ✓ n'existe pas dans les polices du document)
\newcommand{\popcheck}[1]{\tikz[baseline=-0.5ex]\draw[#1,line width=1.6pt,line cap=round,line join=round](-0.13,0.0)--(-0.04,-0.09)--(0.13,0.1);}
\providecommand{\coche}{\popcheck{popgreen}}
\providecommand{\resultat}[1]{\par\smallskip\noindent\fcolorbox{popgreen}{popgreenL}{\parbox{\dimexpr\linewidth-2\fboxsep-2\fboxrule\relax}{\textbf{Résultat.} #1}}\par\smallskip}
\newcommand{\popboxhead}[3]{\poppill{#1}{#2}{white}\ifx\relax#3\relax\else\hspace{6pt}{\popxbold\fontsize{10.6}{12}\selectfont\color{popink}#3}\fi\par\vspace{7pt}}

\newtcolorbox{aretenir}[1][]{popbase,colframe=popgreen,before upper={\popboxhead{À retenir}{popgreen}{#1}}}
\newtcolorbox{exemplebox}[1][]{popbase,colframe=poporange,before upper={\popboxhead{Exemple}{poporange}{#1}}}
\newtcolorbox{definition}[1][]{popbase,colframe=popblue,before upper={\popboxhead{Définition}{popblue}{#1}}}
\newtcolorbox{propriete}[1][]{popbase,colframe=poppurple,before upper={\popboxhead{Propriété}{poppurple}{#1}}}
\newtcolorbox{methode}[1][]{popbase,colframe=popgold,before upper={\popboxhead{Méthode}{popgold}{#1}}}
\newtcolorbox{attention}[1][]{popbase,colframe=poppink,before upper={\popboxhead{Attention}{poppink}{#1}}}
\newtcolorbox{experience}[1][]{popbase,colframe=popblue,colback=popblueL,before upper={\popboxhead{Expérience}{popblue}{#1}}}
% « Le savais-tu ? » : carte violette pleine (contexte, culture, Cameroun)
\newtcolorbox{savaistu}[1][]{popbase,colback=poppurple,colframe=popink,
  fontupper=\popsemi\fontsize{10.4}{14.2}\selectfont\color{white},
  before upper={\poppill{Le savais-tu\,?}{white}{poppurple}\ifx\relax#1\relax\else\hspace{6pt}{\popxbold\color{white}#1}\fi\par\vspace{7pt}}}
% Exercice résolu : énoncé en haut, \tcblower puis la solution sur fond crème
\newcounter{popexo}\newcounter{popexr}
\newtcolorbox{exoresolu}[1][]{popbase,colframe=poporange,colbacklower=popcream,
  segmentation style={popink,line width=1.2pt,dashed},
  before upper={\stepcounter{popexr}\popboxhead{Exercice résolu \thepopexr}{poporange}{#1}},
  before lower={\poppill{Solution}{popink}{popcream}\par\vspace{6pt}}}
% Exercice (non résolu en place) — la correction est dans la partie Corrigés
\newtcolorbox{exercice}[1][]{popbase,colframe=popink,
  before upper={\stepcounter{popexo}\popboxhead{Exercice \thepopexo}{popink}{#1}}}
% Corrigé
\newtcolorbox{corrige}[1][]{popbase,colframe=popgreen,colback=popgreenL,
  before upper={\popboxhead{Corrigé}{popgreen}{#1}}}
% Objectifs / prérequis / bilan
\newtcolorbox{objectifs}{popbase,colframe=popink,colback=poporangeL,
  before upper={\popboxhead{Objectifs de la leçon}{popink}{}}}
\newtcolorbox{prerequis}{popbase,colframe=popink,colback=popblueL,
  before upper={\popboxhead{Prérequis}{popblue}{}}}
\newtcolorbox{bilan}{popbase,colframe=popink,colback=white,boxrule=2.8pt,
  before upper={\popboxhead{Fiche bilan}{poporange}{L'essentiel à connaître pour l'examen}}}
% Figure encadrée avec légende
\newtcolorbox{popfigure}[1][]{enhanced,breakable=false,colback=white,colframe=popline,boxrule=1.4pt,arc=8pt,
  left=10pt,right=10pt,top=10pt,bottom=8pt,before skip=10pt,after skip=12pt,halign=center,
  lower separated=false,halign lower=center,
  fontlower=\popsemi\fontsize{9}{11.5}\selectfont\color{popmuted},
  #1}
\newcommand{\legende}[1]{\par\vspace{6pt}{\popsemi\fontsize{9}{11.5}\selectfont\color{popmuted}#1}}

% ============================================================
% Signature OnBuch+
% ============================================================
\newcommand{\poplogoinline}[1]{{\poplogo\fontsize{#1}{#1}\selectfont\color{popink}OnBuch\color{poporange}+}}
\newcommand{\signatureonbuch}{%
  \begin{tcolorbox}[enhanced,colback=popink,colframe=popink,boxrule=2.2pt,arc=10pt,
      drop shadow=poporange,left=14pt,right=14pt,top=10pt,bottom=10pt,before skip=0pt,after skip=0pt]
    \tikz[baseline=-0.7ex]\node[circle,fill=popgreen,draw=popcream,line width=1.3pt,minimum size=22pt,inner sep=0]{\popcheck{white}};%
    \hspace{9pt}\begin{minipage}[c]{0.62\linewidth}
      {\popxbold\fontsize{12}{13}\selectfont\color{popcream}Signé\ \textbullet\ {\poplogo OnBuch\color{poporange}+}}\par\vspace{2pt}
      {\raggedright\popsemi\fontsize{8}{10.5}\selectfont\color{popline}\MakeUppercase{Équipe pédagogique OnBuch+}\\\MakeUppercase{Conforme au programme officiel MINESEC}\par}
    \end{minipage}\hfill
    {\poplogo\fontsize{22}{22}\selectfont\color{popcream}OnBuch\color{poporange}+}
  \end{tcolorbox}%
}

% ============================================================
% Page de garde
% #1 titre · #2 matière · #3 niveau · #4 module · #5 n° leçon · #6 durée ·
% #7 description / objectifs (texte ou itemize)
% ============================================================
\newcommand{\pagegarde}[7]{%
  \thispagestyle{empty}
  \newgeometry{a4paper,margin=1.7cm,top=1.6cm,bottom=1.4cm}%
  \begin{tikzpicture}[remember picture,overlay]
    \fill[poporange!18] ([xshift=-3.2cm,yshift=-3.6cm]current page.north east) circle (4.6cm);
    \fill[poppurple!14] ([xshift=2.4cm,yshift=7.6cm]current page.south west) circle (3.8cm);
  \end{tikzpicture}%
  {\poplogo\fontsize{30}{30}\selectfont\color{popink}OnBuch\color{poporange}+}\hfill
  \raisebox{0.6ex}{\poppill{Cours signé \textbullet\ Premium}{popink}{popcream}}\par
  \vspace{1pt}\popsquiggle{poppurple}\par
  \vspace{8pt}
  \begin{tcolorbox}[enhanced,colback=poporange,colframe=popink,boxrule=2.8pt,arc=13pt,
      drop shadow=popink,left=18pt,right=18pt,top=12pt,bottom=13pt,after skip=10pt]
    \poppill{#2 \textbullet\ #3}{popink}{popcream}\hspace{5pt}\poppill{#4}{white}{popink}\par\vspace{8pt}
    {\popdisplay\fontsize{14}{14}\selectfont\color{popink}LEÇON #5}\par\vspace{4pt}
    {\raggedright\hyphenpenalty=10000\exhyphenpenalty=10000\popdisplay\fontsize{25}{27}\selectfont\color{popcream}\MakeUppercase{#1}\par}
    \vspace{8pt}
    {\popxbold\fontsize{9.5}{9.5}\selectfont\color{popink}\MakeUppercase{Durée conseillée : #6}}
  \end{tcolorbox}
  \begin{tcolorbox}[enhanced,colback=white,colframe=popink,boxrule=2.2pt,arc=10pt,
      drop shadow=popink,left=16pt,right=16pt,top=10pt,bottom=10pt,after skip=8pt]
    \poppill{Dans cette leçon}{poppurple}{white}\par\vspace{5pt}
    \popsemi\fontsize{9.3}{12.6}\selectfont\color{popink}#7
  \end{tcolorbox}
  \noindent\begin{minipage}[t]{0.487\linewidth}
    \begin{tcolorbox}[enhanced,colback=popgreen,colframe=popink,boxrule=2.2pt,arc=10pt,
        drop shadow=popink,left=11pt,right=11pt,top=10pt,bottom=10pt,height=3.1cm,valign=center]
      \tcbox[colback=white,colframe=popink,boxrule=1.3pt,arc=5pt,boxsep=0pt,nobeforeafter,
        left=3pt,right=3pt,top=3pt,bottom=3pt,valign=center]{\qrcode[height=1.9cm]{https://play.google.com/store/apps/details?id=com.luvvix.onbuch&hl=fr}}%
      \hspace{8pt}\begin{minipage}[c]{0.5\linewidth}
        {\popdisplay\fontsize{11}{12.5}\selectfont\color{white}\MakeUppercase{Télécharge OnBuch}}\par\vspace{4pt}
        {\raggedright\popsemi\fontsize{8.4}{11}\selectfont\color{white}Gratuit \textbullet\ Google Play\\Tuteur IA, annales, concours, résultats\par}
      \end{minipage}
    \end{tcolorbox}
  \end{minipage}\hfill
  \begin{minipage}[t]{0.487\linewidth}
    \begin{tcolorbox}[enhanced,colback=poppurple,colframe=popink,boxrule=2.2pt,arc=10pt,
        drop shadow=popink,left=11pt,right=11pt,top=10pt,bottom=10pt,height=3.1cm,valign=center]
      \tikz[baseline=-0.7ex]\node[circle,fill=white,draw=popink,line width=1.3pt,minimum size=1.6cm,inner sep=0]{\popdisplay\fontsize{24}{24}\selectfont\color{poppurple}+};%
      \hspace{8pt}\begin{minipage}[c]{0.6\linewidth}
        {\popdisplay\fontsize{11}{12.5}\selectfont\color{white}\MakeUppercase{Passe à OnBuch+}}\par\vspace{4pt}
        {\raggedright\popsemi\fontsize{8.4}{11}\selectfont\color{white}Tous les cours signés, Tuteur IA illimité, fiches PDF — onglet OnBuch+ de l'app\par}
      \end{minipage}
    \end{tcolorbox}
  \end{minipage}
  \vspace{8pt}
  \popcontenu
  \vfill
  \signatureonbuch
  \restoregeometry
  \newpage
}

% Rangée « Ce cours contient » (page de garde)
\newcommand{\poptile}[3]{% #1 couleur #2 gros texte #3 légende
  \begin{tcolorbox}[enhanced,colback=#1,colframe=popink,boxrule=2pt,arc=9pt,shadow={2.5pt}{-2.5pt}{0pt}{popink},
      left=6pt,right=6pt,top=9pt,bottom=9pt,halign=center,valign=center,height=1.75cm,width=\linewidth,nobeforeafter]
    {\popdisplay\fontsize{12.5}{13}\selectfont\color{white}\MakeUppercase{#2}}\par\vspace{3pt}
    {\centering\popsemi\fontsize{7.8}{9.5}\selectfont\color{white}#3\par}
  \end{tcolorbox}}
\newcommand{\popcontenu}{%
  \par\noindent{\popxboldls\fontsize{7.5}{7.5}\selectfont\color{popmuted}CE COURS SIGNÉ CONTIENT}\par\vspace{7pt}
  \noindent\begin{minipage}[t]{0.235\linewidth}\poptile{popblue}{Cours}{complet, illustré, pas à pas}\end{minipage}\hfill
  \begin{minipage}[t]{0.235\linewidth}\poptile{poporange}{Méthodes}{et exercices résolus}\end{minipage}\hfill
  \begin{minipage}[t]{0.235\linewidth}\poptile{poppink}{Exercices}{type examen + corrigés}\end{minipage}\hfill
  \begin{minipage}[t]{0.235\linewidth}\poptile{popgreen}{Fiche bilan}{l'essentiel à retenir}\end{minipage}\par}

% Sommaire
\newtcolorbox{sommaire}{enhanced,colback=white,colframe=popink,boxrule=2.2pt,arc=10pt,
  drop shadow=popink,left=18pt,right=18pt,top=13pt,bottom=13pt,before skip=6pt,after skip=20pt,
  before upper={\poppill{Sommaire}{poppurple}{white}\par\vspace{9pt}\popsemi\fontsize{10.6}{14.5}\selectfont\color{popink}}}

% Signature de fin de cours — poussée tout en bas de la dernière page
\newcommand{\findecours}{%
  \par\vfill\nopagebreak
  \begin{center}\popsquiggle{poporange}\end{center}\vspace{4pt}
  \signatureonbuch
}
\AtBeginDocument{\def\llbracket{\mathopen{[\mkern-3.5mu[}}\def\rrbracket{\mathclose{]\mkern-3.5mu]}}} % intervalles d'entiers (glyphe ⟦ absent de la police)
% Glyphes absents de Fira Math : redéfinis à partir de glyphes disponibles
\AtBeginDocument{%
  \def\setminus{\mathbin{\backslash}}\let\smallsetminus\setminus
  \def\vdots{\mathord{\vbox{\baselineskip3.2pt\lineskiplimit0pt\kern4pt\hbox{.}\hbox{.}\hbox{.}}}}%
  \def\ddots{\mathinner{\mkern1mu\raise6.5pt\hbox{.}\mkern2mu\raise3.6pt\hbox{.}\mkern2mu\raise0.7pt\hbox{.}\mkern1mu}}%
  \def\triangle{\mathord{\scalebox{0.9}{$\symup{\Delta}$}}}%
  \def\nabla{\mathord{\raisebox{\depth}{\scalebox{1}[-1]{$\symup{\Delta}$}}}}%
  \def\checkmark{\popcheck{popdark}}%
  \def\star{\mathbin{\ast}}\def\bigstar{\mathord{\ast}}%
  \def\diamond{\mathbin{\scalebox{0.6}{$\Diamond$}}}%
  \def\bigcup{\mathop{\mathchoice{\raisebox{-0.3em}{\scalebox{1.7}{$\cup$}}}{\scalebox{1.25}{$\cup$}}{\cup}{\cup}}\displaylimits}%
  \def\bigcap{\mathop{\mathchoice{\raisebox{-0.3em}{\scalebox{1.7}{$\cap$}}}{\scalebox{1.25}{$\cap$}}{\cap}{\cap}}\displaylimits}%
  \def\lesseqqgtr{\mathrel{\lessgtr}}\def\bigoplus{\mathop{\oplus}\displaylimits}%
}
\providecommand{\ssi}{si et seulement si} % abréviation fréquente
\AtBeginDocument{\providecommand{\wideparen}{\overparen}} % arc de cercle

%% FIGFIX-BEGIN v4 — correctifs globaux des figures (docs/onbuch-generation/tools/figfix)
\usepackage{adjustbox}\usepackage{etoolbox}
% toute figure TikZ/pgfplots plus large que la ligne est réduite à la largeur disponible
\BeforeBeginEnvironment{tikzpicture}{\begin{adjustbox}{max width=\linewidth}}
\AfterEndEnvironment{tikzpicture}{\end{adjustbox}}
% légendes pgfplots lisibles par-dessus les courbes
\pgfplotsset{every axis/.append style={legend style={fill=white,fill opacity=0.92,text opacity=1,draw=black!15,inner sep=3pt}}}
% étiquettes d'arêtes (syntaxe edge["..."]) avec fond blanc
\tikzset{every edge quotes/.append style={fill=white,inner sep=1.2pt}}
%% FIGFIX-END
\begin{document}

\pagegarde{Amélioration de la santé du système immunitaire}{SVT}{1ère A}{Module 2 : L'Éducation à la Santé}{NA04}{3 h}{Cette leçon présente l'organisation du système immunitaire (organes et cellules), les notions d'antigène, d'anticorps et de réaction immunitaire, puis le rôle de la vaccination dans le renforcement de l'immunité. Elle débouche sur l'identification des facteurs qui affaiblissent les défenses et sur la proposition de comportements favorables à une bonne santé immunitaire.
\vspace{4pt}
\begin{multicols}{2}\small
\begin{itemize}
\item À la fin de cette leçon, je sais citer les organes et les cellules du système immunitaire
\item À la fin de cette leçon, je sais définir les notions d'antigène, d'anticorps et de réaction immunitaire
\item À la fin de cette leçon, je sais expliquer le principe général d'une réaction immunitaire (reconnaissance, défense, mémoire)
\item À la fin de cette leçon, je sais distinguer immunité naturelle et immunité acquise
\item À la fin de cette leçon, je sais justifier l'intérêt de la vaccination à partir du principe de la mémoire immunitaire
\item À la fin de cette leçon, je sais citer les facteurs qui affaiblissent le système immunitaire (malnutrition, stress, maladies)
\end{itemize}\end{multicols}}

\begin{sommaire}
\begin{multicols}{2}
\begin{enumerate}
\item Le système immunitaire : organes et cellules
\item Antigènes, anticorps et réaction immunitaire
\item Les types d'immunité
\item La vaccination, moyen de renforcement de l'immunité
\item Facteurs affaiblissant le système immunitaire
\item Comportements favorables au renforcement de l'immunité
\item Activité d'intégration
\item Méthodes et exercices résolus
\item Exercices
\item Corrigés des exercices
\item Fiche bilan
\end{enumerate}
\end{multicols}
\end{sommaire}

\begin{objectifs}
\begin{itemize}
\item À la fin de cette leçon, je sais citer les organes et les cellules du système immunitaire
\item À la fin de cette leçon, je sais définir les notions d'antigène, d'anticorps et de réaction immunitaire
\item À la fin de cette leçon, je sais expliquer le principe général d'une réaction immunitaire (reconnaissance, défense, mémoire)
\item À la fin de cette leçon, je sais distinguer immunité naturelle et immunité acquise
\item À la fin de cette leçon, je sais justifier l'intérêt de la vaccination à partir du principe de la mémoire immunitaire
\item À la fin de cette leçon, je sais citer les facteurs qui affaiblissent le système immunitaire (malnutrition, stress, maladies)
\item À la fin de cette leçon, je sais proposer des comportements favorables au renforcement de l'immunité
\item À la fin de cette leçon, je sais interpréter un schéma ou un document sur la réaction immunitaire
\end{itemize}
\end{objectifs}

\begin{prerequis}
\begin{itemize}
\item Notion de microbe et d'agent pathogène (virus, bactéries) vue en classe de Seconde
\item Constitution du sang : globules rouges, globules blancs, plaquettes (Seconde)
\item Notion de cellule et d'organes du corps humain
\item Notions d'hygiène et de maladies transmissibles (début du Module 2)
\item Lecture et interprétation de schémas et de tableaux de données
\end{itemize}
\end{prerequis}

\newpage

\coursec{Le système immunitaire : organes et cellules}

\textbf{Situation d'accroche.} Pendant la saison des pluies à Yaoundé, le paludisme frappe plusieurs élèves du lycée : certains tombent malades plusieurs fois, d'autres jamais, alors qu'ils vivent dans le même quartier, fréquentent la même cour de récréation et sont piqués par les mêmes moustiques. Au centre de santé du quartier, l'infirmier fait remarquer deux autres faits troublants : les enfants vaccinés contre la rougeole la contractent rarement, et les élèves mal nourris ou stressés en période d'examens attrapent plus facilement les infections (rhumes, angines, gastro-entérites). Trois questions se posent alors à nous : pourquoi notre organisme résiste-t-il différemment aux microbes ? Comment la vaccination protège-t-elle durablement ? Quels comportements renforcent nos défenses naturelles ? Pour répondre à ces questions, il faut d'abord découvrir l'« armée intérieure » de notre corps : le système immunitaire, ses organes et ses cellules. C'est l'objet de cette première section.

\courssub{Qu'est-ce que le système immunitaire ?}

\textbf{Observation et intuition.} Chaque jour, nous respirons, buvons, mangeons et touchons des millions de microbes : bactéries, virus, champignons, parasites. Pourtant, la plupart du temps, nous ne tombons pas malades. Mieux : quand tu te coupes légèrement au doigt en épluchant un manioc, la plaie rougit parfois un peu, puis guérit seule en quelques jours. Quelque chose, à l'intérieur du corps, a donc repéré les microbes entrés par la plaie, les a combattus et éliminés. Ce « quelque chose » n'est pas un organe unique comme le cœur ou le foie : c'est un ensemble coordonné d'organes, de cellules et de molécules répartis dans tout le corps.

\begin{definition}[Le système immunitaire]
Le \cle{système immunitaire} est l'ensemble des \textbf{organes} (organes lymphoïdes), des \textbf{cellules} (leucocytes ou globules blancs) et des \textbf{molécules} (anticorps, substances de défense) qui assurent la \cle{défense de l'organisme} contre les \cle{agents pathogènes} (microbes capables de provoquer une maladie : virus, bactéries, champignons, parasites) et contre les cellules anormales de l'organisme lui-même (cellules cancéreuses).
\end{definition}

\begin{aretenir}[Les trois missions du système immunitaire]
\begin{itemize}
\item \textbf{Reconnaître} : distinguer ce qui appartient au corps (le « soi ») de ce qui est étranger (le « non-soi » : microbes, mais aussi cellules d'un greffon).
\item \textbf{Défendre} : détruire ou neutraliser les agents étrangers pénétrant dans l'organisme.
\item \textbf{Mémoriser} : garder le souvenir d'un microbe déjà rencontré, pour réagir plus vite et plus fort lors d'un contact ultérieur. C'est cette mémoire qui explique l'efficacité durable de la vaccination (section 4).
\end{itemize}
\end{aretenir}

\courssub{Les organes du système immunitaire : les organes lymphoïdes}

\textbf{Observation.} Lors d'une angine, le médecin palpe le cou du malade et regarde sa gorge. Il cherche des ganglions gonflés et des amygdales rouges. Pourquoi ces organes réagissent-ils à une infection ? Parce qu'ils font partie des \cle{organes lymphoïdes}, les « casernes » et les « postes de contrôle » de notre armée intérieure. On distingue deux catégories d'organes lymphoïdes.

\textbf{Les organes lymphoïdes primaires (ou centraux)}

Ce sont les organes où les cellules immunitaires sont \textbf{produites} et où elles \textbf{achèvent leur maturation} (leur « formation militaire »).

\begin{itemize}
\item \textbf{La moelle osseuse} : tissu rouge situé à l'intérieur des os (os longs, côtes, sternum, bassin). C'est l'\textbf{usine de fabrication} : elle produit en permanence toutes les cellules sanguines, dont tous les leucocytes, à partir de cellules souches hématopoïétiques.
\item \textbf{Le thymus} : organe situé dans le thorax, derrière le sternum, très actif chez l'enfant et l'adolescent, qui s'involutionne (diminue de taille) à l'âge adulte. C'est l'\textbf{école de maturation} d'une catégorie de lymphocytes : les lymphocytes T (la lettre T vient d'ailleurs de \emph{thymus}). Les lymphocytes T y apprennent à reconnaître le « non-soi » sans attaquer le « soi » (tolérance centrale).
\end{itemize}

\textbf{Les organes lymphoïdes secondaires (ou périphériques)}

Ce sont les organes où les cellules immunitaires, devenues matures, se \textbf{postent en embuscade} pour rencontrer les microbes et déclencher les réactions de défense.

\begin{itemize}
\item \textbf{La rate} : organe situé dans l'hypocondre gauche, sous les côtes. Elle filtre le \textbf{sang}, détruit les globules rouges usés et capture les microbes circulant dans le sang.
\item \textbf{Les ganglions lymphatiques} : petits organes en forme de haricot (quelques millimètres à un centimètre), répartis le long des vaisseaux lymphatiques, notamment dans le cou, les aisselles et l'aine. Ils filtrent la \textbf{lymphe} (liquide clair qui baigne les tissus) et retiennent les microbes venus des tissus.
\item \textbf{Les amygdales} : situées au fond de la gorge (amygdales palatines, tubaires, linguale et pharyngée), elles montent la garde à l'entrée des voies digestives et respiratoires, première porte d'entrée des microbes (air, aliments).
\item \textbf{L'appendice} et les plaques de Peyer de l'intestin : ils surveillent les microbes du tube digestif (immunité muqueuse).
\end{itemize}

\begin{popfigure}
\begin{center}
\begin{tikzpicture}[scale=1.0, every node/.style={font=\small}]
% Corps humain simplifié
\draw[thick, popdark, fill=popcream] (0,4.6) circle (0.55); % tête
\draw[thick, popdark, fill=popcream] (-0.9,3.9) .. controls (-1.1,3.2) and (-1.1,1.8) .. (-0.8,1.2) -- (-0.55,0.0) -- (-0.15,0.0) -- (0,1.0) -- (0.15,0.0) -- (0.55,0.0) -- (0.8,1.2) .. controls (1.1,1.8) and (1.1,3.2) .. (0.9,3.9) -- cycle; % tronc+jambes
\draw[thick, popdark, fill=popcream] (-0.95,3.7) -- (-1.7,2.4) -- (-1.45,2.2) -- (-0.8,3.2); % bras gauche
\draw[thick, popdark, fill=popcream] (0.95,3.7) -- (1.7,2.4) -- (1.45,2.2) -- (0.8,3.2); % bras droit
% Thymus
\draw[fill=poporange, draw=popdark] (0,3.35) ellipse (0.28 and 0.18);
% Amygdales
\draw[fill=poppink, draw=popdark] (-0.18,4.05) circle (0.09);
\draw[fill=poppink, draw=popdark] (0.18,4.05) circle (0.09);
% Rate
\draw[fill=poppurple, draw=popdark] (-0.55,2.55) ellipse (0.22 and 0.32);
% Appendice
\draw[fill=popgreen, draw=popdark] (0.45,1.55) circle (0.10);
% Ganglions (cou, aisselles, aine)
\draw[fill=popblue, draw=popdark] (-0.30,3.85) circle (0.08);
\draw[fill=popblue, draw=popdark] (0.30,3.85) circle (0.08);
\draw[fill=popblue, draw=popdark] (-0.85,3.35) circle (0.08);
\draw[fill=popblue, draw=popdark] (0.85,3.35) circle (0.08);
\draw[fill=popblue, draw=popdark] (-0.35,1.35) circle (0.08);
\draw[fill=popblue, draw=popdark] (0.35,1.35) circle (0.08);
% Moelle osseuse (os du bras et de la jambe)
\draw[fill=popgold, draw=popdark] (-1.55,2.95) ellipse (0.09 and 0.30);
\draw[fill=popgold, draw=popdark] (-0.35,0.55) ellipse (0.09 and 0.30);
% Légendes pointées
\node[anchor=west] at (2.1,4.05) {\textbf{Amygdales}};
\draw[thin, popmuted] (0.28,4.05) -- (2.05,4.05);
\node[anchor=west] at (2.1,3.35) {\textbf{Thymus}};
\draw[thin, popmuted] (0.30,3.35) -- (2.05,3.35);
\node[anchor=east] at (-2.1,3.85) {\textbf{Ganglions lymphatiques}};
\draw[thin, popmuted] (-0.40,3.85) -- (-2.05,3.85);
\node[anchor=east] at (-2.1,2.55) {\textbf{Rate}};
\draw[thin, popmuted] (-0.80,2.55) -- (-2.05,2.55);
\node[anchor=west] at (2.1,1.55) {\textbf{Appendice}};
\draw[thin, popmuted] (0.57,1.55) -- (2.05,1.55);
\node[anchor=east] at (-2.1,0.55) {\textbf{Moelle osseuse}};
\draw[thin, popmuted] (-0.46,0.55) -- (-2.05,0.55);
\end{tikzpicture}
\end{center}
\legende{Localisation simplifiée des organes lymphoïdes dans le corps humain. Organes lymphoïdes \textbf{primaires} : moelle osseuse (production de tous les leucocytes) et thymus (maturation des lymphocytes T). Organes lymphoïdes \textbf{secondaires} : rate, ganglions lymphatiques, amygdales, appendice (lieux de rencontre entre cellules immunitaires et antigènes).}
\end{popfigure}

\begin{savaistu}[Une usine qui ne s'arrête jamais]
La moelle osseuse d'un adulte produit chaque jour des \textbf{milliards de globules blancs} (environ \num{1e11} neutrophiles par jour), ainsi que des centaines de milliards de globules rouges et de plaquettes. Chez l'enfant, presque tous les os contiennent de la moelle rouge active ; chez l'adulte, elle se concentre surtout dans le bassin, le sternum, les côtes et le crâne. C'est pourquoi, lors de certaines maladies du sang (leucémies, aplasies), les médecins prélèvent un peu de moelle au niveau du bassin (crête iliaque) ou du sternum pour l'analyser : c'est le \emph{myélogramme}, pratiqué dans les hôpitaux de référence du Cameroun (CHU de Yaoundé, Douala, etc.).
\end{savaistu}

\courssub{Les cellules du système immunitaire : les leucocytes}

\textbf{Observation.} Au laboratoire d'analyses médicales, on réalise parfois une « NFS » (numération formule sanguine) : on compte les cellules du sang. Le résultat distingue les globules rouges (environ \num{5e6} par mm$^3$ chez l'homme) et les globules blancs (environ \num{4000} à \num{10000} par mm$^3$). Lors d'une infection bactérienne, le nombre de globules blancs augmente nettement (leucocytose) : c'est le signe que l'armée de défense se mobilise.

\begin{definition}[Leucocyte]
Un \cle{leucocyte} (ou \cle{globule blanc}) est une cellule sanguine qui participe à la \textbf{défense de l'organisme} contre les agents pathogènes. Contrairement aux globules rouges, les leucocytes possèdent un noyau et peuvent sortir des vaisseaux sanguins pour gagner les tissus infectés (diapédèse).
\end{definition}

On distingue deux grandes familles de leucocytes, selon leur mode d'action.

\textbf{1. Les phagocytes : les « mangeurs » de microbes (immunité innée).} Ce sont des cellules capables de réaliser la \cle{phagocytose} : elles englobent le microbe dans leur cytoplasme (formation d'un phagosome) puis le digèrent grâce aux enzymes lysosomales (formation d'un phagolysosome). On y trouve :
\begin{itemize}
\item les \textbf{polynucléaires neutrophiles} (ou neutrophiles), cellules à noyau découpé en plusieurs lobes, très abondantes dans le sang (50 à 70 \% des leucocytes) ; ce sont les premières à arriver massivement sur le lieu d'une infection (chimiotaxie) ;
\item les \textbf{monocytes} du sang (2 à 8 \% des leucocytes), qui deviennent des \textbf{macrophages} lorsqu'ils pénètrent dans les tissus : de grosses cellules au cytoplasme abondant, capables de phagocyter de nombreux microbes, de nettoyer les débris cellulaires et de présenter des antigènes aux lymphocytes.
\end{itemize}

\textbf{2. Les lymphocytes : les « spécialistes » de la reconnaissance spécifique (immunité adaptative).} Ce sont de petites cellules à gros noyau arrondi, cytoplasme peu abondant, capables de reconnaître \emph{un} antigène précis grâce à leurs récepteurs de membrane. On distingue :
\begin{itemize}
\item les \textbf{lymphocytes B}, qui achèvent leur maturation dans la moelle osseuse (chez l'homme) et qui, une fois activés, se différencient en \textbf{plasmocytes} producteurs d'\textbf{anticorps} (étudiés en section 2) ;
\item les \textbf{lymphocytes T}, qui migrent de la moelle osseuse vers le \textbf{thymus} pour y mûrir (sélection positive et négative), et qui se différencient en lymphocytes T cytotoxiques (destruction directe des cellules infectées) ou T auxiliaires (coordination de la défense par cytokines).
\end{itemize}

\begin{center}
\begin{tabularx}{\linewidth}{|l|X|X|}
\hline
\rowcolor{popblueL} \textbf{Cellule} & \textbf{Lieu de production / maturation} & \textbf{Rôle principal} \\
\hline
Polynucléaires neutrophiles & Moelle osseuse & Phagocytose rapide des microbes ; premiers intervenants (innée) \\
\hline
Monocytes $\rightarrow$ Macrophages & Moelle osseuse (monocytes) ; Tissus (macrophages) & Phagocytose, nettoyage, présentation d'antigènes (innée / lien adaptative) \\
\hline
Lymphocytes B & Moelle osseuse (production et maturation) & Différenciation en plasmocytes sécrétant des anticorps spécifiques (adaptative humorale) \\
\hline
Lymphocytes T & Moelle osseuse (production), Thymus (maturation) & Cytotoxicité (T CD8+) ou aide/coordination (T CD4+) (adaptative cellulaire) \\
\hline
\end{tabularx}
\end{center}

\begin{attention}[Ne pas confondre globules blancs et globules rouges]
C'est l'erreur la plus fréquente au Probatoire ! Les \textbf{globules rouges} (hématies ou érythrocytes) transportent le \textbf{dioxygène} des poumons vers les organes grâce à l'hémoglobine ; ils n'ont pas de noyau (anucleés chez le mammifère adulte) et ne participent \emph{pas} à la défense. Les \textbf{globules blancs} (leucocytes) assurent la \textbf{défense} de l'organisme ; ils ont un noyau et sont beaucoup moins nombreux (ratio ~ 1/1000). Retiens : \emph{rouge = respiration (transport de \ce{O2}) ; blanc = bouclier (défense)}. Une anémie (manque de globules rouges ou d'hémoglobine) n'est donc pas une baisse des défenses immunitaires, bien qu'elle affaiblisse l'organisme globalement.
\end{attention}

\courssub{Où circulent les cellules immunitaires ?}

\textbf{Observation.} Un microbe peut pénétrer n'importe où : par une coupure au pied, par les voies respiratoires, par un aliment contaminé. Comment les cellules de défense, fabriquées dans les os, peuvent-elles intervenir partout dans le corps ? Réponse : elles circulent en permanence, comme des patrouilles, dans deux réseaux complémentaires.

\begin{itemize}
\item \textbf{Le sang} : les leucocytes y circulent (marge vasculaire) et peuvent traverser la paroi des capillaires sanguins (vénules post-capillaires) pour gagner les tissus infectés : c'est la \cle{diapédèse} (ou extravasation), guidée par des molécules d'adhésion et des chémokines. Le sang irrigue tous les organes, donc la patrouille est permanente.
\item \textbf{La lymphe} : liquide clair (ultrafiltrat du plasma) qui s'écoule des tissus dans des vaisseaux lymphatiques. Ces vaisseaux traversent les \textbf{ganglions lymphatiques}, qui filtrent la lymphe avant qu'elle ne rejoigne le sang veineux (angle veineux jugulo-sous-clavier). Les lymphocytes circulent sans cesse entre le sang, la lymphe et les organes lymphoïdes secondaires (recirculation lymphocytaire), ce qui maximise leurs chances de rencontrer leur antigène cogné.
\end{itemize}

\begin{exemplebox}[Pourquoi les ganglions gonflent-ils pendant une angine ?]
Lors d'une infection de la gorge (angine bactérienne à \emph{Streptococcus} ou virale), les microbes pénètrent dans les tissus du pharynx. La lymphe qui draine cette zone transporte les microbes (ou leurs antigènes) et les cellules présentatrices d'antigènes jusqu'aux \textbf{ganglions cervicaux} (du cou). Ces ganglions jouent leur rôle de \textbf{sentinelles} : ils filtrent la lymphe, retiennent les antigènes, et les lymphocytes spécifiques qu'ils contiennent s'activent, prolifèrent (clonal expansion) et se différencient. Cette affluence massive de cellules fait \textbf{gonfler} le ganglion (adénopathie), qui devient palpable et parfois douloureux sous la mâchoire. Un ganglion gonflé est donc le signe visible d'une défense immunitaire adaptative en cours — c'est exactement ce que recherche le médecin en palpant le cou du malade.
\end{exemplebox}

\begin{exoresolu}[Identifier les organes et cellules mis en jeu]
\textbf{Énoncé.} Junior, élève de Première A, se blesse au pied avec un clou rouillé en jouant au football. Deux jours plus tard, la plaie est rouge, chaude, douloureuse et contient du pus. Le médecin du centre de santé palpe l'aine de Junior et sent un ganglion gonflé.
\begin{enumerate}
\item Quelles cellules immunitaires sont intervenues en premier sur la plaie, et quel mécanisme ont-elles réalisé ?
\item Expliquer pourquoi un ganglion de l'aine (et non du cou) a gonflé.
\item Préciser l'organe dans lequel ces cellules immunitaires ont été produites.
\end{enumerate}
\tcblower
\textbf{Solution détaillée.}
\begin{enumerate}
\item Les premières cellules intervenues sont les \textbf{phagocytes de l'immunité innée} : les polynucléaires neutrophiles (arrivés par le sang via chimiotaxie et diapédèse) et les macrophages résidents du tissu conjonctif. Ils ont réalisé la \textbf{phagocytose} : ils ont englobé puis digéré les bactéries (ex. \emph{Clostridium tetani}, staphylocoques) entrées par la plaie. Le pus est l'accumulation de microbes détruits, de phagocytes morts (cellules en décomposition), de débris tissulaires et de liquide inflammatoire.
\item La lymphe des tissus du pied et de la jambe est drainée par les vaisseaux lymphatiques vers les \textbf{ganglions inguinaux} (de l'aine), qui filtrent cette lymphe. Les antigènes bactériens transportés par la lymphe y sont retenus et présentés aux lymphocytes ; ceux-ci se multiplient alors activement (prolifération clonale), ce qui fait gonfler le ganglion. Les ganglions du cou (cervicaux) drainent la tête et la gorge : ils ne sont pas concernés par une infection du membre inférieur.
\item Toutes ces cellules (phagocytes : polynucléaires, monocytes/macrophages ; lymphocytes B et T) ont été \textbf{produites dans la moelle osseuse rouge} (myéloïse et lymphopoïèse B), l'organe lymphoïde primaire où naissent tous les leucocytes à partir de cellules souches hématopoïétiques. Les lymphocytes T ont ensuite migré vers le thymus pour leur maturation.
\end{enumerate}
\end{exoresolu}

\begin{aretenir}[L'essentiel de la section]
\begin{itemize}
\item Le système immunitaire = organes lymphoïdes + cellules (leucocytes) + molécules de défense ; ses missions : reconnaître (soi/non-soi), défendre, mémoriser.
\item Organes lymphoïdes \textbf{primaires} : moelle osseuse (production de tous les leucocytes ; maturation des lymphocytes B) et thymus (maturation des lymphocytes T, tolérance centrale).
\item Organes lymphoïdes \textbf{secondaires} : rate (filtre sanguin), ganglions lymphatiques (filtre lymphatique), amygdales, appendice/plaques de Peyer (immunité muqueuse) — lieux de rencontre entre lymphocytes naïfs et antigènes.
\item Deux grandes familles de leucocytes : les \textbf{phagocytes} (polynucléaires neutrophiles, macrophages : immunité innée, non spécifique, rapide) et les \textbf{lymphocytes} (B : immunité humorale/anticorps ; T : immunité cellulaire/cytotoxicité et aide : immunité adaptative, spécifique, mémoire).
\item Les leucocytes patrouillent dans le \textbf{sang} (diapédèse vers tissus) et la \textbf{lymphe} (recirculation via ganglions) ; les ganglions sont des sentinelles qui gonflent (adénopathie) lors de la réaction adaptative locale.
\end{itemize}
\end{aretenir}

Maintenant que nous connaissons les soldats et leurs casernes, il reste à comprendre \emph{comment} ils reconnaissent l'ennemi avec précision : c'est l'objet de la section suivante, consacrée aux antigènes, aux anticorps et à la réaction immunitaire.

\coursec{Antigènes, anticorps et réaction immunitaire}

Dans la section précédente, nous avons découvert les \textbf{organes} (moelle osseuse, thymus, ganglions, rate, amygdales) et les \textbf{cellules} (lymphocytes B et T, phagocytes) qui constituent l'armée de notre défense. Mais une armée ne sert à rien si elle ne sait pas \emph{qui} combattre ni \emph{comment} le faire. Comment l'organisme distingue-t-il le « soi » du « non-soi » ? Quelle est la molécule qui sert de passeport à l'ennemi ? Et comment nos lymphocytes fabriquent-ils des armes sur mesure ? C'est ce que nous allons élucider maintenant.

\courssub{Les agents pathogènes : des envahisseurs variés}

Avant de comprendre la riposte, identifions l'ennemi. Un \cle{agent pathogène} est tout micro-organisme capable de provoquer une maladie. Au Cameroun, comme partout, nous côtoyons quotidiennement quatre grands types :
\begin{itemize}
    \item Les \textbf{virus} (ex. : VIH, virus de la grippe, de la rougeole, de l'hépatite B) : acellulaires, parasites intracellulaires stricts.
    \item Les \textbf{bactéries} (ex. : \emph{Salmonella} -- fièvre typhoïde, \emph{Vibrio cholerae} -- choléra, \emph{Mycobacterium tuberculosis} -- tuberculose) : unicellulaires procaryotes, souvent extracellulaires.
    \item Les \textbf{champignons microscopiques} (ex. : \emph{Candida albicans} -- candidose, dermatophytes -- teignes) : eucaryotes, affectionnant les milieux humides.
    \item Les \textbf{parasites} (ex. : \emph{Plasmodium} -- paludisme, vers intestinaux comme \emph{Ascaris}) : eucaryotes pluricellulaires ou unicellulaires, cycles complexes.
\end{itemize}
Tous ces agents portent à leur surface (ou libèrent) des molécules qui les trahissent : les \textbf{antigènes}.

\begin{popfigure}
\begin{tikzpicture}[scale=0.9, every node/.style={font=\small}]
    % Virus
    \node[draw, thick, rounded corners, fill=popblueL, minimum width=3.5cm, minimum height=2.5cm] (vbox) at (0,0) {};
    \node[draw, circle, fill=popblue, minimum size=1.2cm] (virus) at (0,0.3) {};
    \foreach \ang/\r in {0/1.5, 45/1.2, 90/1.5, 135/1.2, 180/1.5, 225/1.2, 270/1.5, 315/1.2} {
        \draw[popblue, thick] (virus) -- ++(\ang:\r cm) node[circle, fill=poporange, minimum size=3pt] {};
    }
    \node[align=center] at (0,-1.1) {\textbf{Virus} \\ Capside + glycoprotéines (Ag)};
    
    % Bactérie
    \node[draw, thick, rounded corners, fill=popgreenL, minimum width=3.5cm, minimum height=2.5cm] (bbox) at (5,0) {};
    \node[draw, ellipse, fill=popgreen!50, minimum width=2cm, minimum height=1cm] (bact) at (5,0.3) {};
    \foreach \ang/\r in {10/1.3, 50/1.1, 90/1.3, 130/1.1, 170/1.3, 210/1.1, 250/1.3, 290/1.1, 330/1.1} {
        \draw[popgreen, thick] (bact) -- ++(\ang:\r cm) node[circle, fill=poppink, minimum size=3pt] {};
    }
    \node[align=center] at (5,-1.1) {\textbf{Bactérie} \\ Paroi + flagelles/pilis (Ag)};
    
    % Champignon
    \node[draw, thick, rounded corners, fill=poppurpleL, minimum width=3.5cm, minimum height=2.5cm] (cbox) at (10,0) {};
    \node[draw, ellipse, fill=poppurple!50, minimum width=1.5cm, minimum height=0.8cm] (champ) at (10,0.3) {};
    \draw[poppurple, thick] (champ) -- ++(90:1.2) node[above] {Spores (Ag)};
    \draw[poppurple, thick] (champ) -- ++(-90:1.2) node[below] {Mycélium};
    \node[align=center] at (10,-1.1) {\textbf{Champignon} \\ Paroi chitineuse (Ag)};
    
    % Parasite (Plasmodium)
    \node[draw, thick, rounded corners, fill=poporangeL, minimum width=3.5cm, minimum height=2.5cm] (pbox) at (15,0) {};
    \node[draw, circle, fill=poporange!50, minimum size=1.2cm] (para) at (15,0.3) {};
    \draw[poporange, thick] (para) -- ++(120:1.3) node[circle, fill=popblue, minimum size=3pt] {};
    \draw[poporange, thick] (para) -- ++(60:1.3) node[circle, fill=popblue, minimum size=3pt] {};
    \node[align=center] at (15,-1.1) {\textbf{Parasite} \\ Protéines de surface (Ag)};
    
    % Légende antigènes
    \node[draw, dashed, thick, fill=popcream, rounded corners, align=center, font=\footnotesize] at (7.5,-2.5) {
        \textbf{Point commun :} À leur surface, tous portent des \textbf{antigènes (Ag)} \\
        (glycoprotéines, protéines, polysaccharides) reconnus comme « non-soi ».
    };
\end{tikzpicture}
\legende{Principaux types d'agents pathogènes et localisation de leurs antigènes (représentés par des points colorés).}
\end{popfigure}

\courssub{La notion d'antigène : le signal d'alarme}

\begin{definition}[Antigène]
Un \textbf{antigène (Ag)} est toute substance \textbf{étrangère à l'organisme} (généralement une protéine, une glycoprotéine ou un polysaccharide de haut poids moléculaire) capable de \textbf{déclencher une réaction immunitaire spécifique} et de se fixer sur les produits de cette réaction (anticorps ou récepteurs lymphocytaires).
\end{definition}

\begin{attention}[Piège : Antigène $\neq$ Agent pathogène entier]
L'antigène n'est \textbf{pas} le microbe tout entier, mais une \textbf{molécule} (souvent un fragment de protéine appelé \emph{épitope} ou \emph{déterminant antigénique}) portée par le microbe. Un même virus porte des centaines d'antigènes différents. De plus, une toxine bactérienne, un pollen, ou un organe greffé sont aussi des antigènes sans être des agents pathogènes vivants.
\end{attention}

\courssub{La notion d'anticorps : les missiles guidés}

Si l'antigène est la cible, l'anticorps est le projectile. 

\begin{definition}[Anticorps / Immunoglobuline]
Un \textbf{anticorps (Ac)} (ou \textbf{immunoglobuline}) est une \textbf{protéine} de la famille des globulines, fabriquée et sécrétée par les \textbf{lymphocytes B différenciés en plasmocytes}. Il possède la propriété de se \textbf{fixer spécifiquement} sur l'antigène qui a provoqué sa synthèse, pour le neutraliser ou le marquer en vue de sa destruction.
\end{definition}

\begin{popfigure}
\begin{tikzpicture}[scale=1.1, every node/.style={font=\small}]
    % Structure Y
    \draw[very thick, popdark] (0,0) -- (0,2.5); % Tronc
    \draw[very thick, popdark] (0,2.5) -- (-1,4); % Bras G
    \draw[very thick, popdark] (0,2.5) -- (1,4);  % Bras D
    
    % Domaines constants (Fc)
    \draw[fill=popblue!30, draw=popblue, thick] (-0.4,0) rectangle (0.4,1.5) node[midway, align=center] {Domaine \textbf{constant (Fc)}\\\textit{Fixation phagocytes / complément}};
    
    % Domaines variables (Fab)
    \draw[fill=poporange!30, draw=poporange, thick] (-1.4,3) rectangle (-0.2,4.2) node[midway, align=center] {Domaine \textbf{variable (Fab)}\\\textit{Site de fixation Ag}};
    \draw[fill=poporange!30, draw=poporange, thick] (0.2,3) rectangle (1.4,4.2) node[midway, align=center] {Domaine \textbf{variable (Fab)}\\\textit{Site de fixation Ag}};
    
    % Chaînes lourdes et légères
    \draw[<->, popmuted] (-1.8,0.5) -- node[left] {Chaîne lourde} (-1.8,2.5);
    \draw[<->, popmuted] (-1.8,3) -- node[left] {Chaîne légère} (-1.8,4);
    
    % Ponts disulfure
    \foreach \y in {0.5, 1.2, 2.0, 3.5} {
        \draw[popmuted, dashed] (-0.2,\y) -- (0.2,\y);
    }
    \node[popmuted, font=\tiny] at (0.7,1.8) {Ponts disulfure};
    
    % Antigène fixable
    \node[draw, circle, fill=popgreen!50, minimum size=0.8cm] (ag) at (-1.8, 3.6) {Ag};
    \node[draw, circle, fill=popgreen!50, minimum size=0.8cm] (ag2) at (1.8, 3.6) {Ag};
    \draw[popgreen, thick, ->] (ag) -- (-1.1, 3.6);
    \draw[popgreen, thick, ->] (ag2) -- (1.1, 3.6);
    
    \node[align=center, font=\footnotesize] at (0,-0.7) {\textbf{Structure en Y d'un anticorps (IgG)} \\ 2 chaînes lourdes + 2 chaînes légères};
\end{tikzpicture}
\legende{Structure schématique d'une immunoglobuline G (IgG). Les deux bras (Fab) assurent la reconnaissance spécifique ; le tronc (Fc) recrute les effecteurs (phagocytes, complément).}
\end{popfigure}

\courssub{La spécificité antigène-anticorps : le modèle clé-serrure}

C'est le cœur de l'immunité adaptative : \textbf{un anticorps ne reconnaît qu'un seul épitope}. Cette spécificité repose sur la complémentarité de forme et de charges électriques entre le site de liaison de l'anticorps (paratope) et l'éitope de l'antigène.

\begin{popfigure}
\begin{tikzpicture}[scale=1.2, every node/.style={font=\small}]
    % Antigène (Serrure)
    \draw[fill=popblue!20, draw=popblue, thick, rounded corners=5pt] (0,0) rectangle (3,2);
    % Épitope : encoche
    \draw[fill=white, draw=popblue, thick] (1.2,2) -- (1.5,2.8) -- (1.8,2) -- cycle;
    \node[popblue, font=\bfseries] at (1.5,1) {Antigène};
    \node[popblue, font=\tiny] at (1.5,2.4) {Épitope};
    
    % Flèche
    \draw[thick, popdark, ->] (3.5,1) -- (5,1) node[midway, above] {Reconnaissance};
    
    % Anticorps (Clé)
    \draw[fill=poporange!20, draw=poporange, thick, rounded corners=5pt] (6,0) rectangle (9,2);
    % Paratope : pointe complémentaire
    \draw[fill=poporange!50, draw=poporange, thick] (7.2,2) -- (7.5,1.2) -- (7.8,2) -- cycle;
    \node[poporange, font=\bfseries] at (7.5,1) {Anticorps (Fab)};
    \node[poporange, font=\tiny] at (7.5,1.6) {Paratope};
    
    % Complexe
    \draw[thick, popdark, ->] (10,1) -- (11.5,1) node[midway, above] {Fixation};
    
    \draw[fill=popgreen!20, draw=popgreen, thick, rounded corners=5pt] (12.5,0) rectangle (15.5,2);
    \draw[fill=white, draw=popblue, thick] (13.7,2) -- (14,2.8) -- (14.3,2) -- cycle;
    \draw[fill=poporange!50, draw=poporange, thick] (13.7,2) -- (14,1.2) -- (14.3,2) -- cycle; % Overlap visual
    \node[popgreen, font=\bfseries] at (14,1) {Complexe Ag-Ac};
    \node[font=\tiny, align=center] at (14,2.5) {Emboîtement parfait};
    
    % Légende bas
    \node[draw, fill=popcream, rounded corners, align=center, font=\footnotesize] at (7.5,-1.2) {
        \textbf{Modèle clé-serrure :} La fixation est \textbf{spécifique}, \textbf{réversible} et \textbf{non covalente} \\
        (liaisons hydrogène, forces de Van der Waals, interactions ioniques, hydrophobes).
    };
\end{tikzpicture}
\legende{Illustration du modèle clé-serrure : complémentarité structurale entre le paratope de l'anticorps et l'éitope de l'antigène.}
\end{popfigure}

\courssub{Les étapes de la réaction immunitaire : une défense coordonnée}

Comment passe-t-on de la pénétration d'un germe à la production massive d'anticorps ? Suivons une démarche d'investigation inspirée des travaux classiques (Metchnikoff, Ehrlich, Burnet).

\begin{methode}[Démarche d'investigation : Découvrir les étapes de la réponse immunitaire]
\begin{enumerate}
    \item \textbf{Observation :} Après une primo-infection (ex. rougeole), le sujet guérit et devient résistant à une réinfection. Le sérum du convalescent peut protéger un autre sujet (transfert passif).
    \item \textbf{Hypothèses :} 
    \begin{itemize}
        \item H1 : L'antigène modifie directement une protéine existante pour en faire un anticorps (théorie instructionnelle -- fausse).
        \item H2 : L'antigène sélectionne un lymphocyte préexistant porteur du bon récepteur, qui se multiplie et sécrète l'anticorps (théorie clonale -- Burnet, 1957).
    \end{itemize}
    \item \textbf{Expérience / Données clés :} 
    \begin{itemize}
        \item Marquage radioactif : les anticorps ne sont pas des protéines modifiées, mais \emph{neo-synthétisées}.
        \item Limitation par dilution : un seul lymphocyte peut reconstituer la réponse (preuve du clonage).
        \item Séquençage : le gène de l'anticorps résulte d'une réarrangement somatique aléatoire (VDJ) \emph{avant} la rencontre avec l'Ag.
    \end{itemize}
    \item \textbf{Interprétation :} L'antigène ne \emph{crée} pas la spécificité, il \emph{sélectionne} le clone lymphocytaire préexistant porteur du récepteur complémentaire. Ce clone prolifère (expansion clonale) et se différencie.
    \item \textbf{Conclusion :} La réaction immunitaire suit 4 étapes majeures : \textbf{Reconnaissance $\rightarrow$ Activation/Prolifération $\rightarrow$ Différenciation/Production $\rightarrow$ Élimination}.
\end{enumerate}
\end{methode}

\begin{popfigure}
\begin{tikzpicture}[scale=0.85, every node/.style={font=\small}]
    % Styles
    \tikzset{
        stepbox/.style={draw, thick, rounded corners, fill=popcream, minimum width=3.8cm, minimum height=2.8cm, align=center},
        arrow/.style={thick, popdark, ->, >=Stealth},
        cell/.style={circle, draw, thick, minimum size=0.6cm},
        ag/.style={circle, fill=popblue, minimum size=3pt},
        ac/.style={draw, poporange, thick, -{Triangle[length=2mm]}}
    }
    
    % ÉTAPE 1 : Pénétration / Reconnaissance
    \node[stepbox, align=center] (e1) at (0,0){
        \textbf{ÉTAPE 1 : Reconnaissance} \\
        \vspace{0.2cm}
        \begin{tikzpicture}[baseline=-0.5ex, scale=0.6]
            \node[cell, fill=popgreen!30] (ph) at (0,0) {Phagocyte};
            \node[cell, fill=poppurple!30] (lb) at (2,0) {Lymph B};
            \foreach \x in {-0.3,0,0.3} \node[ag] at (\x,0.6) {};
            \draw[arrow] (ph) -- node[above, font=\tiny] {Présentation Ag} (lb);
        \end{tikzpicture} \\
        \textit{Phagocytose + présentation de l'Ag (via CMH II) au Lymphocyte B (et T4 helper).}
    };
    
    % ÉTAPE 2 : Activation / Prolifération
    \node[stepbox, align=center] (e2) at (5.5,0){
        \textbf{ÉTAPE 2 : Activation \& Prolifération} \\
        \vspace{0.2cm}
        \begin{tikzpicture}[baseline=-0.5ex, scale=0.6]
            \node[cell, fill=poppurple!30] (lb) at (0,0) {Lymph B activé};
            \node[cell, fill=poppurple!30] at (1.5,0.5) {Lymph B};
            \node[cell, fill=poppurple!30] at (1.5,-0.5) {Lymph B};
            \node[cell, fill=poppurple!30] at (2.5,0.8) {Lymph B};
            \node[cell, fill=poppurple!30] at (2.5,-0.8) {Lymph B};
            \draw[arrow] (lb) -- ++(1,0);
        \end{tikzpicture} \\
        \textit{Co-stimulation (Lymph T4 + cytokines IL-2) $\rightarrow$ Mitoses rapides = \textbf{Expansion clonale}.}
    };
    
    % ÉTAPE 3 : Différenciation / Production
    \node[stepbox, align=center] (e3) at (11,0){
        \textbf{ÉTAPE 3 : Différenciation \& Production} \\
        \vspace{0.2cm}
        \begin{tikzpicture}[baseline=-0.5ex, scale=0.6]
            \node[cell, fill=poporange!30] (pl) at (0,0) {Plasmocyte};
            \node[cell, fill=popblue!30] (mem) at (2,0) {Lymph Mémoire};
            \foreach \ang in {60, 80, 100, 120, 140, 160} \draw[ac] (pl) -- ++(\ang:0.8);
            \node[font=\tiny] at (0,-0.8) {$\sim$ 2000 Ac/s/cellule};
        \end{tikzpicture} \\
        \textit{Majorité $\rightarrow$ \textbf{Plasmocytes} (usines à Ac). Minorité $\rightarrow$ \textbf{Lymphocytes Mémoire}.}
    };
    
    % ÉTAPE 4 : Élimination
    \node[stepbox, align=center] (e4) at (16.5,0){
        \textbf{ÉTAPE 4 : Élimination de l'agent} \\
        \vspace{0.2cm}
        \begin{tikzpicture}[baseline=-0.5ex, scale=0.6]
            \node[circle, draw=popred, thick, minimum size=1cm] (bact) at (0,0) {Bactérie};
            \foreach \ang in {0,45,...,315} \draw[ac] (bact) -- ++(\ang:0.7);
            \node[cell, fill=popgreen!30] (ph) at (2.5,0) {Phagocyte};
            \draw[arrow, popgreen] (bact) -- node[above, font=\tiny] {Opsonisation} (ph);
        \end{tikzpicture} \\
        \textit{Ac fixés = \textbf{Opsonisation} $\rightarrow$ Phagocytose facilitée. Neutralisation toxines/virus. Activation Complément.}
    };
    
    % Flèches entre étapes
    \foreach \i/\j in {e1/e2, e2/e3, e3/e4} {
        \draw[arrow] (\i.east) -- (\j.west) node[midway, above, font=\bfseries] {\small Suite};
    }
    
    % Titre global
    \node[font=\bfseries\large, popdark] at (8.25, 2.2) {Les 4 étapes de la réponse immunitaire humorale (médiée par les anticorps)};
\end{tikzpicture}
\legende{Chronologie de la réponse immunitaire humorale : de la reconnaissance de l'antigène par les cellules présentatrices (APC) à l'élimination du pathogène par les effecteurs (anticorps, phagocytes, complément).}
\end{popfigure}

\courssub{La phagocytose : le nettoyage final}

La phagocytose est le mécanisme ancestral (immunité innée) qui devient ultra-efficace grâce aux anticorps (immunité adaptative). C'est le \textbf{point de rencontre} entre les deux immunités.

\begin{experience}[Observation microscopique de la phagocytose (ex. rate de souris inoculée)]
\begin{itemize}
    \item \textbf{Matériel :} Lame de frottis de rate de souris injectée avec des bactéries (ex. \emph{Staphylococcus aureus}) il y a 30 min ; coloration May-Grünwald-Giemsa ; microscope optique (G$\times$100).
    \item \textbf{Protocole :} Observer les gros noyaux lobés (neutrophiles) ou noyaux ronds excentrés (macrophages).
    \item \textbf{Observations :} 
    \begin{enumerate}
        \item \textbf{Adhérence :} La membrane du phagocyte épouse la bactérie (récepteurs Fc pour la partie constante des Ac, récepteurs du complément C3b).
        \item \textbf{Englobement :} Formation de pseudopodes $\rightarrow$ fermeture $\rightarrow$ \textbf{phagosome} (vacuole intracellulaire contenant la bactérie).
        \item \textbf{Fusion :} Le phagosome fusionne avec des \textbf{lysosomes} (riches en enzymes hydrolytiques : lysozyme, protéases, NADPH oxydase) $\rightarrow$ \textbf{phagolysosome}.
        \item \textbf{Digestion :} Milieu acide (pH $\approx$ 4,5) + enzymes $\rightarrow$ dégradation de la bactérie en acides aminés, sucres, acides gras.
        \item \textbf{Rejet :} Résidus indigestes (corps résiduels) expulsés par exocytose.
    \end{enumerate}
    \item \textbf{Interprétation :} L'anticorps agit comme un \textbf{opsonine} (marqueur « mange-moi ») qui multiplie par 10 à 100 la vitesse de phagocytose.
\end{itemize}
\end{experience}

\courssub{La mémoire immunitaire : une protection durable}

C'est la propriété la plus extraordinaire : le système immunitaire « apprend » et « retient ».

\begin{propriete}[Mémoire immunitaire (admise, illustrée par la cinétique des réponses)]
La réaction immunitaire possède une \textbf{mémoire}. Lors d'un \textbf{second contact} avec le \textbf{même antigène} :
\begin{itemize}
    \item La réponse est \textbf{plus rapide} (délai de latence raccourci de 7--10 j à 1--3 j).
    \item La réponse est \textbf{plus intense} (taux d'anticorps 10 à 100 fois plus élevé).
    \item La réponse est \textbf{plus durable} (persistance des Ac mois/années).
    \item Les anticorps produits sont d'affinité plus forte (maturation de l'affinité) et de classe IgG (vs IgM en primo-infection).
\end{itemize}
Cette mémoire est portée par les \textbf{lymphocytes B mémoire} (et T mémoire), cellules de longue durée (années, vie entière) issues de l'expansion clonale, qui ne sécrètent pas d'anticorps mais portent le récepteur (BCR) spécifique à leur surface.
\end{propriete}

\begin{popfigure}
\begin{tikzpicture}
\begin{axis}[
    width=\linewidth, height=9cm,
    xlabel={Temps (jours)}, ylabel={Taux d'anticorps (unités arbitraires)},
    xmin=0, xmax=60, ymin=1, ymax=10000,
    xtick={0,7,14,21,28,35,42,49,56},
    ytick={1,10,100,1000,10000},
    ymode=log,
    log ticks with fixed point,
    legend pos=north west,
    grid=major, grid style={popline},
    axis line style={popdark},
    tick style={popdark},
    label style={font=\small},
    tick label style={font=\small},
    legend style={font=\small, fill=popcream, draw=popdark},
    popcurve/.style={thick, mark=none},
]
% Primo-infection (IgM puis IgG)
\addplot[popblue, popcurve] coordinates {
    (0,1) (7,1) (10,10) (14,1000) (21,500) (28,200) (35,100) (42,50) (49,20) (56,5)
};
\addlegendentry{Primo-infection (1er contact) : IgM $\rightarrow$ IgG}

% Second contact (Rappel)
\addplot[poporange, popcurve] coordinates {
    (35,1) (36,10) (38,1000) (40,10000) (42,8000) (45,5000) (49,2000) (56,500)
};
\addlegendentry{Second contact (Rappel) : IgG rapide \& massif}

% Annotations
\node[popblue, font=\small, align=center] at (axis cs:14,2000) {Pic IgM $\rightarrow$ IgG};
\node[poporange, font=\small, align=center] at (axis cs:40,15000) {Pic IgG (x100)};
\draw[popmuted, dashed] (axis cs:35,1) -- (axis cs:35,10000) node[above, font=\tiny, popmuted] {Réinfection};
\end{axis}
\end{tikzpicture}
\legende{Cinétique comparative de la réponse anticorps lors du premier et du second contact avec un même antigène. Échelle logarithmique en ordonnée : la réponse secondaire est environ 100 fois plus intense et précède l'installation de la maladie.}
\end{popfigure}

\begin{exemplebox}[La varicelle : l'immunité à vie]
Après une varicelle (virus Varicella-Zoster, VZV), l'enfant guérit en 10--15 jours. Son organisme a produit des anticorps anti-VZV et, surtout, a généré des \textbf{lymphocytes B et T mémoire} spécifiques du VZV. 
\begin{itemize}
    \item Si le même virus tente de réinfecter cet enfant (ou l'adulte qu'il deviendra), les lymphocytes mémoire déclenchent une réponse secondaire en 24--48 h : le virus est éliminé \textbf{avant même l'apparition de symptômes}. C'est pourquoi on ne fait \textbf{qu'une seule fois} la varicelle dans sa vie (sauf immunodépression sévère).
    \item \textbf{Attention :} Le virus persiste latent dans les ganglions sensoriels. Chez le sujet âgé ou stressé, il peut se réactiver localement : c'est le \textbf{zona} (herpès zoster), pas une nouvelle varicelle.
\end{itemize}
\end{exemplebox}

\begin{attention}[Erreurs fréquentes au Probatoire]
\begin{enumerate}
    \item \textbf{Confondre antigène et anticorps :} L'antigène est \textbf{sur le microbe} (ou la toxine) ; l'anticorps est \textbf{dans le sang/lymphe}, produit par \textbf{nous}.
    \item \textbf{Dire « l'anticorps tue le microbe » :} L'anticorps \textbf{ne tue pas directement} (sauf activation du complément pour certaines bactéries). Il \textbf{neutralise} (bloque entrée cellulaire), \textbf{opsonise} (marque pour phagocytose), \textbf{agglutine} (regroupe). C'est le \textbf{phagocyte} qui tue et digère.
    \item \textbf{Oublier la spécificité :} Un anticorps anti-tétanos ne protège \textbf{pas} contre la diphtérie. Un vaccin protège contre \textbf{un} agent (ou quelques souches proches), pas contre « les maladies » en général.
    \item \textbf{Confondre lymphocytes mémoire et plasmocytes :} Les plasmocytes meurent en quelques semaines ; les lymphocytes mémoire persistent des décennies.
\end{enumerate}
\end{attention}

\courssub{Exercice résolu : Interpréter une courbe de réponse immunitaire}

\begin{exoresolu}[Analyse d'une sérologie post-vaccinale]
\textbf{Énoncé :} Un élève de Première reçoit un vaccin tétanos (anatoxine tétanique) à J0. Un rappel est fait à J28. On dose les anticorps anti-tétanos (IgG) dans son sang à intervalles réguliers. Résultats (taux en UI/mL) :

\begin{center}
\begin{tabular}{|c|c|c|c|c|c|c|c|c|c|}\hline
\textbf{Jour} & 0 & 7 & 14 & 21 & 28 & 35 & 42 & 49 & 56 \\ \hline
\textbf{Taux IgG} & 0 & 5 & 45 & 120 & 80 & 350 & 800 & 600 & 400 \\ \hline
\end{tabular}
\end{center}

\textbf{Questions :}
\begin{enumerate}
    \item Représenter graphiquement l'évolution du taux d'IgG (qualitatif : forme de la courbe).
    \item Identifier les phases de primo-réponse et de réponse de rappel.
    \item Expliquer pourquoi le taux à J28 (80 UI/mL) est plus élevé qu'à J0, mais que la montée est bien plus brutale après J28.
    \item Quel type de lymphocytes est responsable de cette différence ? Préciser leur devenir.
\end{enumerate}

\tcblower
\textbf{Solution détaillée :}

\begin{enumerate}
    \item \textbf{Graphique :} La courbe présente deux pics. Premier pic vers J21 (120 UI/mL), puis baisse. Second pic vers J42 (800 UI/mL), plus haut, plus précoce, plus large. (Voir courbe pgfplots ci-dessus pour le modèle théorique).
    
    \item \textbf{Phases :}
    \begin{itemize}
        \item \textbf{J0--J28 : Primo-réponse (1er contact).} Latence $\approx$ 7 j, montée lente, pic modéré (J21), descente.
        \item \textbf{J28--J56 : Réponse de rappel (2e contact).} Latence quasi nulle (montée dès J35), pic très haut (J42, 800 vs 120), maintien élevé.
    \end{itemize}
    
    \item \textbf{Explication :} À J28, l'organisme \textbf{a déjà rencontré l'antigène} (vaccin J0). Des \textbf{lymphocytes B mémoire} ont été générés lors de la primo-réponse. Ils sont plus nombreux, ont des récepteurs de plus forte affinité, et n'ont pas besoin de la phase lente d'activation/différenciation en plasmocytes. Ils se transforment en plasmocytes \emph{immédiatement} et massivement.
    
    \item \textbf{Lymphocytes concernés :} \textbf{Lymphocytes B mémoire} (CD27+). 
    \begin{itemize}
        \item Issus de l'expansion clonale de J0--J21.
        \item Ne sécrètent pas d'Ac au repos (BCR membranaire uniquement).
        \item Durée de vie : plusieurs années (voire vie entière).
        \item Au contact de l'Ag (J28) : prolifération rapide $\rightarrow$ différenciation en \textbf{plasmocytes} (usines à IgG de haute affinité) + nouveau stock de \textbf{lymphocytes mémoire} (entretien du capital).
    \end{itemize}
\end{enumerate}

\textbf{Conclusion :} Ce profil sérologique est la preuve biologique de la \textbf{mémoire immunitaire} et justifie le \textbf{calendrier vaccinal} (primovaccination + rappels) pour maintenir un taux d'anticorps protecteur ($> 0,1$ UI/mL pour le tétanos).
\end{exoresolu}

\begin{savaistu}[Le « passeport » moléculaire : le CMH]
Pour qu'un lymphocyte T reconnaisse un antigène, celui-ci doit être \textbf{présenté} par une molécule du \textbf{CMH (Complexe Majeur d'Histocompatibilité)} : 
\begin{itemize}
    \item \textbf{CMH I} (sur toutes cellules nucléées) présente des peptides \emph{intracellulaires} (virus, cancer) aux \textbf{Lymphocytes T8 cytotoxiques} (CD8+).
    \item \textbf{CMH II} (sur cellules présentatrices : macrophages, dendritiques, lymphocytes B) présente des peptides \emph{extracellulaires} (bactéries) aux \textbf{Lymphocytes T4 helpers} (CD4+).
\end{itemize}
C'est la \textbf{restriction par le CMH} : un lymphocyte T ne « voit » l'antigène que s'il est porté par \emph{son} CMH (d'où le rejet de greffe : CMH différent = « non-soi » massif). Les lymphocytes B, eux, reconnaissent l'antigène \textbf{natif} (non traité) via leur BCR (anticorps membranaire).
\end{savaistu}

\coursec{Les types d'immunité}

Dans la section précédente, nous avons vu que l'organisme réagit à la pénétration d'un \cle{antigène} en produisant des \cle{anticorps} spécifiques : c'est la \cle{réaction immunitaire}. Mais une question se pose alors naturellement : tous les individus se défendent-ils de la même façon ? Et un même individu se défend-il de la même manière à la naissance, après une maladie, ou après une vaccination ?

Prenons une observation de la vie courante au Cameroun. Dans un quartier de Yaoundé, deux enfants sont exposés au virus de la rougeole lors d'une épidémie. Le premier, qui a déjà eu la rougeole à l'âge de 3 ans, ne tombe pas malade. Le second, qui n'a jamais rencontré ce virus et n'a pas été vacciné, développe la maladie. Par ailleurs, un nourrisson de 4 mois, allaité par sa mère, traverse l'épidémie sans être atteint. Comment expliquer ces trois situations différentes ? C'est ce que nous allons comprendre en classant les \cle{types d'immunité}.

\courssub{Immunité innée et immunité acquise : les deux grandes stratégies de défense}

\textbf{Observation et questionnement.} Un enfant naît avec une peau intacte, des larmes qui lavent ses yeux, un estomac qui sécrète un liquide très acide. Avant même d'avoir rencontré le moindre microbe, il possède déjà des défenses. En revanche, la protection contre la rougeole de notre premier enfant n'existait pas à sa naissance : elle s'est construite après la maladie. Il existe donc deux grandes catégories d'immunité.

\begin{definition}[Immunité innée]
L'\cle{immunité innée} (ou immunité non spécifique) est l'ensemble des défenses présentes dès la naissance, qui agissent immédiatement contre \textbf{tout} agent étranger, sans le reconnaître de façon spécifique.
\end{definition}

L'immunité innée est la première ligne de défense de l'organisme. Elle comprend :

\begin{itemize}
\item des \textbf{barrières physiques} : la \cle{peau}, véritable muraille étanche, et les \cle{muqueuses} (revêtements internes du nez, de la bouche, des bronches, de l'intestin) qui piègent les microbes dans le mucus ;
\item des \textbf{barrières chimiques} : les \cle{larmes} et la salive contiennent une enzyme, le \cle{lysozyme}, qui détruit la paroi de nombreuses bactéries ; l'\cle{acidité gastrique} (pH voisin de 2) tue la plupart des microbes avalés avec les aliments ;
\item des \textbf{cellules phagocytes} : si un microbe franchit les barrières, des globules blancs appelés \cle{phagocytes} (polynucléaires neutrophiles, macrophages) l'englobent et le digèrent : c'est la \cle{phagocytose}, réaction rapide mais non spécifique.
\end{itemize}

\begin{definition}[Immunité acquise]
L'\cle{immunité acquise} (ou immunité adaptative) est une défense \textbf{spécifique} d'un antigène déterminé, qui se met en place \textbf{après un premier contact} avec cet antigène. Elle fait intervenir les lymphocytes et les anticorps, et elle laisse une \textbf{mémoire} immunitaire.
\end{definition}

L'immunité acquise est plus lente à se mettre en place (quelques jours lors du premier contact), mais elle est beaucoup plus efficace car elle est \emph{ciblée} : les anticorps produits contre le virus de la rougeole ne protègent pas contre la varicelle, et inversement. Surtout, elle possède une \cle{mémoire} : lors d'un second contact avec le même antigène, la réponse est plus rapide et plus intense. C'est pourquoi on ne fait la rougeole qu'une seule fois dans sa vie.

\begin{attention}[Ne pas confondre les deux niveaux de défense]
L'immunité innée est \textbf{immédiate mais non spécifique} ; l'immunité acquise est \textbf{lente à s'installer mais spécifique et durable}. Dire « les anticorps interviennent dans l'immunité innée » est une erreur classique : les anticorps relèvent toujours de l'immunité acquise. La phagocytose, en revanche, appartient à l'immunité innée (même si elle peut être facilitée par les anticorps via l'opsonisation).
\end{attention}

\begin{popfigure}
\begin{tikzpicture}[
  barriere/.style={rounded corners=3pt, draw=popdark, thick, fill=popblueL, minimum width=3.2cm, minimum height=1.5cm, align=center, font=\small},
  microbe/.style={circle, draw=poporange, fill=poporange!40, inner sep=1.5pt},
  fleche/.style={-{Stealth[length=2.5mm]}, thick, popdark}
]
% Microbes extérieurs
\node[microbe] (m1) at (0,4.6) {};
\node[microbe] (m2) at (1.1,5.0) {};
\node[microbe] (m3) at (2.2,4.5) {};
\node[microbe] (m4) at (3.3,5.0) {};
\node[font=\small\itshape, popmuted] at (1.6,5.6) {Microbes de l'extérieur};

% Barrières
\node[barriere, align=center] (peau) at (1.6,3.0){\textbf{Peau}\\ barrière physique étanche};
\node[barriere, align=center] (muq) at (6.2,3.0){\textbf{Muqueuses}\\ mucus piégeant les microbes};
\node[barriere, align=center] (secr) at (1.6,0.8){\textbf{Sécrétions}\\ larmes, salive (lysozyme)\\ suc gastrique acide};
\node[barriere, fill=popgreenL, align=center] (phag) at (6.2,0.8){\textbf{Phagocytes}\\ phagocytose des microbes\\ ayant franchi les barrières};

% Flèches bloquées
\draw[fleche, poporange] (m1) -- (0.6,3.7);
\draw[fleche, poporange] (m3) -- (2.2,3.7);
\node[font=\small, poporange] at (3.9,4.0) {passage bloqué};

\draw[fleche] (peau.east) -- (muq.west);
\draw[fleche] (peau.south) -- (secr.north);
\draw[fleche] (muq.south) -- (phag.north);
\draw[fleche] (secr.east) -- (phag.west);
\end{tikzpicture}
\legende{Les barrières naturelles de l'immunité innée : la peau, les muqueuses et les sécrétions (larmes, salive, acidité gastrique) empêchent l'entrée des microbes ; les phagocytes détruisent ceux qui ont réussi à franchir ces barrières.}
\end{popfigure}

\courssub{Immunité acquise active naturelle et immunité acquise active artificielle}

\textbf{Observation.} Revenons à notre premier enfant : il est protégé contre la rougeole parce qu'il \emph{a eu la maladie}. Mais un enfant vacciné contre la rougeole au Centre de Santé Intégré de son quartier est lui aussi protégé, sans jamais avoir été malade. Dans les deux cas, l'organisme a rencontré l'antigène et fabriqué ses propres anticorps, mais les circonstances du contact diffèrent.

\begin{itemize}
\item L'\cle{immunité acquise active naturelle} s'établit après une \textbf{maladie} : l'organisme rencontre naturellement le microbe, déclenche une réaction immunitaire, fabrique ses anticorps et garde en mémoire l'antigène. Exemple : après la rougeole, la varicelle ou les oreillons, on est en général protégé pour toute la vie.
\item L'\cle{immunité acquise active artificielle} s'établit après une \textbf{vaccination} : on introduit dans l'organisme un antigène tué, atténué ou fragmenté, incapable de provoquer la maladie mais capable de déclencher la fabrication d'anticorps et de cellules mémoires.
\end{itemize}

Dans les deux cas, l'organisme \textbf{fabrique lui-même} ses anticorps : on parle donc d'\cle{immunité active}. Cette immunité est durable, car elle s'accompagne d'une mémoire immunitaire.

\courssub{Immunité passive : la protection prêtée}

\textbf{Situation concrète.} Le nourrisson de 4 mois de notre observation n'a jamais rencontré le virus de la rougeole et n'a pas encore été vacciné (le vaccin rougeole est administré à 9 mois dans le Programme Élargi de Vaccination du Cameroun). Pourtant, il résiste. D'où vient sa protection ?

L'enquête est simple : pendant la grossesse, les anticorps de la mère (principalement des IgG) traversent le \cle{placenta} et passent dans le sang du fœtus ; après la naissance, le \cle{lait maternel}, surtout le colostrum (premier lait, épais et jaunâtre), est très riche en anticorps maternels (principalement IgA) qui protègent le tube digestif du nourrisson. Le bébé n'a \textbf{rien fabriqué lui-même} : il a \emph{reçu} des anticorps tout faits.

\begin{definition}[Immunité passive]
L'\cle{immunité passive} est une protection obtenue par la \textbf{réception d'anticorps produits par un autre individu}, sans que l'organisme les ait fabriqués lui-même. Elle est \textbf{immédiate mais temporaire}.
\end{definition}

On distingue :
\begin{itemize}
\item l'\textbf{immunité passive naturelle} : transmission des anticorps de la mère à l'enfant (placenta pendant la grossesse, lait maternel après la naissance) ;
\item l'\textbf{immunité passive artificielle} : injection d'anticorps tout faits, c'est la \cle{sérothérapie}. Exemple : la sérothérapie antivenimeuse administrée d'urgence à l'hôpital après une morsure de serpent, ou le sérum antitétanique après une blessure profonde chez une personne non vaccinée ou dont le statut vaccinal est inconnu.
\end{itemize}

\begin{exemplebox}[Le nouveau-né allaité]
Un nouveau-né allaité au sein reçoit, par le placenta puis par le lait maternel, les anticorps de sa mère contre de nombreux microbes (rougeole, tétanos, infections digestives). Il s'agit d'une \textbf{immunité passive naturelle} : immédiate, précieuse, mais qui s'épuise en quelques mois (demi-vie des IgG maternelles ~ 3 semaines). C'est l'une des raisons pour lesquelles l'Organisation Mondiale de la Santé et le Ministère de la Santé publique du Cameroun recommandent l'allaitement maternel exclusif pendant les six premiers mois de la vie.
\end{exemplebox}

\begin{attention}[L'immunité passive ne laisse pas de mémoire]
L'immunité passive \textbf{ne laisse aucune mémoire immunitaire} : les anticorps reçus sont progressivement détruits et éliminés, et la protection disparaît en quelques semaines à quelques mois. C'est pourquoi le nourrisson protégé par les anticorps maternels doit quand même être vacciné selon le calendrier vaccinal : seule la vaccination (immunité active) lui donnera une protection durable. Ne jamais écrire qu'un sérum « vaccine » : un sérum protège immédiatement mais brièvement, un vaccin protège durablement mais après un délai.
\end{attention}

\courssub{Classer les types d'immunité : une démarche rigoureuse}

\begin{methode}[Deux questions pour classer une immunité]
Face à une situation concrète (sujet d'examen, cas de la vie courante), poser successivement deux questions :
\begin{enumerate}
\item \textbf{L'organisme fabrique-t-il lui-même ses anticorps ?}
  \begin{itemize}
  \item Oui $\rightarrow$ immunité \textbf{active} (contact avec l'antigène : maladie ou vaccination) ;
  \item Non, il reçoit des anticorps tout faits $\rightarrow$ immunité \textbf{passive} (mère-enfant ou sérothérapie).
  \end{itemize}
\item \textbf{Le contact avec l'antigène ou la réception des anticorps est-elle naturelle ou provoquée par l'intervention humaine ?}
  \begin{itemize}
  \item Maladie, allaitement, grossesse $\rightarrow$ \textbf{naturelle} ;
  \item Vaccination, injection de sérum $\rightarrow$ \textbf{artificielle}.
  \end{itemize}
\end{enumerate}
On obtient alors quatre combinaisons pour l'immunité acquise : active naturelle (maladie), active artificielle (vaccination), passive naturelle (anticorps maternels), passive artificielle (sérothérapie). L'immunité innée, elle, est présente dès la naissance et ne dépend d'aucun contact préalable.
\end{methode}

\begin{popfigure}
\begin{tikzpicture}[
  grow=right,
  level 1/.style={sibling distance=4.8cm, level distance=3.4cm},
  level 2/.style={sibling distance=2.4cm, level distance=3.6cm},
  level 3/.style={sibling distance=1.3cm, level distance=3.8cm},
  every node/.style={rounded corners=2pt, draw=popdark, thick, align=center, font=\small, inner sep=3pt, text width=2.2cm},
  edge from parent/.style={draw=popmuted, thick}
]
\node[fill=popblueL] {Immunité}
  child { node[fill=popgreenL, align=center]{Acquise\\ (spécifique, mémoire)}
    child { node[fill=popgreenL, align=center]{Passive\\ (anticorps reçus,\\ pas de mémoire)}
      child { node[fill=white, align=center]{Artificielle\\ sérothérapie} }
      child { node[fill=white, align=center]{Naturelle\\ anticorps maternels} }
    }
    child { node[fill=popgreenL, align=center]{Active\\ (anticorps fabriqués,\\ mémoire)}
      child { node[fill=white, align=center]{Artificielle\\ vaccination} }
      child { node[fill=white, align=center]{Naturelle\\ après une maladie} }
    }
  }
  child { node[fill=poporangeL, align=center]{Innée\\ (présente à la naissance,\\ non spécifique)\\ peau, muqueuses, phagocytes} };
\end{tikzpicture}
\legende{Arbre de classification des types d'immunité : l'immunité innée est non spécifique ; l'immunité acquise est spécifique et se divise en immunité active (l'organisme fabrique ses anticorps, mémoire) et passive (il reçoit des anticorps, pas de mémoire), chacune pouvant être naturelle ou artificielle.}
\end{popfigure}

\begin{center}
\begin{tabularx}{\linewidth}{|l|X|X|X|}
\hline
\rowcolor{popblueL} \textbf{Critère} & \textbf{Immunité innée} & \textbf{Immunité active} & \textbf{Immunité passive} \\
\hline
Origine & Présente dès la naissance & Contact avec l'antigène (maladie ou vaccin) & Anticorps reçus d'un autre individu \\
\hline
Spécificité & Non spécifique & Spécifique d'un antigène & Spécifique (anticorps reçus) \\
\hline
Fabrication d'anticorps & Non & Oui, par l'organisme & Non, anticorps reçus tout faits \\
\hline
Rapidité d'action & Immédiate & Lente (quelques jours au 1er contact) & Immédiate \\
\hline
Durée & Toute la vie & Longue (mémoire immunitaire) & Courte (quelques mois, pas de mémoire) \\
\hline
Exemples & Peau, larmes, acidité gastrique, phagocytose & Rougeole guérie (naturelle) ; vaccination (artificielle) & Lait maternel, placenta (naturelle) ; sérum antivenimeux (artificielle) \\
\hline
\end{tabularx}
\end{center}

\begin{exoresolu}[Classer des situations concrètes]
\textbf{Énoncé.} Pour chacune des situations suivantes, préciser le type d'immunité mis en jeu, en justifiant : (a) un enfant guéri de la varicelle ne la contracte plus ; (b) un bébé de 2 mois, allaité, est protégé contre certaines diarrhées ; (c) une personne mordue par un serpent reçoit un sérum antivenimeux à l'hôpital de district ; (d) une fillette vaccinée contre le tétanos est protégée durablement.
\tcblower
\textbf{Solution détaillée.}
\begin{itemize}
\item (a) L'enfant a rencontré naturellement le virus et a fabriqué ses propres anticorps et des cellules mémoires : il s'agit d'une \textbf{immunité acquise active naturelle}.
\item (b) Le bébé n'a rien fabriqué : il a reçu les anticorps de sa mère par le lait : \textbf{immunité passive naturelle}. Elle disparaîtra en quelques mois.
\item (c) Le malade reçoit des anticorps tout faits, produits par un animal immunisé : \textbf{immunité passive artificielle} (sérothérapie antivenimeuse). Elle agit immédiatement mais ne laisse pas de mémoire. Par ailleurs, la morsure étant une blessure profonde, une prophylaxie tétanique (sérum antitétanique et/ou rappel vaccinal) sera décidée selon le carnet de vaccination.
\item (d) Le vaccin a provoqué la fabrication d'anticorps et de cellules mémoires par l'organisme lui-même, sans maladie : \textbf{immunité acquise active artificielle}.
\end{itemize}
\end{exoresolu}

\begin{aretenir}[L'essentiel de la section]
\begin{itemize}
\item L'\textbf{immunité innée}, présente dès la naissance, est non spécifique : peau, muqueuses, sécrétions (lysozyme, acidité gastrique) et phagocytes.
\item L'\textbf{immunité acquise} est spécifique et dotée de mémoire ; elle est \textbf{active} quand l'organisme fabrique ses anticorps (naturelle après une maladie, artificielle après une vaccination) et \textbf{passive} quand il reçoit des anticorps tout faits (naturelle : mère vers enfant ; artificielle : sérothérapie).
\item L'immunité passive est immédiate mais temporaire et sans mémoire ; seule l'immunité active protège durablement.
\end{itemize}
\end{aretenir}

Maintenant que nous savons distinguer ces types d'immunité, nous allons étudier plus en détail, dans la section suivante, l'outil qui permet de provoquer artificiellement une immunité active sans risquer la maladie : la \cle{vaccination}.

\coursec{La vaccination, moyen de renforcement de l'immunité}

Dans la section précédente, nous avons distingué les \cle{types d'immunité} : l'immunité naturelle (innée ou acquise après une infection) et l'immunité artificielle. Nous avons noté qu'attraper une maladie permet parfois de devenir immunisé... mais au prix de souffrances, voire de la mort. Comment obtenir cette immunité durable \emph{sans} passer par la maladie ? C'est toute l'idée de la \cle{vaccination}, l'une des plus grandes victoires de la médecine.

\courssub{Situation-problème : pourquoi vacciner un enfant en bonne santé ?}

\begin{exemplebox}[Une observation à l'hôpital de district]
Au centre de santé de Bafoussam, une maman hésite : « Mon bébé n'est pas malade, pourquoi lui faire une piqûre ? » Quelques mois plus tard, une épidémie de rougeole touche le quartier. Parmi les enfants hospitalisés, la grande majorité n'était pas vaccinée ; les enfants vaccinés, eux, sont restés indemnes ou n'ont eu qu'une forme très légère. Comment une injection faite \emph{avant} la maladie a-t-elle pu protéger ces enfants ?
\end{exemplebox}

\textbf{Hypothèse de travail :} l'injection contiendrait quelque chose qui « apprend » au système immunitaire à reconnaître le microbe, sans provoquer la maladie. Vérifions cette hypothèse en définissant rigoureusement la vaccination, puis en observant les données.

\courssub{Définition et principe de la vaccination}

\begin{definition}[Vaccin et vaccination]
Un \cle{vaccin} est une préparation contenant un \cle{antigène inoffensif} (microbe tué, microbe vivant atténué, \cle{anatoxine} -- toxine inactivée -- ou simple fragment du microbe) qui entraîne le système immunitaire à fabriquer des \cle{anticorps} et des \cle{cellules mémoires}, sans provoquer la maladie. La \cle{vaccination} est l'acte d'introduire ce vaccin dans l'organisme (par injection, ou parfois par voie orale comme pour le vaccin antipoliomyélitique).
\end{definition}

Le principe repose sur un point capital étudié dans les sections précédentes : la \cle{mémoire immunitaire}. Le vaccin \emph{simule une première infection} :

\begin{enumerate}
\item L'antigène du vaccin est reconnu comme étranger par les lymphocytes.
\item Une \cle{réaction immunitaire primaire} se déclenche : production d'anticorps spécifiques (réponse lente, en quelques jours à quelques semaines) et fabrication de \cle{lymphocytes mémoires} (lymphocytes B et T mémoires).
\item L'antigène étant inoffensif, la maladie ne se déclare pas (au pire, une légère fièvre passagère, signe de l'activation immunitaire).
\item Lors d'un \emph{vrai} contact ultérieur avec le microbe pathogène, les cellules mémores déclenchent une \cle{réponse secondaire} : rapide (quelques heures à 2--3 jours), massive et efficace. Le microbe est détruit avant de provoquer la maladie.
\end{enumerate}

\begin{propriete}[Immunité conférée par la vaccination (admise)]
La vaccination confère une \cle{immunité acquise artificielle} durable grâce à la mémoire immunitaire : l'organisme vacciné répond à une infection réelle par une réponse secondaire plus rapide et plus intense que la réponse primaire d'un organisme non vacciné. Cette propriété est admise ; on doit savoir la justifier par la courbe des taux d'anticorps.
\end{propriete}

\begin{experience}[Observation des données : la courbe qui prouve la mémoire]
\textbf{Protocole :} on mesure régulièrement le taux d'anticorps antitétanos dans le sang de deux groupes de personnes. Groupe A : personnes vaccinées contre le tétanos (primo-vaccination effectuée), mises en contact ultérieurement avec la toxine tétanique (rappel ou infection). Groupe B : personnes non vaccinées, infectées pour la première fois.

\textbf{Observations :} dans le groupe B, le taux d'anticorps monte lentement et faiblement (réponse primaire) ; la maladie a le temps de se déclarer. Dans le groupe A, le taux d'anticorps s'élève très vite et atteint un niveau élevé : la toxine est neutralisée avant tout dégât.

\textbf{Interprétation :} la vaccination a créé des cellules mémoires qui accélèrent et amplifient la réponse.

\textbf{Conclusion :} le vaccin ne protège pas en agissant directement contre le microbe, mais en \emph{préparant} le système immunitaire.
\end{experience}

\begin{popfigure}
\begin{center}
\begin{tikzpicture}[scale=0.95, every node/.style={font=\small}]
% Courbe : taux d'anticorps
\begin{scope}
\draw[->, thick, popdark] (0,0) -- (9.6,0) node[below left]{Temps (semaines)};
\draw[->, thick, popdark] (0,0) -- (0,4.4) node[above, align=center]{Taux\\d'anticorps};
% réponse primaire (non vacciné)
\draw[very thick, poporange, domain=0:2.6, samples=60] plot (\x, {1.6*(1-exp(-1.6*\x))});
\draw[very thick, poporange, domain=2.6:4.2, samples=40] plot (\x, {1.6*exp(-0.9*(\x-2.6))});
% réponse secondaire (vacciné)
\draw[very thick, popblue, domain=4.2:5.4, samples=60] plot (\x, {3.9*(1-exp(-2.2*(\x-4.2)))});
\draw[very thick, popblue, domain=5.4:9.2, samples=40] plot (\x, {3.9*exp(-0.35*(\x-5.4))});
% marqueurs
\draw[dashed, popmuted] (0.6,0) -- (0.6,0.95);
\node[poporange, below, align=center] at (0.6,0) {Infection\\(non vacciné)};
\draw[dashed, popmuted] (4.2,0) -- (4.2,0.9);
\node[popblue, below, align=center] at (4.2,0) {Infection\\(vacciné)};
\node[poporange, align=center] at (2.1,2.1) {Réponse primaire\\lente et faible};
\node[popblue, align=center] at (7.3,4.05) {Réponse secondaire\\rapide et forte};
\end{scope}
\end{tikzpicture}
\end{center}
\legende{Comparaison des taux d'anticorps après infection chez un individu non vacciné (réponse primaire, lente et faible) et chez un individu vacciné (réponse secondaire, rapide et intense, grâce aux cellules mémoires). L'axe des abscisses représente le temps en semaines ; l'axe des ordonnées le taux d'anticorps (unités arbitraires).}
\end{popfigure}

\begin{popfigure}
\begin{center}
\begin{tikzpicture}[font=\small, box/.style={draw, rounded corners=2pt, align=center, minimum height=1cm, text width=3.1cm}]
% Colonne non vacciné
\node[box, fill=poporangeL, draw=poporange, align=center] (a1) at (0,0){\textbf{Individu non vacciné}\\Premier contact avec le microbe};
\node[box, fill=poporangeL, draw=poporange, below=0.5cm of a1, align=center] (a2){Réponse primaire lente\\(anticorps en 1 à 3 semaines)};
\node[box, fill=poporange!40, draw=poporange, below=0.5cm of a2, align=center] (a3){Le microbe se multiplie\\$\Rightarrow$ \textbf{maladie déclarée},\\parfois grave ou mortelle};
% Colonne vacciné
\node[box, fill=popgreenL, draw=popgreen, align=center] (b1) at (7.2,0){\textbf{Individu vacciné}\\Vaccin = antigène inoffensif};
\node[box, fill=popgreenL, draw=popgreen, below=0.5cm of b1, align=center] (b2){Fabrication d'anticorps\\et de \textbf{cellules mémoires}};
\node[box, fill=popgreen!40, draw=popgreen, below=0.5cm of b2, align=center] (b3){Infection réelle : réponse\\secondaire rapide $\Rightarrow$\\microbe détruit, \textbf{pas de maladie}};
\draw[-{Stealth}, thick, poporange] (a1) -- (a2);
\draw[-{Stealth}, thick, poporange] (a2) -- (a3);
\draw[-{Stealth}, thick, popgreen] (b1) -- (b2);
\draw[-{Stealth}, thick, popgreen] (b2) -- (b3);
\end{tikzpicture}
\end{center}
\legende{Comparaison du destin d'un organisme non vacciné et d'un organisme vacciné face à un même agent pathogène : seul le vacciné possède une mémoire immunitaire prête à réagir.}
\end{popfigure}

\begin{attention}[Un vaccin ne soigne pas, il prévient]
Erreur fréquente au Probatoire : écrire que « le vaccin guérit la maladie ». Faux ! Un vaccin n'a \emph{aucun effet} sur une maladie déjà déclarée, car la réponse immunitaire qu'il induit demande des jours à des semaines. La vaccination est un acte de \cle{prévention}, réalisé sur une personne \emph{saine}, avant tout contact avec le microbe. De même, le vaccin ne contient pas d'anticorps : c'est l'organisme qui les fabrique lui-même (immunité active).
\end{attention}

\courssub{Intérêt individuel et intérêt collectif de la vaccination}

\textbf{Intérêt individuel.} La vaccination protège la personne vaccinée contre des maladies graves, parfois mortelles ou laissant de lourdes séquelles : la \cle{rougeole} (encore responsable d'épidémies au Cameroun), le \cle{tétanos} (contracté par une simple blessure souillée de terre), la \cle{poliomyélite} (paralysies irréversibles), la \cle{tuberculose} et la \cle{fièvre jaune} (dont le vaccin est exigé pour voyager dans de nombreux pays).

\textbf{Intérêt collectif.} Un microbe se transmet de personne à personne. Si la plupart des individus sont vaccinés, le microbe ne trouve presque plus d'hôte sensible : il \emph{circule difficilement} et l'épidémie s'éteint. On parle de \cle{protection de la communauté} (ou immunité collective). Ainsi, se faire vacciner, c'est aussi protéger ceux qui ne peuvent pas l'être : les nourrissons trop jeunes, les malades dont le système immunitaire est affaibli.

\begin{savaistu}[La variole, une maladie effacée de la surface de la Terre]
La variole a tué des millions de personnes au cours de l'histoire. Grâce à une campagne mondiale de vaccination coordonnée par l'Organisation mondiale de la Santé, le dernier cas naturel a été détecté en 1977 et la maladie a été officiellement déclarée \textbf{éradiquée en 1980}. C'est la première — et à ce jour la seule — maladie humaine éliminée par l'action de l'Homme. La vaccination est donc capable de faire disparaître une maladie de la planète entière.
\end{savaistu}

\courssub{Le calendrier vaccinal au Cameroun : le Programme Élargi de Vaccination (PEV)}

Depuis 1976, le Cameroun applique le \cle{Programme Élargi de Vaccination} (PEV), qui offre gratuitement aux enfants les vaccins essentiels, administrés dans les hôpitaux et centres de santé, dès la naissance. Le carnet de santé de l'enfant enregistre chaque vaccin reçu.

\begin{center}
\begin{tabularx}{\linewidth}{|l|l|X|}
\hline
\rowcolor{popblueL} \textbf{Âge} & \textbf{Vaccin} & \textbf{Maladie(s) visée(s)} \\
\hline
À la naissance & BCG & Tuberculose \\
\hline
À la naissance & Polio 0 (voie orale) & Poliomyélite \\
\hline
6, 10 et 14 semaines & Pentavalent & Diphtérie, tétanos, coqueluche, hépatite B, infections à \emph{Haemophilus influenzae} b \\
\hline
6, 10 et 14 semaines & Polio (voie orale) & Poliomyélite \\
\hline
9 mois & Rougeole & Rougeole \\
\hline
9 mois & Fièvre jaune & Fièvre jaune \\
\hline
\end{tabularx}
\end{center}

\begin{aretenir}[L'essentiel du PEV]
Le PEV camerounais protège gratuitement les enfants contre la tuberculose (BCG), la poliomyélite, la diphtérie, le tétanos, la coqueluche, l'hépatite B, la rougeole et la fièvre jaune. Le respect complet du calendrier (toutes les doses, aux âges prévus) est indispensable à une protection efficace.
\end{aretenir}

\courssub{Primo-vaccination et rappels de vaccins}

La mémoire immunitaire peut s'atténuer avec le temps : le nombre de cellules mémoires diminue et le taux d'anticorps baisse. Pour certaines maladies, une seule injection ne suffit pas à établir une mémoire solide ; on pratique alors une \cle{primo-vaccination} en plusieurs doses (ex. : trois doses du pentavalent à 6, 10 et 14 semaines) pour « booster » progressivement la réponse. Par la suite, des \cle{rappels} (injections ultérieures du même vaccin, espacées de plusieurs années) \emph{stimulent à nouveau} la mémoire immunitaire. Chaque rappel provoque une réponse secondaire : rapide, forte, et qui reconstitue le stock de cellules mémoires. C'est pourquoi le rappel du vaccin antitétanique (souvent associé à la diphtérie et la coqueluche : DTC) est recommandé régulièrement tout au long de la vie (ex. : à 15-18 mois, 6 ans, puis tous les 10-20 ans chez l'adulte).

\courssub{Complément exigible : vaccination et sérothérapie}

\begin{definition}[Sérothérapie]
La \cle{sérothérapie} consiste à injecter à un malade des \cle{anticorps déjà fabriqués} (prélevés dans le sang d'une personne ou d'un animal immunisé, ou produits par génie génétique). L'action est \emph{immédiate} mais \emph{temporaire}, car ces anticorps sont détruits en quelques semaines et l'organisme n'a fabriqué aucune cellule mémoire (immunité passive artificielle).
\end{definition}

\begin{center}
\begin{tabularx}{\linewidth}{|l|X|X|}
\hline
\rowcolor{popblueL} \textbf{Critère} & \textbf{Vaccination} & \textbf{Sérothérapie} \\
\hline
Ce qu'on injecte & Antigène inoffensif (microbe tué, atténué, anatoxine, fragment) & Anticorps déjà fabriqués (immunoglobulines) \\
\hline
Qui fabrique les anticorps ? & L'organisme lui-même (immunité active) & Un autre organisme (donneur) -- immunité passive \\
\hline
Début d'action & Lent (jours à semaines) & Immédiat \\
\hline
Durée de protection & Longue (années, vie entière grâce à la mémoire) & Courte (quelques semaines) \\
\hline
But & Prévenir la maladie & Traiter d'urgence (morsure de serpent, tétanos déclaré, rage, diphtérie) \\
\hline
\end{tabularx}
\end{center}

\begin{methode}[Justifier l'intérêt d'un vaccin en trois idées]
Pour rédiger une justification complète au Probatoire, articule toujours :
\begin{enumerate}
\item \textbf{Antigène inoffensif} : le vaccin introduit un antigène incapable de provoquer la maladie (préciser le type : atténué, inactivé, anatoxine, sous-unité) ;
\item \textbf{Mémoire immunitaire} : l'organisme fabrique des anticorps et surtout des cellules mémoires (lymphocytes B et T mémoires) ;
\item \textbf{Réponse rapide} : lors de l'infection réelle, la réponse secondaire détruit le microbe avant l'apparition de la maladie.
\end{enumerate}
Conclure en citant l'intérêt individuel (protection efficace) et collectif (blocage de la circulation du microbe, protection des non-vaccinables).
\end{methode}

\begin{exoresolu}[Vaccin ou sérum ?]
Énoncé. Un enfant non vacciné contre le tétanos se blesse profondément au pied avec un clou rouillé. Le médecin lui injecte un sérum antitétanique, puis programme une vaccination. (1) Pourquoi le sérum est-il injecté en urgence ? (2) Pourquoi faut-il malgré tout vacciner ensuite ?
\tcblower
Solution. (1) La blessure souillée risque d'introduire le bacille du tétanos (\emph{Clostridium tetani}) ; la maladie peut se déclarer avant que l'organisme ne fabrique ses propres anticorps. Le sérum apporte des anticorps antitétaniques \emph{déjà fabriqués} (immunoglobulines), dont l'action est immédiate : ils neutralisent la toxine tétanique tout de suite. (2) Mais ces anticorps injectés disparaissent en quelques semaines et l'enfant n'a aucune cellule mémoire. La vaccination, en introduisant l'antigène tétanique inoffensif (anatoxine tétanique), pousse son organisme à fabriquer ses propres anticorps et des cellules mémoires, garantissant une protection durable contre les blessures futures.
\end{exoresolu}

\begin{exercice}[À toi de jouer]
Lors d'une épidémie de rougeole dans un lycée de Yaoundé, on constate que 90 \% des cas sont des élèves non vaccinés. (1) Explique, en trois étapes, pourquoi les élèves vaccinés sont protégés. (2) Explique pourquoi une forte couverture vaccinale protège aussi les rares élèves non vaccinés.
\end{exercice}

\begin{corrige}[Exercice]
(1) Le vaccin antirougeoleux a introduit un antigène inoffensif (virus atténué) ; l'organisme a fabriqué des anticorps et des cellules mémoires ; au contact du virus sauvage, la réponse secondaire, rapide et intense, détruit le virus avant que la maladie ne se déclare. (2) Quand la grande majorité des élèves est vaccinée, le virus ne trouve presque plus d'hôte sensible : il circule difficilement, la chaîne de transmission est interrompue, et même les élèves non vaccinés ont peu de chances de rencontrer le virus. C'est la protection de la communauté (immunité collective).
\end{corrige}

\begin{aretenir}[Bilan de la section]
La vaccination introduit un antigène inoffensif (microbe tué, atténué, anatoxine ou fragment) qui provoque une réaction immunitaire avec mémoire, sans provoquer la maladie : c'est une immunité acquise artificielle, durable, à visée préventive. Elle protège l'individu et, par la réduction de la circulation du microbe, toute la communauté (immunité collective). Le PEV camerounais organise cette protection dès la naissance (BCG, polio, pentavalent, rougeole, fièvre jaune), complétée par des rappels tout au long de la vie. Contrairement à la sérothérapie, qui apporte des anticorps tout faits à action immédiate mais brève (immunité passive), la vaccination fait travailler notre propre système immunitaire (immunité active).
\end{aretenir}

\coursec{Facteurs affaiblissant le système immunitaire}

Dans la section précédente, nous avons vu que la vaccination permet de \emph{renforcer} l'immunité en préparant l'organisme à se défendre contre un agent pathogène précis. Mais tout système de défense, aussi bien entraîné soit-il, peut être affaibli. Au Cameroun, les services de santé le constatent chaque jour : deux personnes exposées au même microbe ne tombent pas toujours malades de la même façon. L'une guérit en quelques jours, l'autre enchaîne les infections. Pourquoi ? Parce que l'efficacité du système immunitaire dépend de nombreux facteurs : l'alimentation, le stress, la fatigue, certaines maladies, l'âge ou encore le mode de vie. Cette section recense ces facteurs, explique leurs mécanismes et montre leurs conséquences sur la santé.

\courssub{Observation et problème de santé}

Dans un centre de santé de Yaoundé, le médecin remarque que certains patients reviennent très souvent : angines à répétition, diarrhées persistantes, plaies qui guérissent mal. En les interrogeant, il constate des points communs : alimentation pauvre en protéines, fatigue chronique, parfois une infection par le VIH non traitée. L'hypothèse se dessine : \textbf{le système immunitaire de ces patients est affaibli}.

\begin{definition}[Immunodépression]
On dit qu'un organisme est en état d'\cle{immunodépression} (ou immunodéficience) lorsque ses défenses immunitaires sont diminuées, de sorte qu'il devient plus vulnérable aux infections. L'immunodépression peut être temporaire (fatigue, stress, carence alimentaire) ou durable (infection par le VIH non traitée, certaines maladies chroniques).
\end{definition}

\courssub{La malnutrition : des défenses privées de matériaux}

\textbf{Situation concrète.} Dans certaines zones rurales, des enfants souffrant de malnutrition (alimentation essentiellement constituée de bouillie de maïs, sans protéines ni légumes) meurent de maladies bénignes comme la rougeole ou une simple diarrhée, alors que des enfants bien nourris guérissent facilement.

\textbf{Explication.} Les cellules immunitaires (lymphocytes, phagocytes) et les anticorps sont fabriqués à partir de \cle{protéines}. Les anticorps sont eux-mêmes des protéines (les immunoglobulines). Une alimentation carencée en protéines réduit donc directement la \emph{production} des cellules et des molécules de défense. De plus, certaines vitamines et certains minéraux sont indispensables au bon fonctionnement de ces cellules :

\begin{center}
\begin{tabularx}{\linewidth}{|l|X|X|}
\hline
\rowcolor{popblueL} \textbf{Nutriment} & \textbf{Rôle dans l'immunité} & \textbf{Sources alimentaires locales} \\
\hline
Protéines & Fabrication des anticorps et des cellules immunitaires & Viande, poisson, œufs, haricots niébé, arachides \\
\hline
Vitamine A & Intégrité des muqueuses (barrières naturelles), maturation des lymphocytes & Carotte, patate douce, huile de palme rouge, foie \\
\hline
Vitamine C & Activité des phagocytes, protection des cellules contre le stress oxydatif & Agrumes, goyave, mangue, légumes verts \\
\hline
Fer & Fonctionnement des cellules immunitaires (attention : excès défavorable) & Viande, folong, épinards, légumineuses \\
\hline
Zinc & Développement et activation des lymphocytes & Viande, poisson, graines, céréales complètes \\
\hline
\end{tabularx}
\end{center}

\begin{aretenir}[Malnutrition et immunité]
Les carences en \cle{protéines}, en \cle{vitamines} (notamment A et C) et en \cle{minéraux} (fer, zinc) réduisent la \textbf{production} et l'\textbf{efficacité} des cellules immunitaires. La malnutrition est ainsi la première cause d'immunodépression dans le monde, et elle aggrave la mortalité des maladies infectieuses de l'enfant.
\end{aretenir}

\courssub{Le stress prolongé et la fatigue : des hormones qui freinent les défenses}

\textbf{Observation.} Beaucoup d'élèves remarquent qu'ils attrapent plus facilement un rhume ou une angine en période de révisions intenses, lorsqu'ils dorment peu et s'alimentent mal. Ce n'est pas un hasard.

\textbf{Mécanisme.} En cas de stress, l'organisme sécrète des hormones, dont le \cle{cortisol} (par les glandes surrénales). À petite dose et sur une courte durée, le cortisol est utile : il mobilise l'énergie. Mais lorsque le stress devient \emph{prolongé}, le taux de cortisol reste élevé en permanence, et cette hormone \textbf{diminue l'activité des défenses immunitaires} : elle réduit la production et l'efficacité des lymphocytes. La fatigue et le manque de sommeil aggravent ce phénomène, car c'est pendant le sommeil que l'organisme « répare » et renouvelle une partie de ses défenses.

\begin{attention}[Nuancer le rôle du stress]
Le stress des examens \emph{seul} ne rend pas malade : un stress bref et modéré est même stimulant. C'est un \textbf{stress prolongé}, combiné à la \textbf{fatigue}, au \textbf{manque de sommeil} et à une \textbf{mauvaise alimentation}, qui diminue réellement les défenses. La bonne stratégie avant le Probatoire : dormir suffisamment, manger équilibré et s'accorder des pauses.
\end{attention}

\courssub{Certaines maladies qui affaiblissent les défenses}

\courssub{Le VIH : destruction des lymphocytes T4 (CD4+)}

Le cas le plus spectaculaire est celui du \cle{VIH} (Virus de l'Immunodéficience Humaine). Ce virus s'attaque précisément aux \cle{lymphocytes T4} (ou lymphocytes T auxiliaires CD4+), les cellules qui coordonnent toute la réaction immunitaire : ce sont elles qui « donnent l'alerte » et dirigent la production d'anticorps et l'action des autres cellules tueuses.

\textbf{Évolution de l'infection non traitée.} Après la contamination, le VIH se multiplie dans les lymphocytes T4 et les détruit progressivement. Pendant plusieurs années, la personne peut ne présenter aucun symptôme (phase dite asymptomatique), mais le nombre de lymphocytes T4 diminue lentement. Lorsque ce taux passe sous un seuil critique, les défenses s'effondrent : c'est le stade du \cle{SIDA} (Syndrome d'ImmunoDéficience Acquise), caractérisé par l'apparition d'infections graves.

\begin{exemplebox}[Un exemple chiffré à connaître]
Chez un adulte sain, le taux de lymphocytes T4 est d'environ \num{800} cellules par mm\textsuperscript{3} de sang (valeurs normales : 500 à \num{1500}/mm\textsuperscript{3}). Chez une personne séropositive \textbf{non traitée}, ce taux chute progressivement. En dessous de \textbf{200 cellules/mm\textsuperscript{3}}, les défenses sont si affaiblies que des \emph{infections opportunistes} apparaissent : c'est le seuil qui définit le stade SIDA.
\end{exemplebox}

\begin{popfigure}
\begin{center}
\begin{tikzpicture}
\begin{axis}[
    ybar,
    width=0.95\linewidth,
    height=7cm,
    bar width=0.8cm,
    ymin=0, ymax=1000,
    xlabel={Temps écoulé depuis la contamination (années)},
    ylabel={Lymphocytes T4 (cellules/mm\textsuperscript{3})},
    xtick={0,2,4,6,8,10},
    xticklabels={0, 2, 4, 6, 8, 10},
    ytick={0,200,400,600,800,1000},
    axis line style={popline},
    tick label style={font=\small},
    label style={font=\small},
    every axis plot/.append style={fill=popblue, draw=popdark},
    nodes near coords,
    nodes near coords style={font=\footnotesize, popdark},
]
\addplot coordinates {(0,800) (2,650) (4,500) (6,350) (8,220) (10,120)};
\draw[dashed, thick, poporange] (0,200) -- (10,200) node[pos=0.02, above, font=\footnotesize, poporange, anchor=west]{Seuil 200/mm\textsuperscript{3} : risque d'infections opportunistes};
\end{axis}
\end{tikzpicture}
\end{center}
\legende{Évolution du taux de lymphocytes T4 chez une personne infectée par le VIH \textbf{non traitée} (données fictives simplifiées). Le taux passe d'environ 800/mm\textsuperscript{3} à moins de 200/mm\textsuperscript{3} en quelques années : c'est le stade du SIDA. Avec un traitement antirétroviral bien suivi, le taux de T4 se maintient et cette chute n'a pas lieu.}
\end{popfigure}

\begin{definition}[Infection opportuniste]
Une \cle{infection opportuniste} est une infection due à un microbe habituellement inoffensif (présent dans l'environnement ou sur notre corps sans nous rendre malades), qui se développe lorsque les défenses immunitaires sont affaiblies. Exemples : tuberculose, pneumonies à \emph{Pneumocystis}, candidoses (muguet), méningites à cryptocoque, diarrhées prolongées.
\end{definition}

\begin{savaistu}[Pourquoi le VIH est-il si redoutable ?]
Le VIH s'attaque précisément aux \textbf{lymphocytes T4}, les « chefs d'orchestre » de la réaction immunitaire. En détruisant les coordinateurs, il paralyse toute l'armée de défense, même si les autres soldats (phagocytes, lymphocytes B) sont encore présents. Heureusement, les traitements antirétroviraux (ARV), distribués gratuitement au Cameroun, bloquent la multiplication du virus : une personne sous traitement peut vivre longtemps avec un système immunitaire fonctionnel. D'où l'importance capitale du \textbf{dépistage précoce}.
\end{savaistu}

\courssub{Le paludisme et le diabète}

\textbf{Le paludisme}, première cause de consultation au Cameroun, affaiblit les défenses de plusieurs façons : le parasite (\emph{Plasmodium}) détruit les globules rouges, provoquant une \emph{anémie} qui diminue l'apport en oxygène aux cellules ; les crises répétées épuisent l'organisme. Un enfant paludéen fréquent devient plus vulnérable aux autres infections.

\textbf{Le diabète} (taux de sucre trop élevé dans le sang) perturbe le fonctionnement des phagocytes : ceux-ci se déplacent moins bien (chimiotaxie altérée) et détruisent moins efficacement les microbes. C'est pourquoi les diabétiques souffrent plus souvent d'infections de la peau et des plaies qui guérissent mal, notamment au niveau des pieds.

\courssub{Autres facteurs : âge, alcool, tabac, manque de sommeil}

\begin{itemize}
\item \textbf{L'âge} : chez le \emph{nourrisson}, le système immunitaire est encore immature (il est partiellement protégé par les anticorps de la mère, transmis pendant la grossesse et par l'allaitement). Chez la \emph{personne âgée}, les défenses s'affaiblissent naturellement (c'est l'\cle{immunosénescence}) : d'où l'importance de la vaccination contre la grippe chez les personnes âgées.
\item \textbf{L'alcool} : consommé régulièrement et en excès, il diminue l'activité des phagocytes et des lymphocytes, et aggrave souvent la malnutrition.
\item \textbf{Le tabac} : il détruit les \cle{cils vibratiles} des bronches (barrière naturelle qui évacuent les microbes) et réduit l'efficacité des défenses respiratoires : les fumeurs ont plus d'infections des voies respiratoires.
\item \textbf{Le manque de sommeil} : dormir moins de 6 heures par nuit de façon prolongée diminue la production de cellules immunitaires et d'anticorps.
\end{itemize}

\begin{popfigure}
\begin{center}
\begin{tikzpicture}[
    facteur/.style={rectangle, rounded corners=3pt, draw=popdark, fill=popblueL, text width=3.2cm, align=center, font=\small, inner sep=4pt},
    centre/.style={circle, draw=popdark, fill=poporangeL, text width=2.8cm, align=center, font=\small\bfseries, inner sep=4pt}
]
\node[centre] (c) {Affaiblissement du système immunitaire};
\node[facteur] (f1) at (90:3.6cm) {Malnutrition (carences en protéines, vitamines A et C, fer, zinc)};
\node[facteur] (f2) at (30:3.9cm) {Stress prolongé, fatigue, manque de sommeil (cortisol)};
\node[facteur] (f3) at (-30:3.9cm) {Maladies : VIH/SIDA, paludisme, diabète};
\node[facteur] (f4) at (-90:3.6cm) {Âge : nourrissons et personnes âgées};
\node[facteur] (f5) at (-150:3.9cm) {Alcool et tabac};
\node[facteur] (f6) at (150:3.9cm) {Conséquences : infections fréquentes, graves, opportunistes};
\foreach \f in {f1,f2,f3,f4,f5,f6} \draw[thick, popmuted] (c) -- (\f);
\end{tikzpicture}
\end{center}
\legende{Diagramme en étoile résumant les principaux facteurs qui affaiblissent le système immunitaire et leur conséquence commune : la multiplication des infections.}
\end{popfigure}

\courssub{Conséquences d'un système immunitaire affaibli}

\begin{aretenir}[Conséquences de l'immunodépression]
Un système immunitaire affaibli se traduit par des infections :
\begin{itemize}
\item plus \textbf{fréquentes} (rhumes, angines, diarrhées à répétition) ;
\item plus \textbf{graves} (une infection bénigne peut évoluer en complication sérieuse) ;
\item plus \textbf{longues} (guérison lente, plaies qui cicatrisent mal) ;
\item parfois \textbf{opportunistes}, c'est-à-dire causées par des microbes normalement inoffensifs, en particulier chez les personnes vivant avec le VIH non traité.
\end{itemize}
\end{aretenir}

\courssub{Lien avec la situation camerounaise : préserver son immunité}

Au Cameroun, deux combats de santé publique contribuent directement à préserver l'immunité des populations :

\begin{enumerate}
\item \textbf{La lutte contre le paludisme} : utilisation des moustiquaires imprégnées d'insecticide (distribuées lors des campagnes nationales), assainissement des eaux stagnantes, traitement rapide des crises. Moins de paludisme, c'est moins d'anémie et des défenses préservées, surtout chez les enfants et les femmes enceintes.
\item \textbf{Le dépistage du VIH} : gratuit et confidentiel dans les centres de santé, il permet de découvrir l'infection tôt et de commencer les antirétroviraux avant que les défenses ne s'effondrent. Une personne dépistée tôt et bien traitée garde un taux de T4 élevé et échappe aux infections opportunistes.
\end{enumerate}

\begin{exoresolu}[Analyse d'un cas clinique]
\textbf{Énoncé.} M. Ndi, 42 ans, consulte pour la troisième fois en six mois : pneumonie, puis diarrhée persistante, puis muguet (candidose de la bouche). Il a perdu 8 kg en un an. Le médecin prescrit une numération des lymphocytes T4 : le résultat est de 150 cellules/mm\textsuperscript{3}.
\begin{enumerate}
\item Pourquoi dit-on que les infections de M. Ndi sont « opportunistes » ?
\item Quelle hypothèse diagnostique le médecin doit-il vérifier, et par quel examen ?
\item Quel seuil du taux de T4 justifie cette inquiétude ?
\end{enumerate}
\tcblower
\textbf{Solution détaillée.}
\begin{enumerate}
\item Ces infections sont dites opportunistes car elles sont causées par des microbes habituellement inoffensifs ou bien combattus chez une personne saine (champignon de la candidose, microbes des voies respiratoires). Leur répétition et leur gravité traduisent un affaiblissement des défenses immunitaires.
\item Le médecin doit rechercher une infection par le VIH, virus qui détruit les lymphocytes T4. L'examen est le \textbf{test de dépistage du VIH} (recherche d'anticorps anti-VIH dans le sang), réalisé avec le consentement du patient.
\item Le seuil critique est de \textbf{200 cellules/mm\textsuperscript{3}} : en dessous, le risque d'infections opportunistes est élevé et on parle de stade SIDA. M. Ndi, avec 150 cellules/mm\textsuperscript{3}, est en dessous de ce seuil. La prise en charge par les antirétroviraux permettra de faire remonter son taux de T4 et de restaurer ses défenses.
\end{enumerate}
\end{exoresolu}

\begin{exercice}[À toi de jouer]
Madame Essomba se demande pourquoi son fils de 3 ans, mal nourri, attrape « toutes les maladies qui passent », alors que sa fille de 8 ans, bien nourrie, est rarement malade. En utilisant les notions de cette section, propose une explication scientifique et deux conseils pratiques pour renforcer l'immunité du garçon.
\end{exercice}

\begin{corrige}[Exercice : À toi de jouer]
\textbf{Explication.} Le garçon souffre probablement de malnutrition : les carences en protéines réduisent la fabrication des anticorps et des cellules immunitaires, et les carences en vitamines A et C, en fer et en zinc diminuent l'efficacité de ces cellules. Son système immunitaire, déjà immature à 3 ans, est donc doublement affaibli, ce qui explique la fréquence de ses infections.

\textbf{Conseils pratiques.}
\begin{enumerate}
\item Améliorer son alimentation : introduire des protéines (poisson, œufs, haricots niébé, arachides), des fruits et légumes riches en vitamines A et C (carotte, patate douce, mangue, agrumes) et des sources de fer et de zinc (folong, viande).
\item Veiller aux autres facteurs de protection : sommeil suffisant, vaccination à jour, lutte contre le paludisme (moustiquaire imprégnée) et consultation rapide en cas d'infection.
\end{enumerate}
\end{corrige}

\begin{aretenir}[L'essentiel de la section]
\begin{itemize}
\item La \textbf{malnutrition} (carences en protéines, vitamines A et C, fer, zinc) réduit la production et l'efficacité des cellules immunitaires.
\item Le \textbf{stress prolongé}, la \textbf{fatigue} et le \textbf{manque de sommeil} agissent via des hormones comme le cortisol, qui diminue l'activité des défenses.
\item Le \textbf{VIH} détruit les lymphocytes T4 (CD4+) ; sous 200 cellules/mm\textsuperscript{3}, les infections opportunistes apparaissent (SIDA). Le \textbf{paludisme} et le \textbf{diabète} affaiblissent aussi les défenses.
\item L'\textbf{âge} (nourrissons, personnes âgées), l'\textbf{alcool} et le \textbf{tabac} sont des facteurs aggravants.
\item Conséquence : des infections plus fréquentes, plus graves et plus longues. Dépistage du VIH et lutte contre le paludisme sont des priorités de santé publique au Cameroun.
\end{itemize}
\end{aretenir}

\coursec{Comportements favorables au renforcement de l'immunité}

Dans la section précédente, nous avons identifié les facteurs qui affaiblissent le système immunitaire : malnutrition, stress prolongé, manque de sommeil, infections chroniques (paludisme, VIH/sida). La relation est réversible : des comportements adaptés entretiennent et renforcent nos défenses. Cette section traduit nos connaissances en \cle{comportements concrets de santé}, applicables au quotidien au Cameroun avec nos ressources locales.

\courssub{Une alimentation équilibrée et variée, carburant des défenses immunitaires}

\textbf{Observation de départ.} Dans un même quartier de Yaoundé, deux enfants du même âge présentent des profils de santé différents : l'un enchaîne diarrhées et rhumes, l'autre est rarement malade. Les enquêtes épidémiologiques (ex. : enquêtes DHS-MICS au Cameroun) confirment que la malnutrition (sous-nutrition ou carences en micronutriments) est corrélée à une morbidité infectieuse plus élevée et plus grave.

\textbf{Hypothèse et interprétation.} Les cellules effectrices de l'immunité (lymphocytes T et B, phagocytes : neutrophiles, macrophages) et les molécules effectrices (anticorps ou immunoglobulines, cytokines, protéines du complément) sont synthétisées à partir d'acides aminés (apport \cle{protéique}). Leur prolifération, leur différenciation et leur activité enzymatique requièrent des \cle{vitamines} (A, C, D, E, B6, B9, B12) et des \cle{oligo-éléments} (zinc, sélénium, fer, cuivre). Une alimentation carencée prive l'organisme de la « matière première » et des cofacteurs de ses défenses.

\begin{aretenir}[Principes d'une alimentation favorable à l'immunité au Cameroun]
\begin{itemize}
\item \cle{Fruits et légumes} à chaque repas (au moins 400 g/jour selon l'OMS) : sources majeures de vitamines (C, A, B9), de minéraux et de fibres. Valorisons nos productions locales : \emph{folong}, \emph{njama-njama}, \emph{okok}, éru, amaranthe, carottes, tomates, poivrons ; mangues, oranges, goyaves, papayes, bananes, ananas, fruits du safoutier.
\item \cle{Protéines} d'origine animale et/ou végétale chaque jour : poisson (machoiron, carpe, capitaine, poisson fumé/séché), viande, œufs, lait ; haricots niébé, arachides, soja, graines de courge. Les anticorps sont des glycoprotéines : sans apport protéique suffisant (environ 0,8 à 1 g/kg/jour), l'hyperglobulinémie de défense est impossible.
\item \cle{Céréales complètes et tubercules} (maïs, riz complet, mil, sorgho, manioc, igname, macabo, patate douce) : apportent l'énergie (glucides complexes), des vitamines du groupe B et du zinc, indispensables à la prolifération lymphocytaire.
\item \cle{Eau potable} en quantité suffisante (1,5 à 2 L/jour, plus en cas de fièvre ou chaleur) : maintient l'hydratation des muqueuses (barrière physique), assure le transport lymphatique et sanguin, élimine les déchets métaboliques.
\item Limiter les excès de graisses saturées, de sucres ajoutés et de sel, facteurs de risque des maladies métaboliques (diabète de type 2, hypertension artérielle) qui installent une inflammation chronique de bas grade et affaiblissent l'immunité à long terme.
\end{itemize}
\end{aretenir}

\begin{popfigure}
\begin{center}
\begin{tikzpicture}[scale=1.0]
% Pyramide alimentaire standard (base énergétique) avec focus immunité
\fill[popblueL] (0,0) -- (6,0) -- (4.8,1.6) -- (1.2,1.6) -- cycle; % Base: Céréales
\fill[popgreenL] (1.2,1.6) -- (4.8,1.6) -- (3.9,2.9) -- (2.1,2.9) -- cycle; % Niveau 1: Fruits/Légumes
\fill[popgoldL] (2.1,2.9) -- (3.9,2.9) -- (3.3,4.0) -- (2.7,4.0) -- cycle; % Niveau 2: Protéines
\fill[poppinkL] (2.7,4.0) -- (3.3,4.0) -- (3,4.9) -- cycle; % Sommet: Graisses/Sucres
\draw[popdark, thick] (0,0) -- (3,4.9) -- (6,0) -- cycle;
\draw[popdark] (1.2,1.6) -- (4.8,1.6);
\draw[popdark] (2.1,2.9) -- (3.9,2.9);
\draw[popdark] (2.7,4.0) -- (3.3,4.0);
% Étiquettes
\node[align=center, font=\small] at (3,0.8) {\textbf{Céréales \& tubercules} \\ maïs, riz, manioc, igname, mil \\ \textit{(Énergie, Vit B, Zn)}};
\node[align=center, font=\small] at (3,2.25) {\textbf{Fruits \& Légumes} \\ folong, njama-njama, mangues, oranges \\ \textit{(Vit A, C, B9, antioxydants)}};
\node[align=center, font=\small] at (3,3.45) {\textbf{Protéines} \\ poisson, haricots, arachides, œufs \\ \textit{(Acides aminés, Fer, Zn)}};
\node[align=center, font=\small] at (3,4.35) {\textbf{Graisses, Sucres, Sel} \\ \textit{Modération}};
% Flèche eau
\draw[->, popblue, very thick] (6.6,0.3) -- (6.6,1.5);
\node[align=center, font=\small, text=popblue] at (7.6,0.9) {Eau potable\\ chaque jour};
\end{tikzpicture}
\end{center}
\legende{Pyramide alimentaire de référence (adaptée au contexte camerounais). La base apporte l'énergie nécessaire à la prolifération cellulaire ; les deux étages suivants concentrent les micronutriments et protéines spécifiquement requis par la réponse immunitaire. L'eau est transversale.}
\end{popfigure}

\begin{methode}[Justifier un comportement nutritionnel par son mécanisme immunologique]
Au Probatoire, une affirmation vague (« manger des fruits renforce l'immunité ») ne suffit pas. Suis la chaîne \textbf{Comportement $\rightarrow$ Nutriment/Substance $\rightarrow$ Mécanisme cellulaire/moléculaire} :
\begin{enumerate}
\item \textbf{Identifier} le comportement précis (ex. : consommer des goyaves ou des oranges quotidiennement).
\item \textbf{Nommer} le nutriment clé apporté (ex. : vitamine C, acide ascorbique).
\item \textbf{Relier} ce nutriment à un mécanisme immunitaire démontré (ex. : la vitamine C stimule la chimiotaxie et la phagocytose des neutrophiles, favorise la prolifération des lymphocytes B et T, et protège les membranes cellulaires contre le stress oxydatif généré par le burst respiratoire des phagocytes).
\end{enumerate}
\textbf{Exemple de rédaction attendue :} « La consommation régulière de goyaves, riches en vitamine C, améliore la chimiotaxie et l'activité microbicide des neutrophiles, premières cellules intervenant lors de l'inflammation aiguë. »
\end{methode}

\courssub{Respecter le calendrier vaccinal (PEV) : l'entraînement ciblé du système immunitaire}

La vaccination est la seule méthode scientifiquement validée pour induire une \cle{mémoire immunitaire} spécifique sans risque de maladie grave. Elle repose sur l'activation des lymphocytes B et T CD4$^+$ naïfs, leur différenciation en cellules effectrices (plasmocytes sécréteurs d'anticorps, lymphocytes T cytotoxiques) et, surtout, en \cle{lymphocytes mémoires} (B mémoires et T mémoires) à longue durée de vie.

\begin{aretenir}[Deux règles d'or pour une protection durable]
\begin{itemize}
\item \textbf{Compléter le calendrier du PEV (Programme Élargi de Vaccination) aux âges requis} : BCG (tuberculose), Polio (OPV/IPV), Pentavalent (Diphtérie-Tétanos-Coqueluche + Hépatite B + \textit{Haemophilus influenzae} type b), Pneumocoque (PCV), Rotavirus, Rougeole-Rubéole (RR), Fièvre jaune, HPV (filles 9-14 ans). Ces vaccins sont \textbf{gratuits} dans les centres de santé publics et hôpitaux de district.
\item \textbf{Effectuer les rappels} : la concentration d'anticorps et le nombre de lymphocytes mémoires déclinent avec le temps. Le rappel (ex. : DTC à 18 mois, 6 ans, puis dT tous les 10-20 ans ; fièvre jaune à 10 ans) provoque une \cle{réponse anamnestique} (secondaire) : plus rapide, plus intense, de plus haute affinité (affinité mature), assurant une protection continue.
\end{itemize}
\end{aretenir}

\begin{attention}[Piège fréquent au Probatoire]
Un calendrier \emph{incomplet} (doses primaires manquantes ou rappels oubliés) ne confère qu'une protection partielle et transitoire. Conserve le \textbf{carnet de vaccination} (ou carte de santé de l'enfant) et présente-le à \emph{chaque} contact avec le système de santé (consultation, campagne de rattrapage). Ne pas confondre « vaccin » (le produit) et « vaccination » (l'acte complet incluant les rappels).
\end{attention}

\courssub{Sommeil suffisant et gestion du stress chronique}

Le sommeil n'est pas une pause passive : c'une phase active de \cle{régulation neuro-immunitaire}. Pendant le sommeil lent profond, la sécrétion de cortisol (immunosuppresseur) est minimale, tandis que celle de l'hormone de croissance (GH) et de la prolactine (immunostimulantes) augmente. Cela favorise la production de cytokines pro-inflammatoires (IL-12, IFN-$\gamma$) nécessaires à l'immunité cellulaire (Th1) et la redistribution des cellules T vers les ganglions lymphatiques. Un élève de Première (16-18 ans) a besoin de \textbf{8 à 10 heures} de sommeil par nuit (recommandations National Sleep Foundation / OMS).

Le \cle{stress chronique} (angoisse scolaire, conflits familiaux, insécurité) maintient une hypercortisolémie. Le cortisol (glucocorticoïde) se lie à son récepteur intracellulaire dans les leucocytes et inhibe la transcription de gènes codant pour des cytokines (IL-2, IFN-$\gamma$), réduit la prolifération lymphocytaire et induit l'apoptose de lymphocytes thymiques et périphériques. Résultat : immunosupression relative.

\begin{exemplebox}[Cumul de facteurs affaiblissants : le piège des compositions]
Pendant la session d'examens, un élève dort 5 h/nuit, saute le petit-déjeuner, vit en angoisse permanente. Il cumule : (1) \textbf{privation de sommeil} $\rightarrow$ baisse IL-12/IFN-$\gamma$, (2) \textbf{déficit nutritionnel aigu} $\rightarrow$ manque substrats synthèse anticorps, (3) \textbf{stress chronique} $\rightarrow$ cortisolémie élevée $\rightarrow$ lymphopénie transitoire. Résultat fréquent : surinfection bactérienne (angine) ou réactivation virale (herpès labial) en pleine épreuve. La performance scolaire \emph{nécessite} la préservation de l'immunité.
\end{exemplebox}

\begin{propriete}[Stratégies de gestion du stress compatibles avec les études]
\begin{itemize}
\item \textbf{Planification} : répartir les révisions sur le long terme (méthode J-30, J-15, J-7) pour éviter le « bachotage » final.
\item \textbf{Activité physique aérobie modérée} (marche rapide 30 min, footing, football) : augmente la circulation des cellules NK et des lymphocytes T, réduit le cortisol.
\item \textbf{Hygiène du sommeil} : horaires fixes coucher/lever, obscurité, pas d'écran 1 h avant le coucher (lumière bleue inhibe mélatonine).
\item \textbf{Soutien social} : discuter avec un pair, un enseignant, un conseiller d'orientation, la famille.
\end{itemize}
\end{propriete}

\courssub{Hygiène corporelle et alimentaire : réduire la charge antigénique}

Le principe est \cle{préventif} : limiter l'entrée des agents pathogènes préserve le capital immunitaire pour les agressions inévitables.

\begin{aretenir}[Gestes barrières quotidiens, validés scientifiquement]
\begin{itemize}
\item \cle{Lavage des mains} à l'eau et au savon (ou solution hydroalcoolique si pas d'eau) : \textbf{avant} de manger/nourrir un enfant/préparer les repas ; \textbf{après} les toilettes, le nettoyage d'un enfant, les transports en commun, le contact avec des animaux/déchets. Le savon solubilise l'enveloppe lipidique des virus (grippe, coronavirus, VIH, hépatites B/C) et décolle mécaniquement bactéries et kystes parasitaires.
\item \cle{Cuisson thorough} (à cœur, $> 70^\circ$C) des viandes, poissons, œufs : détruit bactéries (\textit{Salmonella}, \textit{E. coli}, \textit{Campylobacter}), parasites (\textit{Taenia}, \textit{Ascaris}) et virus.
\item Laver fruits/légumes à l'eau potable (ajouter quelques gouttes d'eau de Javel diluée ou vinaigre si eau non traitée) ; boire de l'eau traitée (ébullition 1 min, pastilles chlore, filtre céramique, bouteille scellée).
\item Protection des aliments (couvertures, réfrigération) contre les mouches (vecteurs mécaniques de \textit{Shigella}, \textit{Vibrio cholerae}, kystes d'amibes).
\end{itemize}
\end{aretenir}

\begin{savaistu}[Preuve d'efficacité : le lavage des mains]
Des essais randomisés en grappes (cluster-RCT) menés au Pakistan, au Népal et au Ghana ont démontré que la promotion du lavage des mains au savon réduit l'incidence de la diarrhée de \textbf{30 à 48\%} et des infections respiratoires aiguës de \textbf{16 à 23\%} chez l'enfant de moins de 5 ans. Au Cameroun, la stratégie « École Amie de l'Hygiène » (UNICEF/MINESEC) a réduit l'absentéisme scolaire lié aux maladies féco-orales. C'est l'intervention de santé publique au \textbf{ratio coût-efficacité le plus élevé}.
\end{savaistu}

\courssub{Activité physique régulière ; éviter l'alcool et le tabac}

L'\cle{activité physique} modérée et régulière (au moins 150 min/semaine d'endurance modérée ou 75 min intense, OMS 2020) améliore la \cle{circulation hémato-lymphatique} : les leucocytes patrouillent plus efficacement dans les tissus (surveillance immunitaire). Elle réduit l'inflammation chronique (baisse CRP-us, IL-6, TNF-$\alpha$) et le stress oxydatif.

À l'inverse, le \cle{tabac} (cigarette, chicha, tabac à priser) est un toxique direct sur l'immunité respiratoire :
\begin{itemize}
\item Destruction des \cle{cils vibratiles} de l'épithélium respiratoire (mucociliary clearance) $\rightarrow$ stagnation du mucus et des microbes.
\item Inhibition de la fonction macrophagique alvéolaire (phagocytose, présentation d'antigène).
\item Inflammation chronique des voies aériennes (bronchite chronique) $\rightarrow$ porte d'entrée facilitée pour \textit{Streptococcus pneumoniae}, \textit{Haemophilus influenzae}, virus grippaux.
\end{itemize}
L'\cle{alcool} (même modéré, et surtout en \emph{binge drinking} / « biture ») :
\begin{itemize}
\item Déprime la fonction neutrophile (chimiotaxie, phagocytose, burst oxydatif).
\item Perturbe l'immunité adaptative (réduction lymphocytes T CD4$^+$ et CD8$^+$, altération réponse vaccinale).
\item Atteint le foie (site majeur de synthèse des protéines de la phase aiguë : CRP, fibrinogène, complément).
\end{itemize}
\textbf{Savoir-être :} savoir refuser la pression de groupe (« un petit verre ne tue pas », « tire une bouffée ») est une compétence de santé majeure.

\courssub{Prévenir les maladies qui épuisent durablement le système immunitaire}

Certaines pathologies endémiques au Cameroun détruisent activement l'architecture ou la fonction immunitaire. Leur prévention \emph{est} une stratégie de renforcement de l'immunité.

\begin{itemize}
\item \textbf{Paludisme (\textit{Plasmodium falciparum})} : chaque accès palustre détruit des globules rouges, active massivement les macrophages (hémophagocytose), induit une dysrégulation cytokinique (tempête inflammatoire) et épuise la rate. \textbf{Prévention} : \cle{Moustiquaire imprégnée d'insecticide à longue durée d'action (MILDA)} chaque nuit, toute l'année ; traitement préventif intermittent (TPI) chez la femme enceinte (SP) ; assainissement (évacuation eaux stagnantes, larvicidage). \emph{Chaque crise évitée = capital immunitaire préservé.}
\item \textbf{Infection à VIH (Virus de l'Immunodéficience Humaine)} : le VIH-1 (groupe M au Cameroun) infecte et détruit sélectivement les \cle{lymphocytes T CD4$^+$} (chefs d'orchestre de l'immunité adaptative : aide aux lymphocytes B pour commutation de classe/affinité, activation macrophages, activation LT CD8$^+$). Sans traitement, chute des CD4 $< 200/\mu$L $\rightarrow$ Sida (infections opportunistes : tuberculose, pneumocystose, toxoplasmose, cryptococcose). \textbf{Prévention combinée} : \cle{fidélité} mutuelle testée, \cle{préservatif} masculin/féminin à chaque rapport à risque, \cle{dépistage} volontaire gratuit (TROD ou ELISA) en centre de santé, \cle{PrEP} (prophylaxie pré-exposition) pour populations à risque élevé, \cle{TASP} (traitement comme prévention : charge virale indétectable = intransmissible). Un dépistage précoce permet l'initiation rapide des ARV (gratuit au Cameroun) qui préserve l'immunité.
\item \textbf{Infections bactériennes répétées non traitées} (angines à streptocoque, otites, furoncles, pneumonies) : peuvent signaler un déficit immunitaire primitif (rare) ou secondaire (malnutrition, VIH, diabète). \textbf{Conduite à tenir} : \cle{consultation précoce} au centre de santé (pas d'automédication antibiotique anarchique qui sélectionne les résistances), bilan biologique si récidives (NFS, dosage Ig, sérologie VIH, recherche diabète).
\end{itemize}

\begin{popfigure}
\begin{center}
\begin{tikzpicture}[scale=1.0]
% Deux colonnes comparatives
% Colonne gauche : Risques
\draw[fill=poppinkL, draw=popdark, rounded corners=5pt] (0,0) rectangle (5.6,7.5);
\node[font=\bfseries, text=popdark] at (2.8,7.2) {Comportements \`a risque (Affaiblissent)};
\node[align=left, font=\small, text width=5.0cm] at (2.8,3.8) {
$\bullet$ Alimentation monotone, carenc\'ee (sauter repas)\\[2pt]
$\bullet$ Sommeil $<$ 8 h, stress non g\'er\'e\\[2pt]
$\bullet$ Calendrier vaccinal incomplet, rappels oubli\'es\\[2pt]
$\bullet$ Consommation d'alcool, tabac (chicha, cigarette)\\[2pt]
$\bullet$ Nuits sans MILDA (moustiquaire)\\[2pt]
$\bullet$ Rapports sexuels non prot\'eg\'es, multi-partenaires\\[2pt]
$\bullet$ Autom\'edication, consultation tardive
};
% Colonne droite : Protecteurs
\draw[fill=popgreenL, draw=popdark, rounded corners=5pt] (6.4,0) rectangle (12,7.5);
\node[font=\bfseries, text=popdark] at (9.2,7.2) {Comportements protecteurs (Renforcent)};
\node[align=left, font=\small, text width=5.0cm] at (9.2,3.8) {
$\bullet$ Alimentation vari\'ee : c\'er\'eales, l\'egumes, fruits, prot\'eines locales\\[2pt]
$\bullet$ 8--10 h sommeil, planning r\'evisions, sport, pauses\\[2pt]
$\bullet$ PEV complet \` jour (carnet pr\'esent\'e \` chaque visite)\\[2pt]
$\bullet$ Refus alcool/tabac, choix pairs sains\\[2pt]
$\bullet$ MILDA chaque nuit + assainissement domicile\\[2pt]
$\bullet$ Fid\'elit\'e test\'ee, pr\'eservatif, d\'epistage VIH r\'egulier\\[2pt]
$\bullet$ Consultation pr\'ecoce centre de sant\'e, observance traitements
};
% Flèche centrale
\draw[->, very thick, poporange] (5.7,3.8) -- (6.3,3.8);
\node[font=\bfseries, text=poporange, rotate=90] at (6.0, 3.0) {CHOIX};
\end{tikzpicture}
\end{center}
\legende{Bilan comparatif : chaque comportement à risque (rose) a son pendant protecteur (vert). Le passage de l'un à l'autre est un choix individuel et collectif, soutenu par l'environnement (accès eau, vaccins gratuits, MILDA, centres de santé).}
\end{popfigure}

\begin{attention}[Vigilance : aucun produit « miracle » ou « boosteur »]
Méfie-toi des allégations marketing (réseaux sociaux, marchés, pharmacies de rue) vantant des « boosters d'immunité » (poudres, sirops, compléments alimentaires, décoctions « secrètes ») censés agir en quelques jours. \textbf{Preuve scientifique :} aucun produit isolé ne « booste » un système immunitaire sain ; l'immunité est un système complexe régulé par l'homéostasie. Seule une \cle{hygiène de vie globale} (nutrition, sommeil, hygiène, activité, stress, vaccination, prévention maladies endémiques) maintient la compétence immunitaire. Dépenser ses ressources dans ces produits détourne souvent des vrais gestes efficaces (achat fruits/légumes, transport vers centre de santé).
\end{attention}

\begin{exoresolu}[Analyser une situation et proposer des mesures correctives argumentées]
\textbf{Énoncé.} Marc, 17 ans, élève de Première A, consulte l'infirmière scolaire pour « fatigue permanente et rhumes à répétition (3 ce trimestre) ». L'anamnèse révèle : coucher 1 h du matin / lever 6 h (5 h sommeil), petit-déjeuner sauté, déjeuner « pain-beurre » ou sauté, dîner unique (riz/sauce peu de légumes), 3-4 cigarettes/jour avec pairs, carnet vaccinal : dernière dose DTC à 6 ans, pas de MILDA à la maison (fenêtres sans moustiquaires). Propose \textbf{quatre} comportements prioritaires pour Marc, en justifiant \textbf{chacun} par un mécanisme immunologique précis.
\tcblower
\textbf{Solution détaillée.}
\begin{enumerate}
\item \textbf{Allonger le temps de sommeil à 8--9 h/nuit (coucher 22 h max).} 
\textit{Justification :} Le sommeil lent profond est la phase de pic de sécrétion d'IL-12 et d'IFN-$\gamma$ (polarisation Th1) et de redistribution des lymphocytes T CD4$^+$ et CD8$^+$ vers les organes lymphoïdes secondaires. La restriction de sommeil à 5 h augmente le cortisol et réduit la réponse anticorps post-vaccinale de 50\% (données expérimentales).
\item \textbf{Rétablir 3 repas équilibrés/jour incluant fruits/légumes et protéines à chaque repas (ex. : papaye/petit-déj, folong/poisson/déjeuner, orange/dîner).} 
\textit{Justification :} Fournir les acides aminés indispensables à la synthèse des immunoglobulines (anticorps) par les plasmocytes, et les cofacteurs enzymatiques (Zn, Se, Vit A, C, D, B9) pour la prolifération clonale des lymphocytes et l'activité microbicide des phagocytes (burst respiratoire).
\item \textbf{Arrêter totalement le tabac (cigarette/chicha).} 
\textit{Justification :} La fumée inhibe le battement ciliaire de l'épithélium trachéobronchique (clairance mucociliaire $\approx$ 0), paralyse les macrophages alvéolaires (phagocytose $\downarrow$, présentation d'antigène $\downarrow$) et crée une inflammation chronique neutrophile stérile qui « distrait » le système immunitaire. Risque multiplié par 2 à 4 pour pneumonie et méningite à pneumocoque.
\item \textbf{Se présenter au CPS (Centre de Promotion de la Santé) pour rattrapage vaccinal : rappel dT (diphtérie-tétanos) + vérification Fièvre Jaune / Hépatite B / Rougeole.} 
\textit{Justification :} Le dernier rappel DTC date de 11 ans (6 ans $\rightarrow$ 17 ans) : le taux d'anticorps anti-toxine tétanique est probablement $< 0,01$ UI/mL (seuil protecteur 0,1 UI/mL). Le rappel réactive les lymphocytes B mémoires $\rightarrow$ différenciation en plasmocytes $\rightarrow$ montée exponentielle des IgG anti-toxine (réponse anamnestique en 3-5 jours). Protection retrouvée pour 10-20 ans.
\end{enumerate}
\textbf{Note pédagogique :} Chaque proposition suit le schéma \textbf{Comportement $\rightarrow$ Cible moléculaire/cellulaire $\rightarrow$ Conséquence fonctionnelle}. C'est l'exigence du Probatoire.
\end{exoresolu}

\begin{exercice}[À toi de jouer -- Argumentation vaccinale]
Ta cousine de 10 ans, en classe de CM2, n'a pas reçu sa 2\up{ème} dose de vaccin Rougeole-Rubéole (RR) prévue à 15-18 mois, ni le vaccin HPV (recommandé filles 9-14 ans). Sa mère hésite : « Elle est grande maintenant, elle a déjà eu la rougeole bébé (peut-être), et le HPV c'est pour plus tard ». Rédige un argumentaire structuré (10-15 lignes) pour la convaincre, en mobilisant : la notion de \cle{mémoire immunitaire} et de \cle{réponse anamnestique}, la gravité des complications (encéphalite post-rougeole, cancer du col utérin), la \cle{gratuité} au PEV, et la notion de \cle{protection collective} (immunité de groupe).
\end{exercice}

\begin{corrige}[Exercice -- Argumentaire type]
« Maman, le fait qu'elle soit "grande" ne la protège pas. La rougeole contractée bébé (si c'en était une) n'a pas garanti une mémoire immunitaire solide sans la 2\up{ème} dose : un seul contact antigénique (maladie ou 1 dose) laisse un taux d'anticorps qui chute sous le seuil protecteur en quelques années. La 2\up{ème} dose RR provoque une \textbf{réponse anamnestique} : les lymphocytes B mémoires se réactivent, produisent massivement des IgG de haute affinité en 2-3 jours, assurant une protection \textbf{durable (vie entière pour la plupart)}. Sans ce rappel, elle risque une rougeole de l'adolescent/adulte, bien plus grave (pneumonie, encéphalite 1/1000, mortalité). Pour le HPV, vacciner \textbf{avant} tout début d'activité sexuelle (9-14 ans) induit le taux d'anticorps anti-VPH le plus élevé, prévenant l'infection persistante par les génotypes 16/18 (70\% cancers du col). Les deux vaccins sont \textbf{gratuits} au centre de santé avec sa carte. Enfin, la vacciner protège aussi ses petits frères non encore vaccinés et les enfants immunodéprimés de l'école (\textbf{immunité de groupe}). C'est un acte de responsabilité, gratuit et sans risque majeur. »
\end{corrige}

\begin{aretenir}[Synthèse : Les 4 piliers d'une immunité compétente au quotidien]
\begin{enumerate}
\item \textbf{Nutrition immunonutritive} : Assiette variée (céréales + légumes + fruits + protéines locales) + eau potable. \textit{Mécanisme : substrats synthèse (AA), cofacteurs enzymatiques (Zn, Se, Vit), énergie (ATP).}
\item \textbf{Entraînement vaccinal à jour} : PEV complet + rappels. \textit{Mécanisme : constitution/maintenance pool lymphocytes B/T mémoires spécifiques.}
\item \textbf{Régulation neuro-immunitaire} : Sommeil 8-10 h + gestion stress (sport, planning, soutien). \textit{Mécanisme : profil cytokinique Th1 favorable, cortisol bas, prolifération lymphocytaire préservée.}
\item \textbf{Prévention des agressions évitables} : Hygiène mains/aliments (barrière), MILDA (paludisme), Préservatif/Dépistage/PrEP (VIH), Non-tabac/Non-alcool (toxiques directs), Consultation précoce (éviter séquelles).
\end{enumerate}
\textbf{Rappel :} \cle{Aucun produit miracle} ne remplace cette hygiène de vie globale. La compétence immunitaire se construit et s'entretient \emph{chaque jour}.
\end{aretenir}

\newpage

\coursec{Activité d'intégration}

\coursec{Leçon NA04 : Amélioration de la santé du système immunitaire}

\courssub{Objectifs de l'activité}

À l'issue de cette activité d'intégration, tu seras capable de :
\begin{itemize}
\item mobiliser les savoirs de la leçon (organes et cellules du système immunitaire, notions d'antigène et d'anticorps, types d'immunité) pour analyser une situation réelle de santé communautaire ;
\item \cle{expliquer} le principe général d'une réaction immunitaire et le rôle de la mémoire immunitaire, à l'aide d'un schéma annoté ;
\item \cle{justifier} l'intérêt de la vaccination par un argumentaire structuré et scientifiquement correct ;
\item \cle{proposer} des comportements favorables au renforcement de l'immunité, en reliant chaque mesure à un mécanisme biologique précis ;
\item travailler en équipe, argumenter à l'oral et défendre un point de vue devant un auditoire, comme le ferait un conseiller municipal de santé.
\end{itemize}

\courssub{Situation}

\textbf{Commune de Bafouam, quartier de Tougang (Ouest-Cameroun).}

Lors de la réunion trimestrielle de la chefferie traditionnelle du quartier de Tougang, le chef et les responsables du comité de santé dressent un constat inquiétant : depuis le début de l'année scolaire, le centre de santé intégré (CSI) de la localité voisine enregistre une \cle{recrudescence} des cas de \textbf{rougeole} et de \textbf{diarrhées aiguës} chez les enfants de moins de 10 ans. Sur les 40 enfants vus en consultation le mois dernier, 22 présentaient une rougeole (dont 15 n'étaient pas vaccinés) et 18 une diarrhée avec déshydratation. Deux enfants ont dû être évacués vers l'hôpital de district de Bafouam.

Le comité de santé dispose des données suivantes, recueillies par le relais communautaire du quartier :

\begin{center}
\begin{tabularx}{\linewidth}{|l|X|}
\hline
\rowcolor{popblueL} \textbf{Donnée} & \textbf{Valeur observée dans le quartier} \\
\hline
Taux de vaccination des enfants (rougeole, vaccin VAR) & 55 \% (contre 95 \% recommandés par l'OMS) \\
\hline
Alimentation dominante des ménages & Repas quotidiens à base de maïs et de manioc, pauvres en protéines ; viande ou poisson moins de 2 fois par semaine \\
\hline
Approvisionnement en eau & Un seul point d'eau (forage) pour tout le quartier ; eau non traitée, non bouillie avant consommation \\
\hline
Accès aux soins & Centre de santé intégré situé à 3 km, accessible à pied ; consultations pédiatriques le mardi et le vendredi \\
\hline
Autres observations & Moustiquaires imprégnées utilisées par seulement 30 \% des enfants ; nombreux enfants se couchent tard (veillées, télévision) \\
\hline
\end{tabularx}
\end{center}

Face à cette situation, la chefferie décide de créer un \textbf{conseil municipal de santé des jeunes} et sollicite ta classe de Première A. Par groupes, vous êtes invités à jouer le rôle de conseillers municipaux de santé : votre mission est de comprendre \emph{pourquoi} les enfants du quartier tombent malades, de convaincre les parents de faire vacciner leurs enfants, et de proposer à la chefferie un \textbf{plan d'action} réaliste et justifié scientifiquement.

\courssub{Consignes}

Travail en groupes de 4 à 6 élèves. Chaque groupe produit une affiche et prépare une restitution orale de 3 minutes.

\begin{enumerate}
\item \textbf{Expliquer.} À l'aide d'un schéma annoté, explique pourquoi un enfant \emph{non vacciné} exposé au virus de la rougeole tombe malade, alors qu'un enfant \emph{vacciné} est protégé. Ton schéma devra faire apparaître : l'antigène, les lymphocytes, les anticorps, la réaction immunitaire primaire et la mémoire immunitaire.
\item \textbf{Justifier.} Rédige un argumentaire de 5 lignes maximum, destiné à être lu par un élève lors d'une sensibilisation des parents au marché du quartier, pour convaincre de l'intérêt d'une campagne de vaccination contre la rougeole. Chaque affirmation doit reposer sur un mécanisme immunitaire ou sur une donnée de la fiche.
\item \textbf{Proposer.} Élabore un plan d'action en \textbf{5 mesures classées par ordre de priorité}, couvrant les domaines suivants : vaccination, alimentation, hygiène et eau, sommeil/gestion du stress, prévention du paludisme. Pour chaque mesure, indique : (a) l'action concrète proposée à la chefferie ; (b) le mécanisme immunitaire ou physiologique qui la justifie.
\item \textbf{Restituer.} Présente oralement ton plan en 3 minutes : un membre explique le schéma, un autre lit l'argumentaire, un troisième défend les deux mesures prioritaires. La classe et l'enseignant établissent ensuite une synthèse collective au tableau.
\end{enumerate}

\courssub{Ressources à mobiliser}

Pour réussir cette activité, tu dois réactiver les savoirs essentiels de la leçon :

\begin{definition}[Organisation du système immunitaire]
Les \textbf{organes lymphoïdes} assurent la maturation et le stockage des cellules immunitaires : \textbf{moelle osseuse} (maturation des lymphocytes B, production de toutes les cellules sanguines), \textbf{thymus} (maturation des lymphocytes T), \textbf{rate} (filtration du sang, site de réaction immunitaire), \textbf{ganglions lymphatiques} (filtration de la lymphe, site de réaction immunitaire). Les \textbf{cellules} effectrices principales sont : les \textbf{lymphocytes B} (production d'anticorps), les \textbf{lymphocytes T} (cytotoxiques ou auxiliaires), les \textbf{phagocytes} (macrophages, neutrophiles : ingestion et destruction des agents pathogènes).
\end{definition}

\begin{definition}[Antigène et Anticorps]
Un \textbf{antigène} (Ag) est une substance étrangère à l'organisme (protéine de surface d'un virus, d'une bactérie, toxine) capable d'être reconnue spécifiquement par des récepteurs de lymphocytes ou par des anticorps. Un \textbf{anticorps} (Ac) ou immunoglobuline est une protéine en Y, sécrétée par les \textbf{plasmocytes} (lymphocytes B différenciés), qui se fixe spécifiquement sur un antigène pour le neutraliser, l'opsoniser ou activer le complément.
\end{definition}

\begin{propriete}[Principes de la réaction immunitaire et de la vaccination]
\begin{itemize}
\item \textbf{Réaction immunitaire primaire} (premier contact avec l'Ag) : reconnaissance $\rightarrow$ activation et multiplication clonale des lymphocytes spécifiques $\rightarrow$ différenciation en plasmocytes (production d'Ac) et en \textbf{lymphocytes mémoires} $\rightarrow$ élimination de l'agent. Réponse \emph{lente} (jours à semaines), taux d'Ac faible et transitoire.
\item \textbf{Réaction immunitaire secondaire} (contact ultérieur avec le même Ag) : activation rapide des lymphocytes mémoires $\rightarrow$ production massive et rapide d'Ac à haute affinité. Réponse \emph{rapide} (heures à jours), \emph{intense} et \emph{durable}. Pas de symptômes cliniques.
\item \textbf{Vaccination} : introduction d'un antigène inoffensif (agent tué, atténué, sous-unité protéique, ARN messager) qui mime l'infection et déclenche une réaction primaire \textbf{sans maladie}, dotant l'organisme de lymphocytes mémoires protecteurs.
\end{itemize}
\end{propriete}

\begin{aretenir}[Types d'immunité et facteurs de risque]
\begin{itemize}
\item \textbf{Immunité innée} : non spécifique, présente dès la naissance (barrières, phagocytes, inflammation).
\item \textbf{Immunité acquise} : spécifique, mémoire. \emph{Naturelle} (après maladie) ou \emph{artificielle} (vaccination active = injection d'Ag ; vaccination passive = injection d'Ac, ex. sérum antitetanique).
\item \textbf{Facteurs affaiblissant l'immunité} : \textbf{malnutrition} (carence protéino-énergétique, vitamines A, D, Zn, Fe $\rightarrow$ synthèse d'Ac et prolifération lymphocytaire altérées), \textbf{stress chronique} et \textbf{manque de sommeil} (cortisol $\uparrow$ $\rightarrow$ immunosuppression), \textbf{maladies chroniques/répétées} (paludisme, VIH/Sida, tuberculose $\rightarrow$ épuisement/anémie/destruction lymphocytaire).
\end{itemize}
\end{aretenir}

\courssub{Démarche et corrigé commenté}

\begin{exoresolu}[Corrigé commenté de l'activité]
\textbf{Étape 1 : Expliquer pourquoi les enfants non vaccinés tombent malades (Schéma annoté).}

Chez un enfant \emph{non vacciné}, le virus de la rougeole pénètre dans l'organisme (voies respiratoires) pour la première fois. Ses protéines de surface (hémagglutinine, fusion) jouent le rôle d'\cle{antigènes}. Les cellules présentatrices d'antigènes (macrophages, cellules dendritiques) présentent l'antigène aux lymphocytes T auxiliaires (CD4+), qui activent les lymphocytes B spécifiques. Ceux-ci déclenchent une \textbf{réaction immunitaire primaire} : multiplication clonale, différenciation en \textbf{plasmocytes} (usines à anticorps) et en \textbf{lymphocytes B mémoires}. La production d'anticorps spécifiques (IgM puis IgG) est \emph{lente} (pic vers 7--10 jours) : pendant ce délai, le virus se réplique librement et provoque les symptômes (fièvre, exanthème, complications). L'enfant tombe donc malade avant d'être protégé. La guérison laisse une population de lymphocytes mémoires.

Chez un enfant \emph{vacciné} (vaccin VAR : virus vivant atténué), l'organisme a déjà rencontré l'antigène vers 9 mois. Des \textbf{lymphocytes B et T mémoires} spécifiques circulent. Lors du contact avec le virus sauvage, la \textbf{réaction immunitaire secondaire} est déclenchée : les lymphocytes mémoires se transforment rapidement en plasmocytes et produisent massivement des \textbf{IgG à haute affinité} en \emph{quelques heures/jours}. Le virus est neutralisé avant d'atteindre un seuil pathogène. L'enfant ne tombe pas malade.

\begin{popfigure}
\begin{tikzpicture}[
    node distance=1.2cm and 1.5cm,
    box/.style={rectangle, draw=popdark, thick, rounded corners, fill=popcream, minimum width=2.5cm, minimum height=1cm, align=center, font=\small},
    arrow/.style={-Stealth, thick, popdark},
    title/.style={font=\bfseries\small, text=popblue, anchor=south},
    phase/.style={font=\bfseries\small\itshape, text=poppurple, anchor=north}
]
% --- PANEL A: Non vacciné ---
\node[title] (tA) at (0, 4.5) {Enfant NON VACCINÉ (1er contact)};
\node[phase] (pA1) at (-4, 3.8) {Phase d'incubation / Réponse primaire};
\node[box, align=center] (vA) at (-4, 3){Virus rougeole\\ (Antigènes Ag)};
\node[box, align=center] (apcA) at (-4, 1.5){Cellule présentatrice\\ d'antigène (APC)};
\node[box, align=center] (lbA) at (0, 1.5){Lymphocyte B naïf\\ (récepteur anti-Ag)};
\node[box, align=center] (plA) at (4, 2.2){Plasmocyte\\ (Production Ac)};
\node[box, align=center] (memA) at (4, 0.5){Lymphocyte B\\ Mémoire};
\node[box, align=center] (acA) at (4, -0.8){Anticorps (IgM $\rightarrow$ IgG)\\ \textit{Apparition lente (J7--J10)}};
\node[box, fill=poppinkL, draw=poppink, align=center] (symA) at (-4, -0.8){SYMPTÔMES : Fièvre,\\ éruption, complications};
\node[box, fill=popgreenL, draw=popgreen, align=center] (protA) at (0, -2.2){Guérison +\\ Mémoire immunitaire};

\draw[arrow] (vA) -- (apcA);
\draw[arrow] (apcA) -- (lbA);
\draw[arrow] (lbA) -- (plA);
\draw[arrow] (lbA) -- (memA);
\draw[arrow] (plA) -- (acA);
\draw[arrow, poppink] (vA) -- (symA);
\draw[arrow, popgreen] (acA) -- (protA);
\draw[arrow, popgreen] (memA) -- (protA);

% --- PANEL B: Vacciné ---
\node[title] (tB) at (10, 4.5) {Enfant VACCINÉ (Contact ultérieur)};
\node[phase] (pB1) at (6, 3.8) {Vaccination (J-9 mois)};
\node[box, align=center] (vacB) at (6, 3){Vaccin VAR\\ (Virus atténué, Ag)};
\node[box, align=center] (memB1) at (6, 1.5){Lymphocytes B \& T\\ Mémoires formés};

\node[phase] (pB2) at (14, 3.8) {Infection naturelle (Virus sauvage)};
\node[box, align=center] (vB) at (14, 3){Virus sauvage\\ (Mêmes Ag)};
\node[box, align=center] (memB2) at (14, 1.5){Lymphocytes Mémoires\\ (Reconnaissance immédiate)};
\node[box, align=center] (plB) at (18, 1.5){Plasmocytes\\ (Différenciation rapide)};
\node[box, align=center] (acB) at (18, 0){Anticorps IgG\\ \textit{Massifs, précoces (J1--J3), haute affinité}};
\node[box, fill=popgreenL, draw=popgreen, align=center] (protB) at (14, -1){PROTECTION :\\ Virus neutralisé,\\ Pas de maladie};

\draw[arrow] (vacB) -- (memB1);
\draw[arrow] (vB) -- (memB2);
\draw[arrow] (memB1) -- node[above, sloped, font=\tiny, pos=0.5] {Persistance} (memB2);
\draw[arrow] (memB2) -- (plB);
\draw[arrow] (plB) -- (acB);
\draw[arrow, popgreen] (acB) -- (protB);
\draw[arrow, popgreen] (memB2) -- (protB);

% Braces for timing
\draw[decorate, decoration={brace, amplitude=5pt}, thick, poporange] (-4.5, -1.5) -- (4.5, -1.5) node[midway, below=5pt, font=\small\bfseries, text=poporange] {Délai de plusieurs jours : Fenêtre de vulnérabilité};
\draw[decorate, decoration={brace, amplitude=5pt}, thick, popgreen] (10, -1.5) -- (18, -1.5) node[midway, below=5pt, font=\small\bfseries, text=popgreen] {Réponse en heures : Protection efficace};
\end{tikzpicture}
\legende{Schéma comparatif de la réponse immunitaire chez un enfant non vacciné (panneau gauche) et vacciné (panneau droit). La vaccination préétablit la mémoire immunitaire, permettant une réponse secondaire rapide (IgG massives) qui neutralise le virus avant l'apparition des symptômes.}
\end{popfigure}

\emph{Commentaire pédagogique :} cette explication mobilise directement le savoir-faire « expliquer le principe général d'une réaction immunitaire » et la notion de mémoire immunitaire, fondement de la vaccination. Le schéma attendu oppose la cinétique de la réponse primaire (lente, symptômes) et secondaire (rapide, protection).

\tcblower

\textbf{Étape 2 : Argumentaire de sensibilisation (5 lignes maximum).}

Exemple de texte attendu : « Parents de Tougang, la rougeole a déjà frappé 22 de nos enfants le mois dernier, et 15 d'entre eux n'étaient pas vaccinés. Le vaccin apprend au corps de l'enfant à reconnaître le virus sans le rendre malade : il fabrique des anticorps et garde en mémoire l'ennemi. Ainsi protégé, l'enfant détruit le virus dès son arrivée et ne tombe pas malade. Aujourd'hui, seuls 55 \% de nos enfants sont vaccinés, alors qu'il en faut 95 \% pour protéger tout le quartier (immunité collective). Amenez vos enfants au CSI de la localité voisine (ouvert mardi et vendredi, à 3 km) : le vaccin VAR est gratuit, sûr, et il sauve des vies. »

\emph{Commentaire :} l'argumentaire combine une donnée chiffrée locale (55 \% contre 95 \%), le mécanisme (anticorps + mémoire), la notion d'immunité collective et un appel à l'action réaliste (localisation, jours, gratuité).

\textbf{Étape 3 : Plan d'action en 5 mesures classées par priorité.}

\begin{enumerate}
\item \textbf{Vaccination (Priorité 1 -- Urgence épidémiologique).} \textbf{Action :} Organiser avec l'équipe du CSI une campagne de rattrapage vaccinal « porte-à-porte » et au marché contre la rougeole (VAR) et les autres antigènes du PEV (Polio, Penta, ROR), ciblant les 45 \% d'enfants non vaccinés. \textbf{Justification :} La vaccination confère une \textbf{immunité artificielle active} avec mémoire immunitaire. Atteindre 95 \% de couverture vaccinale interrompt la circulation du virus (seuil d'immunité collective pour la rougeole, $R_0 \approx 12-18$), protégeant aussi les nourrissons trop jeunes pour être vaccinés et les immunodéprimés.
\item \textbf{Hygiène de l'eau et des mains (Priorité 2 -- Cause majeure des diarrhées).} \textbf{Action :} Sensibiliser au traitement de l'eau du forage (ébullition 1 min ou 2 gouttes d'eau de Javel à 2,6 \% par litre, 30 min d'attente) et au lavage des mains au savon (ou cendres) aux moments critiques (après latrines, avant repas, après change). \textbf{Justification :} Réduire l'ingestion de pathogènes fécaux (bactéries, virus, parasites) diminue la \textbf{charge antigénique} intestinale. Cela évite l'activation chronique du système immunitaire muqueux (GALT), la perte de nutriments par diarrhée et l'affaiblissement nutritionnel qui en découle.
\item \textbf{Alimentation enrichie (Priorité 3 -- Soutien structurel de l'immunité).} \textbf{Action :} Promouvoir l'enrichissement quotidien du « couscous » ou « fufu » maïs/manioc avec des sources protéiques locales accessibles : haricots/niébé, arachides (pâte), poisson fumé séché, soja, œufs ; et la consommation de fruits de saison (papaye, mangue, agrumes pour vit. C). \textbf{Justification :} Les \textbf{anticorps sont des glycoprotéines} ; la synthèse des lymphocytes et des médiateurs de l'inflammation requiert des acides aminés essentiels, du zinc, du fer, des vitamines A, D, C, E. La malnutrition protéino-énergétique (MPE) entraîne une atrophie du thymus et des ganglions, une baisse des IgA sécrétoires et une anergie lymphocytaire.
\item \textbf{Prévention du paludisme (Priorité 4 -- Facteur d'épuisement immunitaire).} \textbf{Action :} Campagne de distribution et d'installation correcte de moustiquaires imprégnées à longue durée d'action (MILDA) pour atteindre 100 \% des lits d'enfants ; traitement préventif intermittent (TPI) si zone à transmission saisonnière. \textbf{Justification :} Le paludisme à répétition (\emph{Plasmodium falciparum}) induit une \textbf{activation polyclonale des lymphocytes B} (hypergammaglobulinémie non spécifique) qui « noie » la réponse spécifique, une \textbf{anémie} (destruction hématies, dépression médullaire) et une \textbf{immunosuppression} transitoire favorisant les surinfections bactériennes et virales (rougeole, diarrhées).
\item \textbf{Sommeil et gestion du stress (Priorité 5 -- Régulation neuro-immunitaire).} \textbf{Action :} Sensibilisation des parents sur l'importance d'un coucher régulier (20 h--21 h) pour 10--11 h de sommeil ; limiter les écrans le soir ; créer un cadre apaisant. \textbf{Justification :} Le \textbf{manque de sommeil} et le \textbf{stress chronique} élèvent le \textbf{cortisol} (glucocorticoïde), qui inhibe la prolifération des lymphocytes T, la production d'IL-2 et d'IFN-$\gamma$, et la fonction phagocytaire. Le sommeil favorise au contraire la sécrétion de cytokines pro-inflammatoires (IL-12, TNF-$\alpha$) nécessaires à la réponse adaptative et la formation de la mémoire immunitaire.
\end{enumerate}

\textbf{Étape 4 : Restitution et synthèse collective.} Chaque groupe présente en 3 minutes (schéma, argumentaire, 2 mesures prioritaires). La classe confronte les classements de priorités (débat riche : certains placent l'eau avant la vaccination car cause immédiate des diarrhées ; d'autres la vaccination car urgence épidémique rougeole). L'enseignant guide la synthèse au tableau en reliant chaque mesure à un levier immunitaire précis (mémoire, charge antigénique, substrat protéique, épuisement, régulation hormonale).

\end{exoresolu}

\courssub{Bilan de l'activité}

Cette activité t'a placé dans la peau d'un véritable acteur de santé communautaire. Elle a permis de mettre en évidence plusieurs idées essentielles :

\begin{itemize}
\item La \textbf{mémoire immunitaire} (lymphocytes B et T mémoires) est la clé de la protection : c'est elle qui explique la différence de destin entre un enfant vacciné et un enfant non vacciné confrontés au même virus sauvage.
\item La vaccination n'est pas seulement un acte individuel : au-delà d'un certain taux de couverture (environ 95 \% pour la rougeole, calculé via $1 - 1/R_0$), elle protège l'ensemble de la communauté, y compris les enfants trop jeunes ou trop fragiles pour être vaccinés (\textbf{immunité collective} ou de groupe).
\item Le système immunitaire ne fonctionne pas isolément : son efficacité dépend de l'\textbf{alimentation} (les anticorps et cellules sont faits de protéines, nécessitent micronutriments), de l'\textbf{hygiène} (moins d'agressions antigéniques chroniques à combattre), du \textbf{sommeil} et de l'absence de \textbf{maladies répétées} comme le paludisme (épuisement, anémie, immunosuppression).
\item Les savoirs scientifiques de la leçon ont une utilité sociale directe : ils permettent d'analyser un problème de santé publique réel, de convaincre par l'argument scientifique et de proposer des solutions concrètes, adaptées aux ressources locales (forage, CSI, cultures vivrières, relais communautaires).
\end{itemize}

Tu as ainsi mobilisé les trois savoir-faire officiels de la leçon — \cle{expliquer} une réaction immunitaire, \cle{justifier} la vaccination, \cle{proposer} des comportements favorables à l'immunité — dans une situation authentique, proche de celles que rencontrent chaque jour les comités de santé de nos quartiers et villages camerounais.

\newpage

\coursec{Méthodes et exercices résolus}

\courssub{Méthodes et exercices résolus}

\begin{methode}[Expliquer le principe général d'une réaction immunitaire]
    \textbf{Objectif :} Décrire de façon séquentielle et rigoureuse la mise en place de la réponse immunitaire adaptative (spécifique) à partir de la rencontre avec un antigène.
    \begin{enumerate}
        \item \textbf{Identifier l'agent pathogène et ses antigènes :} Préciser la nature de l'antigène (protéine de surface bactérienne, capside virale, toxine) et son caractère \cle{non-self}.
        \item \textbf{Décrire la phase de reconnaissance et de présentation :} 
        \begin{itemize}
            \item Phagocytose par un macrophage (ou cellule dendritique) $\rightarrow$ digestion lysosomale.
            \item Présentation des fragments antigéniques (péptides) à la surface de la cellule présentatrice d'antigène (CPA) via les molécules du \cle{CMHC} (Complexe Majeur d'Histocompatibilité) de classe II.
        \end{itemize}
        \item \textbf{Décrire l'activation des lymphocytes T auxiliaires (LT4/CD4+) :} Reconnaissance spécifique du complexe \og antigène-CMHC II \fg par le récepteur (TCR) d'un LT4 naïf $\rightarrow$ prolifération clonale $\rightarrow$ différenciation en LT4 effecteurs (sécrétion d'interleukines, notamment IL-2) et LT4 mémoire.
        \item \textbf{Décrire l'activation des lymphocytes B (LB) :} 
        \begin{itemize}
            \item Reconnaissance directe de l'antigène natif par le récepteur de membrane (BCR = IgM/IgD) d'un LB naïf spécifique.
            \item Co-stimulation indispensable par les interleukines (IL-2, IL-4, IL-5, IL-6) sécrétées par le LT4 effecteur (aide lymphocytaire T).
        \end{itemize}
        \item \textbf{Décrire la différenciation des LB :} Prolifération clonale $\rightarrow$ différenciation en \cle{plasmocytes} (usines à anticorps) et \cle{lymphocytes B mémoire}.
        \item \textbf{Préciser l'action effectrice des anticorps :} Neutralisation, opsonisation, activation du complément, ADCC (cytotoxicité à médiation cellulaire dépendante des anticorps).
        \item \textbf{Conclure sur la mémoire immunitaire :} Persistance des lymphocytes mémoire (T et B) assurant une réponse secondaire plus rapide, plus intense et plus durable.
    \end{enumerate}
    \textbf{Reflexe Probatoire :} Utiliser les termes exacts : \og présentation antigénique \fg, \og CMHC II \fg, \og prolifération clonale \fg, \og plasmocyte \fg, \og mémoire immunologique \fg. Ne pas confondre antigène (sur l'agent) et anticorps (produit par l'organisme).
\end{methode}

\begin{exoresolu}[Analyse des phases de la réponse immunitaire adaptative]
    \textbf{Énoncé :} Lors d'une primo-infection par \emph{Streptococcus pneumoniae} (pneumocoque) chez un élève de 16 ans, on observe l'apparition d'anticorps spécifiques dans le sérum au 7\textsuperscript{e} jour, avec un pic au 14\textsuperscript{e} jour. Au cours d'une ré-infection par la même souche un an plus tard, le pic d'anticorps est atteint dès le 3\textsuperscript{e} jour et son taux est 100 fois plus élevé.
    \begin{enumerate}
        \item En t'appuyant sur tes connaissances, explique les mécanismes cellulaires et moléculaires qui se déroulent entre le jour 0 et le jour 7 de la primo-infection pour aboutir à la production d'anticorps.
        \item Interprète les différences observées entre la primo-infection et la ré-infection.
        \item Pourquoi l'administration d'un sérum antivenimeux (anticorps polyvalents) après une morsure de serpent ne confère-t-elle pas de protection durable ?
    \end{enumerate}
    \tcblower
    \textbf{Solution détaillée :}
    \begin{enumerate}
        \item \textbf{Mécanismes de la primo-réponse (J0 à J7) :}
        \begin{itemize}
            \item \textbf{Reconnaissance innée et présentation :} Les pneumocoques sont phagocytés par des macrophages (et cellules dendritiques) au niveau du tissu conjonctif ou des ganglions lymphatiques drainants. Les antigènes bactériens (capsule polysaccharidique, protéines de surface) sont dégradés en peptides dans les phagolysosomes.
            \item \textbf{Présentation aux LT4 :} Les peptides sont présentés à la surface des CPA via les molécules du \textbf{CMHC de classe II}. Un lymphocyte T auxiliaire naïf (LT4/CD4+) dont le récepteur (TCR) est spécifique de ce complexe \og peptide-CMHC II \fg le reconnaît. Cette reconnaissance, renforcée par des molécules d'adhésion (CD4, CD28/B7), active le LT4.
            \item \textbf{Activation et prolifération clonale des LT4 :} Le LT4 activé sécrète de l'\textbf{IL-2} (autocrine) et d'autres interleukines (IL-4, IL-5, IL-6). Il subit une \textbf{prolifération clonale} massive et se différencie en \textbf{LT4 effecteurs} (sécréteurs de lymphokines) et en \textbf{LT4 mémoire}.
            \item \textbf{Activation des LB :} Parallèlement, un lymphocyte B naïf spécifique reconnaît l'antigène \emph{natif} (conformationnel) via son BCR (IgM/IgD de membrane). Il internalise l'antigène, le traite et présente des peptides sur son CMHC II.
            \item \textbf{Coopération T-B (Aide lymphocytaire) :} Le LT4 effecteur spécifique reconnaît le complexe \og peptide-CMHC II \fg sur le LB. Par contact cellulaire (CD40L/CD40) et sécrétion d'interleukines (IL-4, IL-5, IL-6), il délivre le \textbf{second signal} indispensable à l'activation complète du LB.
            \item \textbf{Différenciation en plasmocytes :} Le LB activé prolifère (prolifération clonale) et se différencie principalement en \textbf{plasmocytes}. Ces cellules, riches en réticulum endoplasmique rugueux et en appareil de Golgi, synthétisent et sécrètent massivement des \textbf{anticorps spécifiques (IgM d'abord, puis IgG après commutation de classe)} dans le sang et la lymphe. C'est l'arrivée de ces anticorps dans le sérum qui est détectée vers J7.
        \end{itemize}
        \item \textbf{Interprétation des différences (Mémoire immunologique) :}
        \begin{itemize}
            \item \textbf{Délai raccourci (J3 vs J7) :} Lors de la ré-infection, les \textbf{lymphocytes B mémoire} et \textbf{LT4 mémoire} (générés lors de la primo-infection et persistant depuis un an) sont déjà présents en grand nombre et pré-activés. Ils ne nécessitent pas la phase lente de présentation/activation/coopération naïve. Ils se différencient immédiatement en plasmocytes.
            \item \textbf{Amplitude accrue (x100) :} Le nombre de cellules mémoire spécifiques est bien supérieur au nombre de cellules naïves initiales. De plus, les anticorps produits sont majoritairement des \textbf{IgG} (affinité plus forte grâce à la maturation de l'affinité survenue lors de la primo-réponse) plus efficaces que les IgM initiaux.
        \end{itemize}
        \item \textbf{Nature de la protection par sérum antivenimeux :} C'est une \textbf{immunité passive artificielle}. Les anticorps (IgG purifiées de cheval ou mouton) sont administrés directement. Ils neutralisent le venin immédiatement mais :
        \begin{itemize}
            \item Ils sont \textbf{exogènes} : l'organisme du patient ne les a pas produits, il n'y a pas eu activation de ses propres lymphocytes B ni prolifération clonale.
            \item Ils ont une \textbf{durée de vie limitée} (demi-vie des IgG $\approx$ 21-28 jours) et sont éliminés par le catabolisme normal.
            \item \textbf{Aucun lymphocyte mémoire} n'est généré chez le receveur.
        \end{itemize}
        Donc, la protection disparaît en quelques semaines sans laisser de mémoire.
    \end{enumerate}
    \textbf{Conclusion :} La réponse immune adaptative repose sur la sélection clonale de lymphocytes spécifiques, leur amplification et leur différenciation en cellules effectrices (plasmocytes) et cellules mémoire, expliquant la spécificité, le délai de la primo-réponse et l'efficacité de la réponse secondaire.
\end{exoresolu}

\begin{methode}[Identifier et situer les organes et cellules du système immunitaire]
    \textbf{Objectif :} Relier la structure (localisation, histologie) à la fonction (site de maturation, site de rencontre antigène/lymphocyte, site de stockage).
    \begin{enumerate}
        \item \textbf{Classer les organes lymphoïdes :}
        \begin{itemize}
            \item \textbf{Primaires (organes centraux) :} Site de genèse et maturation des lymphocytes à partir de cellules souches hématopoïétiques.
                \begin{itemize}
                    \item \textbf{Moelle osseuse rouge (côtes, sternum, bassin, vertèbres) :} Maturation des \textbf{Lymphocytes B} (chez l'homme et les mammifères) et des cellules NK. Hématopoïèse générale.
                    \item \textbf{Thymus (médiastin antérieur, volumineux chez l'enfant, involution graisseuse à l'âge adulte) :} Maturation des \textbf{Lymphocytes T} (sélection positive/négative $\rightarrow$ tolérance au self).
                \end{itemize}
            \item \textbf{Secondaires (organes périphériques) :} Sites de rencontre entre antigènes et lymphocytes naïfs, déclenchement de la réponse immune, stockage des cellules mémoire.
                \begin{itemize}
                    \item \textbf{Ganglions lymphatiques (cervicaux, axillaires, inguinaux, mésentériques) :} Filtrent la \textbf{lymphe}. Structure : corticale (follicules lymphoïdes = zone B, paracortex = zone T) et médullaire (cordons lymphoïdes = plasmocytes). Entrée : vaisseaux lymphatiques afférents ; Sortie : vaisseau efférent.
                    \item \textbf{Rate (fosse iliaque gauche) :} Filtre le \textbf{sang}. Pulpe blanche (follicules B, gaine lymphoïde périaartérielle = zone T) et pulpe rouge (stockage érythrocytes, élimination globules rouges vieillis). Rôle majeur contre bactéries encapsulées (pneumocoque, méningocoque, \emph{Haemophilus}).
                    \item \textbf{Tissu lymphoïde associé aux muqueuses (MALT) :} Amygdales (Waldeyer), plaques de Peyer (iléon), appendice, voies respiratoires (BALT), urogénitales. Première ligne face aux antigènes muqueux. Production d'IgA sécrétoires (dimériques + pièce sécrétoire).
                \end{itemize}
        \end{itemize}
        \item \textbf{Identifier les cellules effectrices clés :}
        \begin{itemize}
            \item \textbf{Phagocytes professionnels :} Macrophages (tissus), Neutrophiles (sang, premier recours), Cellules dendritiques (pont inné/adapté, meilleures CPA).
            \item \textbf{Lymphocytes :} LT4 (CD4+, chef d'orchestre), LT8 (CD8+, cytotoxiques tueurs de cellules infectées), LB (producteurs d'anticorps), NK (immunité innée, tueurs de cellules sans CMHC I).
            \item \textbf{Autres :} Mastocytes/basophiles (allergie, inflammation), Éosinophiles (parasites, allergie).
        \end{itemize}
    \end{enumerate}
    \textbf{Reflexe Probatoire :} Ne pas confondre \og maturation \fg (primaires) et \og activation \fg (secondaires). La rate filtre le \textbf{sang}, les ganglions filtrent la \textbf{lymphe}. L'absence de rate (splénectomie) expose aux infections à germes encapsulés $\rightarrow$ vaccination antipneumococcique/anti-Hib/antiméningococcique obligatoire avant chirurgie.
\end{methode}

\begin{exoresolu}[Localisation et rôle des organes lymphoïdes]
    \textbf{Énoncé :} Un enfant de 8 ans vivant à Douala présente une splénomégalie (rate augmentée de volume) et une adénopathie cervicale (ganglions cervicaux augmentés) lors d'un paludisme aigu à \emph{Plasmodium falciparum}. Par ailleurs, un jeune adulte de 22 ans victime d'un accident de la route subit une splénectomie totale (ablation de la rate) en urgence.
    \begin{enumerate}
        \item Explique les mécanismes physiologiques à l'origine de la splénomégalie et de l'adénopathie cervicale lors de l'accès palustre.
        \item Quels sont les risques infectieux majeurs pour le jeune adulte splénectomisé ? Justifie par le rôle anatomique et fonctionnel de la rate.
        \item Cite deux mesures prophylactiques indispensables à mettre en place pour ce patient après l'intervention.
    \end{enumerate}
    \tcblower
    \textbf{Solution détaillée :}
    \begin{enumerate}
        \item \textbf{Origine de l'hypertrophie ganglionnaire et splénique :}
        \begin{itemize}
            \item \emph{Plasmodium falciparum} infecte les globules rouges (GR). L'hémolyse libère des antigènes parasitiques (hémozoïne, protéines de membrane) dans la circulation sanguine et la lymphe.
            \item \textbf{Rate (filtre sanguin) :} Le sang passe dans les sinusoïdes de la pulpe rouge. Les macrophages de la pulpe rouge phagocytent les GR infectés et les débris. Les antigènes gagnent la pulpe blanche (zone T et follicules B). Cela déclenche une \textbf{réaction immunitaire intense} : prolifération clonale des lymphocytes B et T spécifiques $\rightarrow$ hyperplasie des follicules lymphoïdes (centres germinatifs) et de la gaine lymphoïde périaartérielle $\rightarrow$ \textbf{splénomégalie reactive}.
            \item \textbf{Ganglions cervicaux (filtre lymphatique) :} La lymphe drainant les tissus de la tête et du cou (où peuvent se trouver des parasites ou antigènes drainés) arrive aux ganglions cervicaux. La présentation antigénique et l'activation lymphocytaire y provoquent une \textbf{hyperplasie lymphoïde} (prolifération dans le cortex et le paracortex) $\rightarrow$ \textbf{adénopathie} (ganglion palpable, mou, mobile, douloureux = réactionnel).
        \end{itemize}
        \item \textbf{Risques infectieux post-splénectomie (Syndrome de surinfection post-splénectomie - SSPS) :}
        \begin{itemize}
            \item \textbf{Risque majeur :} Infections invasives à \textbf{germes encapsulés} : \emph{Streptococcus pneumoniae} (pneumocoque) $>$ 50\% des cas, \emph{Neisseria meningitidis} (méningocoque), \emph{Haemophilus influenzae} type b.
            \item \textbf{Justification anatomofonctionnelle :}
                \begin{itemize}
                    \item \textbf{Perte du filtre sanguin mécanique :} La pulpe rouge retient physiquement les particules (GR infectés, bactéries) grâce à la circulation lente dans les sinusoïdes et aux fentes inter-endothéliales. Sans rate, les bactéries circulent librement.
                    \item \textbf{Perte de la fonction opsonique :} La rate est le principal site de production d'opsonines (properdine, tuftsin) et de fixation du complément sur les bactéries encapsulées. Les macrophages spléniques sont très efficaces pour phagocyter les bactéries opsonisées.
                    \item \textbf{Perte de la réponse immune T-dépendante aux antigènes polysaccharidiques :} Les antigènes capsulaires (polysaccharides) sont des antigènes \textbf{T-indépendants} (faible réponse, pas de mémoire, pas de commutation IgM$\rightarrow$IgG) chez le sujet normal. La rate, par sa structure unique (contact sang-lent, forte concentration de cellules marginales zone B1), permet une réponse T-dépendante aux polysaccharides (production d'IgG à haute affinité, mémoire). Sans rate, le patient ne fait que des IgM peu protectrices contre ces germes.
                \end{itemize}
        \end{itemize}
        \item \textbf{Mesures prophylactiques indispensables :}
        \begin{enumerate}
            \item \textbf{Vaccination systématique (avant l'intervention si programmée, sinon dès que possible après) :}
                \begin{itemize}
                    \item Vaccin anti-pneumococcique conjugué (PCV13/PCV15/PCV20) + polysaccharidique (PPSV23) selon schéma séquentiel.
                    \item Vaccin anti-\emph{Haemophilus influenzae} type b (Hib).
                    \item Vaccin anti-méningococcique tétravalent conjugué (ACYW135) + sérogroupe B.
                    \item \emph{Note :} Au Cameroun, le PCV et Hib font partie du PEV (Programme Élargi de Vaccination), le méningocoque est recommandé en zone de méningite (Nord).
                \end{itemize}
            \item \textbf{Antibioprophylaxie au long cours :} \textbf{Pénicilline V per os} (ou Amoxicilline) \textbf{à vie} (ou au moins jusqu'à 18-21 ans, voire 5 ans post-op selon protocoles), surtout chez l'enfant. Alternative si allergie : Érythromycine ou Azithromycine.
            \item (Complément) \textbf{Éducation thérapeutique :} Fièvre $\geq$ 38,5°C = urgence médicale, consultation immédiate pour hémocultures et antibiotiques IV probabilistes (Céphalosporine 3\textsuperscript{e} gén +/- Aminoside) sans attendre les résultats. Carte d'alerte "Patient splénectomisé" à porter en permanence.
        \end{enumerate}
    \end{enumerate}
    \textbf{Conclusion :} La rate est l'organe lymphoïde secondaire spécialisé dans la filtration sanguine et la réponse aux antigènes polysaccharidiques ; son ablation expose à un risque vital de septicémie foudroyante à germes encapsulés, prévenue par la vaccination et l'antibioprophylaxie.
\end{exoresolu}

\begin{methode}[Différencier les types d'immunité (Innée / Acquise ; Active / Passive ; Naturelle / Artificielle)]
    \textbf{Objectif :} Classer une situation immunologique donnée dans le bon croisement de critères (origine + mode d'acquisition).
    \begin{enumerate}
        \item \textbf{Critère 1 : Spécificité et Mémoire (Innée vs Acquise/Adaptative)}
        \begin{itemize}
            \item \textbf{Innée (Non-spécifique) :} Présente à la naissance, pas de mémoire, réponse immédiate (min/heures). Barrières (peau, muqueuses, cils, pH, flore), cellules (macrophages, neutrophiles, NK, mastocytes), molécules (complément, interférons, lysozyme, protéines de phase aiguë : CRP). Reconnaît des \cle{MAP} (Motifs Associés aux Pathogènes) via \cle{PRR} (Récepteurs de Reconnaissance de Formes, ex: TLR).
            \item \textbf{Acquise (Adaptative/Spécifique) :} Se met en place au cours de la vie, \textbf{spécifique} (un antigène $\rightarrow$ un récepteur), \textbf{mémoire}. Lymphocytes T et B, Anticorps. Délai de mise en route (jours).
        \end{itemize}
        \item \textbf{Critère 2 : Origine des anticorps/lymphocytes (Active vs Passive)}
        \begin{itemize}
            \item \textbf{Active :} L'organisme \textbf{produit} lui-même ses anticorps et ses lymphocytes mémoire. $\rightarrow$ Protection \textbf{durable} (années, vie). Délai de mise en place (semaines).
            \item \textbf{Passive :} L'organisme \textbf{reçoit} des anticorps tout faits (ou des lymphocytes). $\rightarrow$ Protection \textbf{immédiate} mais \textbf{transitoire} (semaines/mois, durée de vie des IgG $\approx$ 3 semaines). \textbf{Pas de mémoire}.
        \end{itemize}
        \item \textbf{Critère 3 : Intervention humaine (Naturelle vs Artificielle)}
        \begin{itemize}
            \item \textbf{Naturelle :} Processus physiologique spontané (infection, grossesse, allaitement).
            \item \textbf{Artificielle :} Intervention médicale volontaire (vaccination, sérumthérapie).
        \end{itemize}
        \item \textbf{Tableau de synthèse (à connaître par cœur) :}
        \begin{center}
        \begin{tabular}{|c|c|c|c|}\hline
        \textbf{Type} & \textbf{Exemple} & \textbf{Mécanisme} & \textbf{Durée / Mémoire} \\ \hline
        \textbf{Active Naturelle} & Guérison d'une rougeole, tuberculose & Infection $\rightarrow$ Réponse immune complète & Longue (souvent vie) / \textbf{Oui} \\ \hline
        \textbf{Active Artificielle} & \textbf{Vaccination} (VPI, VPO, BCG, Hépatite B, HPV, COVID-19) & Antigène atténué/inactivé/sous-unitaire/ARNm $\rightarrow$ Réponse immune & Longue / \textbf{Oui} (rappels parfois) \\ \hline
        \textbf{Passive Naturelle} & \textbf{Transfert materno-fœtal (IgG placentaire)}, Allaitement (IgA lait) & Transfert d'anticorps maternels & Courte (mois) / \textbf{Non} \\ \hline
        \textbf{Passive Artificielle} & \textbf{Sérumthérapie} (Antivenimeux, Antitétanique, Anti-Rabique, Anti-D), Immunoglobulines IV & Injection d'IgG purifiées (donneur hyperimmunisé) & Immédiate, Courte (semaines) / \textbf{Non} \\ \hline
        \end{tabular}
        \end{center}
    \end{enumerate}
    \textbf{Reflexe Probatoire :} \og Immunité acquise \fg = \og Immunité adaptative \fg = \og Immunité spécifique \fg. \og Immunité innée \fg = \og Immunité non-spécifique \fg. \textbf{La vaccination est TOUJOURS une immunité active artificielle.} Ne jamais dire \og le vaccin donne des anticorps \fg (c'est le sérum), mais \og le vaccin \textbf{stimule la production} d'anticorps et de lymphocytes mémoire \fg.
\end{methode}

\begin{exoresolu}[Classification des situations immunologiques]
    \textbf{Énoncé :} Classe chacune des situations suivantes dans le tableau à double entrée ci-dessous en cochant la case correspondante. Justifie brièvement chaque classement.
    \begin{itemize}
        \item Situation 1 : Un nourrisson de 2 mois reçoit le vaccin pentavalent (DTC-HepB-Hib) au Centre de Santé Intégré de son quartier.
        \item Situation 2 : Une femme enceinte de 8 mois transmet des anticorps IgG à son fœtus via le placenta.
        \item Situation 3 : Un agriculteur de la région de l'Ouest se fait mordre par un serpent (vipère) ; on lui injecte un sérum antivenimeux polyvalent à l'hôpital de district.
        \item Situation 4 : Un enfant contracte la varicelle à l'école, guérit spontanément en 10 jours.
        \item Situation 5 : Les macrophages alvéolaires phagocytent des bacilles tuberculeux inhalés.
    \end{itemize}
    \begin{center}
    \begin{tabular}{|c|c|c|c|}\hline
    & \textbf{Active} & \textbf{Passive} & \textbf{Justification (1 phrase)} \\ \hline
    \textbf{Naturelle} & & & \\ \hline
    \textbf{Artificielle} & & & \\ \hline
    \end{tabular}
    \end{center}
    \tcblower
    \textbf{Solution détaillée :}
    \begin{center}
    \begin{tabular}{|c|c|c|c|}\hline
    & \textbf{Active} & \textbf{Passive} & \textbf{Justification} \\ \hline
    \textbf{Naturelle} & \textbf{Sit 4} (Varicelle) & \textbf{Sit 2} (Transfert maternel) & Sit 4 : Infection naturelle $\rightarrow$ réponse propre + mémoire. Sit 2 : Transfert physiologique IgG maternelles (placenta), pas de production fœtale, pas de mémoire. \\ \hline
    \textbf{Artificielle} & \textbf{Sit 1} (Vaccin pentavalent) & \textbf{Sit 3} (Sérum antivenimeux) & Sit 1 : Injection d'antigènes (toxines anatoxinées, Ag de surface HBs, capsulaire Hib, coque VPI) $\rightarrow$ stimulation réponse active + mémoire. Sit 3 : Injection d'IgG hétérologues (cheval) $\rightarrow$ neutralisation immédiate venin, pas de production patiente, pas de mémoire. \\ \hline
    \end{tabular}
    \end{center}
    \textbf{Analyse de la Situation 5 (Bonus) :} Les macrophages alvéolaires phagocytant \emph{Mycobacterium tuberculosis} relèvent de l'\textbf{immunité innée (non-spécifique)}. Ce n'est pas de l'immunité acquise (pas de spécificité antigénique fine, pas de mémoire lymphocytaire à ce stade). C'est la première ligne de défense cellulaire.
    \textbf{Conclusion :} La distinction Active/Passive repose sur \og qui produit les effecteurs ? \fg (soi-même vs donneur). La distinction Naturelle/Artificielle repose sur \og intervention médicale ? \fg. La vaccination est le prototype de l'immunité active artificielle.
\end{exoresolu}

\begin{methode}[Justifier l'intérêt de la vaccination (individuel et collectif)]
    \textbf{Objectif :} Construire une argumentation structurée reliant le mécanisme vaccinal à la protection individuelle et à l'épidémiologie (immunité collective).
    \begin{enumerate}
        \item \textbf{Rappeler le principe vaccinal :} Administration d'un \cle{antigène vaccinal} (agent atténué, inactivé, sous-unité, toxoïde, ARNm, vecteur viral) $\rightarrow$ mime l'infection \textbf{sans provoquer la maladie} (inocuité) $\rightarrow$ déclenche une \textbf{immunité active artificielle} (réponse primaire vaccinale : production d'anticorps + lymphocytes mémoire T et B).
        \item \textbf{Argumenter la protection individuelle :}
        \begin{itemize}
            \item \textbf{Réponse secondaire immédiate :} Lors de la rencontre ultérieure avec le pathogène sauvage, les lymphocytes mémoire permettent une réponse en 2-3 jours (vs 7-10 jours), taux d'anticorps 10-100x plus élevé, IgG haute affinité.
            \item \textbf{Prévention de la maladie :} Neutralisation du pathogène avant qu'il n'atteigne la dose infectieuse critique ou n'atteigne les organes cibles.
            \item \textbf{Éviter les séquelles :} Polio (paralysie), Rougeole (encéphalite, surdité), Tétanos (décès), HPV (cancer col utérin), Hépatite B (cirrhose/cancer foie).
        \end{itemize}
        \item \textbf{Argumenter la protection collective (Immunité de groupe / Herd Immunity) :}
        \begin{itemize}
            \item \textbf{Seuil de couverture vaccinale ($S_c$) :} $S_c = 1 - \frac{1}{R_0}$ où $R_0$ = nombre de reproduction de base (contagiosité).
            \item Si couverture $> S_c$ : La chaîne de transmission est rompue. Le pathogène ne trouve plus assez d'hôtes sensibles.
            \item \textbf{Protection des non-vaccinables :} Nourrissons trop jeunes, immunodéprimés (VIH, chimiothérapie), allergiques graves, femmes enceintes (pour vaccins vivants atténués).
            \item \textbf{Éradication / Élimination :} Ex: Variole (éradiquée 1980), Polio (éliminée Afrique 2020, endémique seulement Pakistan/Afghanistan), Rougeole (objectif élimination).
        \end{itemize}
        \item \textbf{Citer le calendrier vaccinal camerounais (PEV - Programme Élargi de Vaccination) :} BCG (naissance), VPO (naissance, 6, 10, 14 sem), Pentavalent (DTC-HepB-Hib) (6, 10, 14 sem), VAR (Rougeole-Rubéole) (9, 15 mois), Fièvre Jaune (9 mois), Vitamine A (6-59 mois), HPV (filles 9-14 ans, 2 doses), COVID-19 (adultes), Méningocoque A (MenAfriVac, ceinture méningite), Choléra (zones à risque), Rotavirus (introduction progressive).
    \end{enumerate}
    \textbf{Reflexe Probatoire :} Ne jamais écrire \og le vaccin tue le microbe \fg. Écrire : \og le vaccin \textbf{induit une mémoire immunitaire} permettant une élimination rapide du microbe lors de l'infection naturelle \fg. Distinguer \textbf{couverture vaccinale} (\% vaccinés dans la population) et \textbf{efficacité vaccinale} (\% protection chez le vacciné).
\end{methode}

\begin{exoresolu}[Argumentation sur la vaccination et immunité collective]
    \textbf{Énoncé :} Dans un district de santé de la région de l'Extrême-Nord, la couverture vaccinale contre la rougeole (VAR 1\textsuperscript{re} dose à 9 mois) est de 65\% depuis 3 ans. Une épidémie de rougeole survient, touchant principalement des enfants non vaccinés de 8 à 14 mois, mais aussi quelques enfants vaccinés et des nourrissons de 4 mois.
    \begin{enumerate}
        \item Calcule le seuil théorique d'immunité collective pour la rougeole sachant que $R_0 \approx 15$. La couverture actuelle est-elle suffisante ?
        \item Explique pourquoi des nourrissons de 4 mois (non éligibles au vaccin) tombent malades.
        \item Explique pourquoi quelques enfants vaccinés (à 9 mois) attrapent quand même la rougeole (échec vaccinal).
        \item Propose deux mesures concrètes pour contrôler l'épidémie et prévenir la suivante, en t'appuyant sur la stratégie OMS/PEV.
    \end{enumerate}
    \tcblower
    \textbf{Solution détaillée :}
    \begin{enumerate}
        \item \textbf{Calcul du seuil ($S_c$) :}
        \[ S_c = 1 - \frac{1}{R_0} = 1 - \frac{1}{15} = 1 - 0,067 = 0,933 = \textbf{93,3\%} \]
        La couverture actuelle (65\%) est \textbf{très inférieure} au seuil requis (93,3\%). L'immunité collective n'est pas assurée, la circulation du virus sauvage est maintenue, d'où l'épidémie.
        \item \textbf{Nourrissons de 4 mois :} Ils ne sont pas encore vaccinés (1\textsuperscript{re} dose à 9 mois). Ils ont perdu les \textbf{anticorps maternels (IgG placentaires)} acquis par immunité passive naturelle, dont le taux a chuté sous le seuil protecteur vers 3-6 mois (demi-vie $\approx$ 3 semaines). Ils sont \textbf{totalement sensibles} (fenêtre de susceptibilité). Ils sont protégés \textbf{uniquement} par l'immunité collective (absence de circulation du virus), qui fait défaut ici.
        \item \textbf{Échecs vaccinaux (enfants vaccinés malades) :} Deux causes principales :
        \begin{itemize}
            \item \textbf{Échec primaire (non-réponse) :} $\approx$ 5-15\% des vaccinés à 9 mois ne séroconvertissent pas. Causes : persistance d'anticorps maternels inhibant le vaccin vivant atténué, malnutrition, immunodépression, chaîne du froid rompue (vaccin inactivé), administration incorrecte.
            \item \textbf{Échec secondaire (perte d'immunité) :} Rareté pour la rougeole (immunité généralement vie), mais possible si immunodépression supervenue.
        \end{itemize}
        \item \textbf{Mesures de contrôle et prévention (Stratégie OMS/PEV Cameroun) :}
        \begin{enumerate}
            \item \textbf{Vaccination de riposte (Campagne de rattrapage) :} Organisation immédiate d'une campagne de vaccination de masse \textbf{VAR (Rougeole-Rubéole)} pour tous les enfants de \textbf{6 mois à 14 ans} (ou 5-15 ans selon épidémiologie) dans le district, \textbf{indépendamment du statut vaccinal antérieur}. Dose à 6-8 mois = dose "zéro" (ne compte pas pour le calendrier, 2\textsuperscript{e} dose à 9 mois obligatoire).
            \item \textbf{Renforcement de la vaccination de routine (PEV) :} Atteindre et maintenir \textbf{$\geq$ 95\% de couverture} pour les 2 doses de VAR (9 mois et 15 mois) dans toutes les aires de santé. Stratégies : Vaccination avancée (mobile), vaccination à domicile, sensibilisation leaders d'opinion/religieux, gestion des refus, surveillance active des cas (déclaration obligatoire, investigation, prélèvement sanguin/salivaire pour confirmation IgM/PCR).
            \item (Complément) \textbf{Prise en charge des cas :} Vitamine A (2 doses à 24h d'intervalle) pour réduire mortalité/séquelles oculaires, antibiotiques si surinfection bactérienne, isolement, nutrition.
        \end{enumerate}
    \end{enumerate}
    \textbf{Conclusion :} Une couverture vaccinale inférieure au seuil d'immunité collective ($1-1/R_0$) expose la population aux épidémies, touchant en priorité les non-vaccinés (nourrissons, refus) et révélant les échecs vaccinaux. La riposte combine vaccination de masse immédiate et renforcement durable de la routine.
\end{exoresolu}

\begin{methode}[Analyser les facteurs affaiblissant le système immunitaire]
    \textbf{Objectif :} Relier un facteur de risque (nutritionnel, psychologique, pathologique, toxique) à son mécanisme d'action sur les composantes du système immunitaire (cellules, molécules, organes).
    \begin{enumerate}
        \item \textbf{Malnutrition (Protéino-énergétique et Micronutriments) :} \textbf{1\textsuperscript{re} cause d'immunodépression secondaire dans le monde (dont Cameroun).}
        \begin{itemize}
            \item \textbf{Atrophie thymique :} Perte de poids $\rightarrow$ involution thymique $\rightarrow$ baisse production LT naïfs (surtout LT4).
            \item \textbf{Déficit phagocytaire :} Baisse chimiotaxie, phagocytose, burst oxydatif (manque ATP, NADPH).
            \item \textbf{Atteinte de l'immunité humorale :} Baisse production IgA (muqueuses $\rightarrow$ diarrhées, infections respiratoires), IgG. Atrophie ganglions/rate.
            \item \textbf{Micronutriments clés :} Vitamine A (intégrité épithéliums, différenciation LT, IgA), Zinc (thymuline, maturation LT, activité thymidine kinase), Fer (prolifération lymphocytaire, burst oxydatif), Sélénium, Vit D, Vit E, Vit B6, B9, B12, Cuivre.
        \end{itemize}
        \item \textbf{Stress chronique / Cortisol :} Activation axe HPA $\rightarrow$ Cortisol (glucocorticoïde).
        \begin{itemize}
            \item \textbf{Lymphopénie :} Redistribution LT (marge vasculaire $\rightarrow$ moelle/ganglions), apoptose thymocytes/LT périphériques.
            \item \textbf{Inhibition cytokine :} Baisse IL-2, IFN-$\gamma$, TNF-$\alpha$ (Th1) $\rightarrow$ bascule vers profil Th2 (allergie, faible défense intracellulaire).
            \item \textbf{Inhibition NK / Macrophages :} Baisse cytotoxicité, présentation antigénique.
        \end{itemize}
        \item \textbf{Maladies infectieuses (Immunodépression acquise) :}
        \begin{itemize}
            \item \textbf{VIH/Sida :} Destruction ciblée des \textbf{LT4/CD4+} (récepteur CD4 = co-récepteur viral). Effondrement de l'aide lymphocytaire $\rightarrow$ anergie LB, défaut activation LT8, défaut activation macrophages. Infections opportunistes (TB, Pneumocystose, Toxoplasmose, Cryptococcose, Candidose).
            \item \textbf{Paludisme grave / Sepsis :} "Paralysie immunitaire" transitoire (tolérance aux endotoxines, désexpression HLA-DR sur monocytes, apoptose lymphocytaire massive).
            \item \textbf{Rougeole :} Immunosuppression prolongée (mois/années) par "amnésie immunitaire" (perte lymphocytes mémoire préexistants).
        \end{itemize}
        \item \textbf{Autres facteurs :} Vieillissement (immunosenescence : involution thymique, répertoire TCR restreint, inflammation chronique "inflammaging"), Toxiques (Alcool : déprime phagocytes, IgA ; Tabac : paralysie ciliaire, déprime macrophages alvéolaires, cancérigène), Médicaments (Corticoïdes, Chimiothérapie, Immunosuppresseurs post-greffe), Splénectomie (vu plus haut).
    \end{enumerate}
    \textbf{Reflexe Probatoire :} Lien direct \textbf{Malnutrition $\leftrightarrow$ Infections} (cercle vicieux). Au Cameroun, le paludisme et la malnutrition sont les deux piliers de la morbidité infantile qui s'entretiennent. Citer \textbf{VIH} comme cause majeure d'immunodépression cellulaire (CD4).
\end{methode}

\begin{exoresolu}[Étude de cas : Immunodépression et cercle vicieux]
    \textbf{Énoncé :} Amadou, 18 mois, vit dans un village de la région du Nord. Il présente un retard de croissance staturo-pondéral (PC < -2 DS, Taille < -2 DS), une anémie ferriprive (Hb = 8 g/dL), et fait des épisodes répétés de diarrhées aiguës et de pneumonies bactériennes depuis 6 mois. Sa mère allaite encore mais l'alimentation de complément est pauvre (bouillie de maïs seule). Le test VIH est négatif.
    \begin{enumerate}
        \item Identifie le diagnostic nutritionnel principal et cite deux carences en micronutriments probables chez Amadou.
        \item Explique, au niveau cellulaire et moléculaire, pourquoi Amadou fait des infections respiratoires et digestives à répétition (mécanismes de l'immunodépression nutritionnelle).
        \item Le médecin prescrit un traitement antibiotique pour la pneumonie actuelle et du Plumpy'Nut® (Aliment Thérapeutique Prêt à l'Emploi - ATPE) pour la malnutrition. Justifie l'intérêt de la supplémentation en Zinc et Vitamine A dans l'ATPE pour la récupération immunitaire.
    \end{enumerate}
    \tcblower
    \textbf{Solution détaillée :}
    \begin{enumerate}
        \item \textbf{Diagnostic :} \textbf{Malnutrition aiguë modérée (MAM) ou sévère (MAS) associée à un retard de croissance chronique (stunting).} PC < -2 DS + Taille < -2 DS = malnutrition globale (aiguë sur chronique).
        \textbf{Carences probables :} \textbf{Fer} (anémie microcytaire hypochrome Hb 8g/dL), \textbf{Zinc} (diarrhées, dermatose, retard croissance, immunodépression), \textbf{Vitamine A} (risque xérophtalmie, immunité muqueuse), \textbf{Iode} (goitre région Nord), \textbf{Vitamine D}.
        \item \textbf{Mécanismes de l'immunodépression nutritionnelle chez Amadou :}
        \begin{itemize}
            \item \textbf{Barrières physiques chimiques altérées :} Atrophie des villosités intestinales (malabsorption, porte d'entrée bactérienne), perte cils respiratoires, sécheresse conjonctivale/cutanée (carence Vit A) $\rightarrow$ perte effet barrière.
            \item \textbf{Immunité innée défaillante :} 
                \begin{itemize}
                    \item \textbf{Phagocytes (Neutrophiles/Macrophages) :} Chimiotaxie réduite, phagocytose diminuée, \textbf{burst oxydatif altéré} (manque NADPH, Fer, Zinc pour superoxyde dismutase). Moins de tueurs intracellulaires.
                    \item \textbf{Complément :} Baisse synthèse hépatique des facteurs du complément (protéines).
                    \item \textbf{NK :} Cytotoxicité diminuée.
                \end{itemize}
            \item \textbf{Immunité adaptative effondrée :} 
                \begin{itemize}
                    \item \textbf{Thymus :} Involution thymique marquée (atrophie graisseuse) $\rightarrow$ \textbf{baisse production LT naïfs (CD4+ et CD8+)}. Répertoire TCR appauvri.
                    \item \textbf{Lymphoïdes périphériques :} Atrophie ganglions, rate, plaques de Peyer $\rightarrow$ zones T et B hypoplasiques.
                    \item \textbf{LB / Anticorps :} Baisse production \textbf{IgA sécrétoire} (muqueuses intestinale/respiratoire $\rightarrow$ porte ouverte aux bactéries/virus), baisse IgG sérique. Réponse vaccinale altérée.
                    \item \textbf{LT4 :} Lymphopénie T, production d'IL-2 diminuée $\rightarrow$ pas d'aide pour LB ni activation macrophages.
                \end{itemize}
            \item \textbf{Cercle vicieux :} Infections $\rightarrow$ anorexie, catabolisme, pertes protéiques (diarrhée), séquestration fer (hépcidine) $\rightarrow$ aggravation malnutrition $\rightarrow$ aggravation immunodépression.
        \end{itemize}
        \item \textbf{Justification Zinc et Vitamine A dans l'ATPE (Plumpy'Nut®) :}
        \begin{itemize}
            \item \textbf{Zinc (Zn) :} 
                \begin{itemize}
                    \item Cofacteur de la \textbf{thymuline} (hormone thymique indispensable à la maturation/différenciation des LT).
                    \item Essentiel pour l'activité de la \textbf{thymidine kinase} (synthèse ADN $\rightarrow$ prolifération clonale lymphocytaire).
                    \item Stabilise les membranes cellulaires, cofacteur superoxyde dismutase (Cu/Zn-SOD) $\rightarrow$ protection phagocytes contre leur propre burst oxydatif.
                    \item \textbf{Résultat :} Restauration de la thymopoïèse, de la prolifération lymphocytaire, de la cytotoxicité NK, de l'activité phagocytaire. Réduit durée/gravité diarrhées et pneumonies (recommandation OMS : 20 mg/jour 10-14 jours lors diarrhée aiguë).
                \end{itemize}
            \item \textbf{Vitamine A (Rétinol) :} 
                \begin{itemize}
                    \item Maintien de l'\textbf{intégrité des épithéliums} (différenciation cellulaire $\rightarrow$ muqueuses respiratoire, digestive, oculaire fonctionnelles = barrière).
                    \item Différenciation des \textbf{LT} (régulation gènes via récepteurs nucléaires RAR/RXR).
                    \item Stimulation production \textbf{IgA} (immunité muqueuse).
                    \item Activation macrophages / NK.
                    \item \textbf{Résultat :} Réduction mortalité rougeole (50\%), diarrhée (30\%), paludisme. Dose thérapeutique OMS : 200 000 UI (6-12 mois : 100 000 UI) J1, J2, J14.
                \end{itemize}
        \end{itemize}
    \end{enumerate}
    \textbf{Conclusion :} La malnutrition protéino-énergétique et micronutritionnelle est une cause majeure d'immunodépression secondaire globale (innée + adaptative, barrières + cellules + anticorps), créant un cercle vicieux infection-malnutrition. La réhabilitation nutritionnelle (ATPE) corrige les carences spécifiques (Zn, Vit A, Fer, Protéines) indispensables à la reconstruction du système immunitaire.
\end{exoresolu}

\begin{methode}[Proposer des comportements favorables au renforcement de l'immunité (Éducation à la santé)]
    \textbf{Objectif :} Formuler des recommandations concrètes, scientifiquement fondées, adaptées au contexte de vie (Cameroun), classées par leviers d'action.
    \begin{enumerate}
        \item \textbf{Levier Nutritionnel (Assiette immunitaire) :}
        \begin{itemize}
            \item \textbf{Allaitement maternel exclusif 6 mois} (IgA, lysozyme, lactoferrine, oligosaccharides, cellules maternelles, facteurs de croissance) $\rightarrow$ poursuite jusqu'à 2 ans + alimentation de complément diversifiée.
            \item \textbf{Alimentation variée, équilibrée, locale et de saison :} 
                \begin{itemize}
                    \item Protéines de qualité (Poisson \emph{Kanga}/\emph{Banga}, viande, œufs, légumineuses \emph{haricots/niébé/soja}, insectes comestibles \emph{larves de palmier/cricket}) $\rightarrow$ acides aminés pour synthèse anticorps/lymphocytes.
                    \item Fruits/Légumes colorés (Mangue, Papaye, Goyave, Ananas, Carotte, Feuxilles \emph{Ndolé}, \emph{Eru}, \emph{Folong}, Tomate, Poivron) $\rightarrow$ \textbf{Vit A, C, B9, Polyphénols} (antioxydants, prolifération lymphocytaire).
                    \item Céréales complètes / Tubercule (Maïs, Mil, Sorgho, Manioc, Igname, Patate douce, Plantain, Taro) $\rightarrow$ Énergie, Fibres (microbiote), Zinc, Fer, Vit B.
                    \item Sources de Zinc/Fer/Sélénium : Abats (foie, rein), fruits de mer (crevettes séchées \emph{crevettes de rivière}), graines (courge, sésame), noix de cajou.
                \end{itemize}
            \item \textbf{Hygiène alimentaire :} Eau potable (traitement Aquatabs/ébullition/filtration), lavage mains au savon (moments clés), cuisson complète viandes/poissons, conservation au frais.
        \end{itemize}
        \item \textbf{Levier Vaccinal (Protection active artificielle) :}
        \begin{itemize}
            \item Respect strict du \textbf{Calendrier PEV Cameroun} (enfants 0-23 mois, filles 9-14 ans HPV, femmes enceintes VAT/VAT2, COVID-19).
            \item Rattrapage vaccinal si retard. Carte de vaccination tenue à jour (preuve pour école, voyage).
            \item Vaccins recommandés hors PEV selon risque : Méningocoque (Nord, pèlerins), Fièvre typhoïde, Hépatite A, Rage (post-exposition), Choléra (zones endémiques), Grippe saisonnière (personnes à risque).
        \end{itemize}
        \item \textbf{Levier Hygiène / Environnement (Réduction charge antigénique) :}
        \begin{itemize}
            \item \textbf{Lavage des mains au savon :} Après toilettes, avant manger/nourrir, après change, après marché/transport. \textbf{Mesure n°1 prévention infections diarrhéiques/réspiratoires.}
            \item Assainissement : Latrines améliorées (VIP), gestion déchets, lutte anti-vectorielle (Moustiquaires Imprégnées à Longue Durée d'Action - MILDA gratuites au Cameroun, destruction gîtes larvaires).
            \item Qualité air : Aération logements, cuisines extérieures ou hottes (fumée bois/charbon = BPCO, cancer, immunodépression locale).
        \end{itemize}
        \item \textbf{Levier Mode de vie / Psychosocial :}
        \begin{itemize}
            \item \textbf{Sommeil suffisant et régulier :} 9-11h (enfant), 8-10h (ado), 7-9h (adulte). Sommeil = pic sécrétion GH, prolactine, mélatonine (immunostimulants), réparation cellulaire, consolidation mémoire immunitaire.
            \item \textbf{Activité physique modérée régulière :} Marche rapide 30 min/jour, sport scolaire, travaux champs. Stimule circulation lymphatique, recrutement NK/neutrophiles, réduit inflammation chronique, gère stress. \textbf{Éviter surentraînement} (fenêtre d'immunosuppression post-effort intense).
            \item \textbf{Gestion du stress :} Soutien familial, communautaire, spirituel. Éviter isolement, violences. Stress chronique = Cortisol $\uparrow$ = Immunité $\downarrow$.
            \item \textbf{Évitement toxiques :} \textbf{Zéro tabac} (actif/passif), \textbf{Alcool modéré/abstinence} (pas d'alcool < 18 ans, femme enceinte), \textbf{Non aux drogues} (Tramadol, cannabis, cocaïne $\rightarrow$ immunodépression sévère).
        \end{itemize}
        \item \textbf{Levier Médical / Suivi :}
        \begin{itemize}
            \item Dépistage précoce VIH (PMTCT, auto-test, CD4/VL si +), Tuberculose, Diabète, HTA.
            \item Prise en charge correcte infections (Observance antibiotiques, antipaludéens ACT).
            \item Supplémentation ciblée si carence avérée (Vit A 6-59 mois, Fer/Folate femme enceinte, Zinc diarrhée, Vit D si faible ensoleillement/enfermement).
        \end{itemize}
    \end{enumerate}
    \textbf{Reflexe Probatoire :} Les recommandations doivent être \textbf{spécifiques, mesurables, atteignables, réalistes, temporelles (SMART)}. Ex: \og Manger équilibré \fg (trop vague) $\rightarrow$ \og Consommer 5 portions fruits/légumes locaux par jour, source protéine à chaque repas \fg. Toujours lier le comportement au mécanisme immunitaire (ex: \og Lavage mains $\rightarrow$ élimination mécanique pathogènes $\rightarrow$ baisse charge antigénique $\rightarrow$ épargne système immunitaire \fg).
\end{methode}

\begin{exoresolu}[Conception d'un message d'éducation pour la santé]
    \textbf{Énoncé :} Tu es infirmier(ère) dans un Centre de Santé Intégré (CSI) en zone rurale dans la région du Centre. Tu dois animer une causerie éducative de 20 minutes pour un groupe de 15 mères d'enfants de 0 à 23 mois, sur le thème : \og Comment protéger naturellement les défenses de mon enfant ? \og. Le taux de couverture vaccinale Pentavalent 3 est de 72\% dans l'aire de santé (objectif 90\%). La malnutrition aiguë globale est à 8\%.
    \begin{enumerate}
        \item Propose un plan structuré de ta causerie en 4 temps (Accroche, Corps du message, Démonstration/Pratique, Engagement/Clôture), en indiquant pour chaque temps : le message clé, l'argument scientifique simplifié, et le support/outil utilisé.
        \item Parmi les comportements proposés, lesquels relèvent de l'immunité active artificielle, de l'immunité passive naturelle, et du renforcement de l'immunité innée/adaptative par le mode de vie ?
        \item Comment évalueras-tu l'impact de ta causerie à 3 mois et à 1 an ?
    \end{enumerate}
    \tcblower
    \textbf{Solution détaillée :}
    \begin{enumerate}
        \item \textbf{Plan de la causerie (Approche participative, langage local + français simple) :}
        \begin{center}
        \begin{tabularx}{\linewidth}{|l|X|X|X|}\hline
        \textbf{Temps} & \textbf{Message Clé (1 phrase)} & \textbf{Argument Scientifique Simplifié} & \textbf{Support / Outil} \\ \hline
        \textbf{1. Accroche (3 min)} & \og Vos enfants tombent moins malades s'ils ont une armée forte dans le corps. \fg & Analogie : Système immunitaire = \textbf{Armée de défense} (Soldats = globules blancs, Armes = anticorps, Casernes = rate/ganglions/thymus, Renseignement = mémoire). Si soldats affamés/fatigués $\rightarrow$ ennemi (microbe) gagne. & Affiche "Armée du corps" (schéma simplifié corps enfant avec icônes rate, ganglion, globules blancs). Question : \og Qui a eu un enfant malade ce mois-ci ? \fg \\ \hline
        \textbf{2. Corps : Les 4 Piliers (12 min)} & \textbf{Pilier 1 : Le Bouclier Vaccin (Immunité Active Artificielle).} \og Le vaccin entraine l'armée sans la blesser. \fg & Vaccin = \textbf{exercice de tir à blanc}. Montre l'ennemi (antigène) $\rightarrow$ armée fabrique armes (Ac) et officiers mémoire. Si vrai ennemi arrive $\rightarrow$ riposte immédiate (J3 vs J10). Protège l'enfant ET le village (immunité collective). & Carnet de santé (carte vaccinale) : Montrer où cocher. Calendrier PEV affiché. Témoignage mère "enfant vacciné vs non vacciné". \\ \cline{2-4}
         & \textbf{Pilier 2 : L'Assiette du Soldat (Nutrition).} \og Manger varié et coloré = armes solides. \fg & Protéines (viande/poisson/œufs/légumineuses) = briques soldats/anticorps. Fruits/Légumes (Mangue, Papaye, Ndolé, Carotte) = Vit A/C = boucliers muqueuses + boost soldats. Zinc (foie, graines) = chef d'orchestre. Fer (feuilles, viande) = carburant. Allaitement 6 mois exclusif = 1er vaccin naturel (IgA maman). & Assiette réelle "Assiette idéale camerounaise" (1/2 légumes, 1/4 protéines, 1/4 féculents + fruit). Démonstration bouillie enrichie (farine maïs/soja + huile rouge + moringa/feuilles + poisson séché pilé). \\ \cline{2-4}
         & \textbf{Pilier 3 : L'Hygiène du Campement (Barrières Innées).} \og Mains propres = ennemi bloqué à la porte. \fg & Lavage mains savon = \textbf{arme mécanique n°1}. Enlève 90\% microbes avant qu'ils n'entrent (bouche, nez, yeux). Eau traitée (Aquatabs/bouillir) = élimine microbes invisibles. Latrines propres = évite contamination sol/eau. MILDA = bloque moustiques (paludisme = fatigue armée). & Démonstration lavage mains (technique OMS 6 étapes) avec seau/robinet "Tippy Tap" local. Flacon Aquatabs. Moustiquaire bien tendue. \\ \cline{2-4}
         & \textbf{Pilier 4 : Le Repos du Guerrier (Mode de vie).} \og Bien dormir, bouger, pas de fumée = armée au top. \fg & Sommeil = usine à soldats (hormones croissance/réparation). Jeu/Activité = circulation soldats (lymphe). Fumée cuisine/tabac = blesse soldats poumons (macrophages paralysés). Stress/Chagrin = hormone cortisol = endort soldats. & Image enfant dormant sous MILDA. Photo cuisine extérieure / foyer amélioré. Conseil : \og Pas de fumée dans la case où dort bébé \fg. \\ \hline
        \textbf{3. Démonstration / Pratique (3 min)} & \og Je sais préparer la bouillie enrichie et laver les mains. \fg & Mise en pratique supervisée par les mères (binômes). Correction par les pairs. & Ingrédients locaux, bols, cuillères, savon, seau Tippy Tap. Fiche recette illustrée remise. \\ \hline
        \textbf{4. Engagement / Clôture (2 min)} & \og Je m'engage pour mon enfant : Carte à jour, Assiette colorée, Mains propres, Sommeil. \fg & Engagement public = levier psychologique (cohérence, norme sociale). Rappel RDV vaccin prochain (Penta 3, VAR). & \textbf{Carnet d'engagement} (petite carte signée/empreinte pouce) avec 4 cases à cocher par mois. Rappel SMS/Whatsapp groupe mères. Chanson/Comptine "L'armée de mon corps". \\ \hline
        \end{tabularx}
        \end{center}
        \item \textbf{Classification des comportements proposés :}
        \begin{itemize}
            \item \textbf{Immunité Active Artificielle :} \textbf{Vaccination (Pentavalent, VAR, VPO, BCG, Fièvre Jaune, HPV, COVID)}. $\rightarrow$ Stimulation production propre Ac + mémoire.
            \item \textbf{Immunité Passive Naturelle :} \textbf{Allaitement maternel exclusif 6 mois (IgA, IgG, cellules, facteurs immuns du colostrum/lait)}. $\rightarrow$ Transfert effecteurs maternels, protection transitoire, pas de mémoire enfant.
            \item \textbf{Renforcement Immunité Innée (Barrières, Phagocytes, NK, Complément, Inflammation) :} 
                \begin{itemize}
                    \item Lavage mains, eau potable, latrines, MILDA $\rightarrow$ \textbf{Baissent la charge antigénique} (évitent de solliciter l'immunité).
                    \item Vitamine A (intégrité épithéliums), Zinc (phagocytose, NK, thymuline), Fer (burst oxydatif), Protéines (complément, cellules) $\rightarrow$ \textbf{Optimisent la fonction effectrice innée}.
                    \item Sommeil, exercice modéré, non-tabac $\rightarrow$ \textbf{Régulent l'inflammation, maintiennent le pool cellulaire}.
                \end{itemize}
            \item \textbf{Renforcement Immunité Adaptative (Lymphocytes T/B, Mémoire, Anticorps) :} 
                \begin{itemize}
                    \item Vaccination (voir actif artificielle).
                    \item Nutrition (Protéines, Zn, Vit A, B9, B12, Fer) $\rightarrow$ \textbf{Permettent prolifération clonale, différenciation plasmocytes, commutation de classe, maturation affinité}.
                    \item Sommeil $\rightarrow$ \textbf{Consolide la mémoire immunologique} (études : sommeil post-vaccination $\uparrow$ taux Ac).
                \end{itemize}
        \end{itemize}
        \item \textbf{Évaluation de l'impact :}
        \begin{itemize}
            \item \textbf{À 3 mois (Processus / Résultats intermédiaires) :}
                \begin{itemize}
                    \item \textbf{Couverture vaccinale Penta 3 / VAR 1 :} Extraction données DHIS2 / Registre CSI $\rightarrow$ Objectif $\geq$ 90\%.
                    \item \textbf{Observance lavage mains :} Enquête KAP (Knowledge, Attitudes, Practices) sur échantillon 30 mères (questionnaire + observation point d'eau/savon).
                    \item \textbf{Prévalence malnutrition aiguë (PB/Taille) :} Pesée mensuelle CSI / Dépistage communautaire (MUAC/PB) $\rightarrow$ Objectif $<$ 5\%.
                    \item \textbf{Taux de complétion carnet d'engagement :} \% mères ramenant carte remplie.
                \end{itemize}
            \item \textbf{À 1 an (Impact / Santé) :}
                \begin{itemize}
                    \item \textbf{Incidence maladies cibles :} Nombre cas \textbf{Diarrhée aiguë, Pneumonie (IRAS), Paludisme, Rougeole} chez les 0-23 mois (Registre CSI / DHIS2) $\rightarrow$ Comparaison année N vs N-1. Objectif : Baisse $\geq$ 20\%.
                    \item \textbf{Mortalité infanto-juvénile (0-5 ans) :} Données district / ENEC / DHS si dispo.
                    \item \textbf{Couverture vaccinale complète (Penta 3 + VAR 1 + 2) :} Objectif 95\%.
                    \item \textbf{État nutritionnel moyen (Score Z Poids/Taille, Taille/Âge) :} Amélioration moyenne.
                \end{itemize}
        \end{itemize}
    \end{enumerate}
    \textbf{Conclusion :} L'éducation pour la santé efficace combine connaissances scientifiques (immunonutrition, vaccinologie), adaptation culturelle (aliments locaux, langues, croyances), participation active (démo, engagement) et suivi d'indicateurs concrets (couverture, morbidité, mortalité).
\end{exoresolu}

\begin{methode}[Analyser une courbe de réponse immunitaire (Kinetics)]
    \textbf{Objectif :} Interpréter graphiquement la dynamique des taux d'anticorps (IgM, IgG) et de lymphocytes lors d'une primo-infection vs ré-infection (ou vaccination).
    \begin{enumerate}
        \item \textbf{Identifier les axes :} 
        \begin{itemize}
            \item Abscisse (X) : \textbf{Temps} (jours, semaines, mois, années).
            \item Ordonnée (Y) : \textbf{Taux sérique} (mg/dL, UI/mL, log10 titer) ou \textbf{Nombre de cellules} (/mm³ ou /µL).
        \end{itemize}
        \item \textbf{Repérer les phases caractéristiques (Primo-réponse / Réponse primaire vaccinale) :}
        \begin{itemize}
            \item \textbf{Phase de latence (J0-J5/7) :} Taux basal (IgM $\approx$ 0, IgG = taux maternel résiduel ou 0). Activation/Proclifération clonale silencieuse.
            \item \textbf{Phase exponentielle / Montée (J5-J14) :} Apparition \textbf{IgM} en premier (pentamères, faible affinité, pas de commutation). Pic IgM vers J7-J10.
            \item \textbf{Phase de plateau / Pic :} Commutation de classe (IgM $\rightarrow$ IgG/IgA/IgE) sous l'effet cytokines (IL-4, IFN-$\gamma$).\textbf{IgG} deviennent majoritaires (monomères, haute affinité, passage placenta, fixation complément, opsonisation). Pic IgG vers J14-J21.
            \item \textbf{Phase de déclin (Catabolisme) :} Élimination plasmocytes effecteurs (apoptose). Baisse IgM rapide (demi-vie 5j), baisse IgG plus lente (demi-vie 21-28j).
            \item \textbf{Phase de mémoire (Niveau basal élevé) :} Persistance \textbf{lymphocytes B mémoire} et \textbf{plasmocytes à longue vie} (moelle osseuse) $\rightarrow$ taux IgG résiduel protecteur pendant années.
        \end{itemize}
        \item \textbf{Comparer avec Réponse Secondaire (Ré-infection / Rappel vaccinal) :}
        \begin{itemize}
            \item \textbf{Latence quasi nulle} (J1-J2).
            \item \textbf{Pente montée très raide} (logarithmique).
            \item \textbf{Pic 10-100x plus haut}.
            \item \textbf{IgG dominantes dès le début} (pas de pic IgM significatif ou très transitoire).
            \item \textbf{Affinité plus forte} (maturation affinité).
        \end{itemize}
        \item \textbf{Distinguer Vaccination (Primaire vs Rappel) :} 
        \begin{itemize}
            \item Primo-vaccination (1 dose ou 2-3 doses espacées 1-2 mois) = mimique primo-infection (courbe primaire).
            \item Rappel (Booster, mois/années après) = mimique réponse secondaire (courbe anamnestique).
        \end{itemize}
    \end{enumerate}
    \textbf{Reflexe Probatoire :} Sur un graphique, \textbf{légender les courbes (IgM, IgG)}. Citer les \textbf{valeurs numériques approximatives} (jours de latence, pic, ratio IgG/IgM). Ne pas confondre \og taux d'anticorps \fg et \og nombre de lymphocytes \fg (courbes parallèles mais décalées : lymphocytes montent avant anticorps).
\end{methode}

\begin{exoresolu}[Interprétation de courbes sérologiques post-vaccinales]
    \textbf{Énoncé :} Le graphique ci-dessous représente l'évolution du taux sérique d'anticorps anti-hépatite B (anti-HBs, en UI/mL) chez un adolescent de 15 ans vacciné selon le schéma 3 doses (J0, M1, M6). Le seuil de protection est fixé à 10 UI/mL.
    \begin{popfigure}
        \begin{tikzpicture}
            \begin{axis}[
                width=\linewidth, height=8cm,
                xlabel={Temps (Mois)}, ylabel={Taux anti-HBs (UI/mL) - Échelle log},
                xmin=0, xmax=14, ymin=0.1, ymax=10000,
                ymode=log, log basis y=10,
                xtick={0,1,2,3,4,5,6,7,8,9,10,11,12,13,14},
                ytick={1,10,100,1000,10000},
                grid=major, grid style={dashed, gray!50},
                legend pos=north west,
                legend cell align=left,
                clip=false
            ]
            % Seuil protection
            \addplot[dashed, thick, poporange, domain=0:14] {10};
            \addlegendentry{Seuil protection (10 UI/mL)}
            
            % Courbe IgG simulée (log scale)
            \addplot[popcurve, very thick, popblue, smooth, samples=200, domain=0:14] 
            coordinates {
                (0, 0.5) (0.5, 0.5) (1, 5) (1.5, 50) (2, 120) (2.5, 100) (3, 80) (4, 40) (5, 20) 
                (5.5, 20) (6, 20) (6.5, 200) (7, 1500) (7.5, 3000) (8, 2500) (9, 1500) (10, 800) (11, 500) (12, 350) (13, 250) (14, 200)
            };
            \addlegendentry{Anti-HBs (IgG)}
            
            % Flèches annotations
            \draw[popblue, thick, ->] (axis cs:0,0.5) -- (axis cs:0,0.1) node[below, font=\tiny, popblue] {Dose 1 (J0)};
            \draw[popblue, thick, ->] (axis cs:1,20) -- (axis cs:1,5) node[below, font=\tiny, popblue] {Dose 2 (M1)};
            \draw[popblue, thick, ->] (axis cs:6,20) -- (axis cs:6,5) node[below, font=\tiny, popblue] {Dose 3 (M6)};
        \end{axis}
        \end{tikzpicture}
        \legende{Évolution du taux d'anticorps anti-HBs (échelle logarithmique) après vaccination hépatite B 3 doses (J0, M1, M6). Seuil de protection = 10 UI/mL.}
    \end{popfigure}
    \begin{enumerate}
        \item Décris l'allure générale de la courbe et identifie les pics correspondant à chaque injection.
        \item Pourquoi l'échelle de l'ordonnée est-elle logarithmique ? Que montre-t-elle sur l'amplitude des réponses ?
        \item Calcule le délai approximatif pour atteindre le seuil de protection après la 1\textsuperscript{re} dose. L'adolescent est-il protégé entre le mois 1 et le mois 6 ?
        \item Explique la différence de cinétique et d'amplitude entre la réponse à la 1\textsuperscript{re} dose et celle à la 3\textsuperscript{e} dose (rappel).
        \item Quelle est l'utilité de la 2\textsuperscript{e} dose (M1) si la 1\textsuperscript{re} dose induit déjà une réponse ?
    \end{enumerate}
    \tcblower
    \textbf{Solution détaillée :}
    \begin{enumerate}
        \item \textbf{Description de la courbe :} La courbe montre une \textbf{échelle logarithmique} (chaque graduation $\times$10). On observe \textbf{trois pics majeurs} correspondant aux trois injections :
        \begin{itemize}
            \item \textbf{Pic 1 (vers M1.5-M2) :} Suite à la 1\textsuperscript{re} dose (J0). Montée lente, pic modéré ($\approx$ 100-120 UI/mL), déclin progressif.
            \item \textbf{Pic 2 (vers M2.5-M3) :} Suite à la 2\textsuperscript{e} dose (M1). Montée plus rapide, pic plus haut ($\approx$ 200-300 UI/mL), déclin.
            \item \textbf{Pic 3 (vers M7) :} Suite à la 3\textsuperscript{e} dose (M6). \textbf{Montée très rapide (réponse anamnestique)}, \textbf{pic très élevé ($\approx$ 3000 UI/mL)}, suivi d'un plateau de maintien élevé ($> 200$ UI/mL à M14).
        \end{itemize}
        \item \textbf{Intérêt de l'échelle logarithmique :} Les variations d'amplitude des taux d'anticorps couvrent plusieurs ordres de grandeur (de $<1$ à $>1000$ UI/mL). Une échelle linéaire écrabouillerait la primo-réponse (presque invisible) et ne montrerait que le pic du rappel. L'échelle log permet de \textbf{visualiser simultanément la cinétique (pentes) et l'amplitude relative (facteurs 10, 100, 1000)} de chaque réponse. Elle montre que le rappel (Dose 3) induit une réponse $\approx$ \textbf{30 fois plus forte} en pic que la primo-réponse (Dose 1).
        \item \textbf{Délai protection et statut M1-M6 :}
        \begin{itemize}
            \item Seuil 10 UI/mL franchi vers \textbf{J10-J12 (M0.3-M0.4)} après la 1\textsuperscript{re} dose.
            \item \textbf{Entre M1 et M6 :} Le taux reste \textbf{bien au-dessus de 10 UI/mL} (entre 20 et 100 UI/mL). L'adolescent \textbf{EST PROTÉGÉ} pendant cet intervalle. La 2\textsuperscript{e} dose à M1 booste le taux pour le maintenir haut et sélectionner les clones à haute affinité.
        \end{itemize}
        \item \textbf{Différence Dose 1 vs Dose 3 (Réponse Primaire vs Secondaire/Anamnestique) :}
        \begin{center}
        \begin{tabular}{|l|c|c|}\hline
        \textbf{Paramètre} & \textbf{1\textsuperscript{re} Dose (Primaire)} & \textbf{3\textsuperscript{e} Dose (Rappel / Secondaire)} \\ \hline
        \textbf{Cellules initiatrices} & LB naïfs (faible fréquence, récepteurs faible affinité) & \textbf{LB Mémoire} (haute fréquence, récepteurs haute affinité) + Plasmocytes longue vie résiduels \\ \hline
        \textbf{Délai (Latence)} & Long (7-14 jours) & \textbf{Très court (1-3 jours)} \\ \hline
        \textbf{Pente de montée} & Faible (linéaire/log) & \textbf{Très raide (exponentielle)} \\ \hline
        \textbf{Pic d'anticorps} & Modeste ($\sim 10^2$ UI/mL) & \textbf{Très élevé ($\sim 10^3-10^4$ UI/mL)} \\ \hline
        \textbf{Isotype dominant} & IgM transitoire puis IgG (faible affinité) & \textbf{IgG haute affinité immédiate} \\ \hline
        \textbf{Durée maintien} & Court (mois) & \textbf{Long (années, décennies)} \\ \hline
        \end{tabular}
        \end{center}
        \item \textbf{Rôle de la 2\textsuperscript{e} dose (M1) :} C'est une \textbf{dose de primage (priming) essentielle}, pas un simple rappel.
        \begin{itemize}
            \item Elle intervient alors que la réponse à la dose 1 est en phase de montée/plateau (lymphocytes activés, centres germinatifs actifs).
            \item Elle \textbf{entretient la stimulation antigénique} dans les centres germinatifs $\rightarrow$ favorise la \textbf{maturation de l'affinité} (hypermutation somatique + sélection) et la \textbf{commutation de classe} (IgM $\rightarrow$ IgG).
            \item Elle \textbf{amplifie le pool de lymphocytes B mémoire} et de plasmocytes à longue vie générés par la dose 1.
            \item Sans dose 2 : Taux d'anticorps s'effondre vite, affinité faible, peu de mémoire $\rightarrow$ échec vaccinal fréquent. Schéma 2 doses (M0, M6) moins immunogène que 3 doses (M0, M1, M6) pour vaccins sous-unitaires (HBsAg) car pas d'adjuvant vivant/réplicatif.
        \end{itemize}
    \end{enumerate}
    \textbf{Conclusion :} Le schéma 3 doses (0, 1, 6 mois) est optimisé : Dose 1 = amorçage, Dose 2 = maturation affinité/amplification mémoire, Dose 3 = rappel anamnestique puissant assurant protection durable ($>20$ ans souvent vie).
\end{exoresolu}

\begin{attention}[Les 5 erreurs qui coûtent des points au Probatoire]
    \begin{enumerate}
        \item \textbf{Confondre "Antigène" et "Anticorps" (ou "Vaccin donne des anticorps").} 
        \begin{itemize}
            \item \textbf{Erreur :} \og Le vaccin injecte des anticorps pour protéger. \fg / \og L'antigène est produit par le lymphocyte B. \fg
            \item \textbf{Correction :} \textbf{Antigène} = substance étrangère (sur microbe/vaccin) qui \textbf{déclenche} la réponse. \textbf{Anticorps} = protéine (Ig) \textbf{produite par le plasmocyte (LB différencié)} en réponse à l'antigène. \textbf{Vaccin = Antigène(s) atténué/inactivé/ARNm} $\rightarrow$ stimule production propre d'anticorps (Immunité Active). \textbf{Sérum = Anticorps tout faits} (Immunité Passive).
        \end{itemize}
        \item \textbf{Omettre le rôle indispensable des LT4 (Aide lymphocytaire) et du CMHC II dans l'activation des LB (Réponse T-dépendante).}
        \begin{itemize}
            \item \textbf{Erreur :} \og Le lymphocyte B rencontre l'antigène et devient plasmocyte. \fg
            \item \textbf{Correction :} Pour les antigènes protéiques (majorité vaccins, bactéries, virus) : \textbf{2 signaux obligatoires}. Signal 1 : Antigène $\times$ BCR. Signal 2 : \textbf{LT4 effecteur} (reconnaissance peptide-CMHC II sur LB + CD40L/CD40 + IL-2/4/5/6). Sans LT4 $\rightarrow$ Anergie ou réponse T-indépendante faible (IgM seulement, pas de mémoire, pas d'affinité). \textbf{Préciser : CMHC II sur CPA et LB ; CMHC I sur toutes cellules nucléées (pour LT8).}
        \end{itemize}
        \item \textbf{Confondre Immunité Innée et Acquise, ou Active et Passive, dans les définitions et exemples.}
        \begin{itemize}
            \item \textbf{Erreur :} Classer "Allaitement" en Immunité Active. Classer "Vaccin" en Immunité Passive. Dire "Macrophage = immunité acquise".
            \item \textbf{Correction :} \textbf{Tableau de référence (Méthode 3) à connaître par cœur.} Mots-clés : \textbf{Active} = Production propre + Mémoire + Délai. \textbf{Passive} = Transfert Ac/Lymphocytes + Immédiat + Transitoire + \textbf{Pas de mémoire}. \textbf{Innée} = Non-spécifique, pas de mémoire, immédiat. \textbf{Acquise} = Spécifique, mémoire, délai.
        \end{itemize}
        \item \textbf{Négliger le seuil d'immunité collective ($1-1/R_0$) et la notion de "Couverture vaccinale" vs "Efficacité vaccinale" dans la justification de la vaccination.}
        \begin{itemize}
            \item \textbf{Erreur :} \og Si 50\% sont vaccinés, l'épidémie s'arrête. \fg / \og Le vaccin protège à 100\%. \fg
            \item \textbf{Correction :} \textbf{Seuil $S_c = 1 - 1/R_0$}. Rougeole ($R_0=15$) $\rightarrow$ $S_c = 93\%$. Pneumocoque ($R_0 \approx 3$) $\rightarrow$ $S_c = 67\%$. \textbf{Couverture} = \% population ayant reçu le vaccin. \textbf{Efficacité} = \% protection individuelle (ex: 95\% pour VAR). \textbf{Impact collectif} = Couverture $\times$ Efficacité. Il faut Couverture $\times$ Efficacité $> S_c$.
        \end{itemize}
        \item \textbf{Réponses imprécises sur la malnutrition / VIH : Mécanismes vagues ("affaiblit le corps") au lieu de cibles cellulaires/moléculaires précises.}
        \begin{itemize}
            \item \textbf{Erreur :} \og La malnutrition affaiblit le système immunitaire. \fg / \og Le VIH tue les globules blancs. \fg
            \item \textbf{Correction :} \textbf{Malnutrition :} Atrophie thymique $\downarrow$ LT naïfs ; $\downarrow$ Phagocytose/Burst oxydatif ; $\downarrow$ IgA muqueuse ; $\downarrow$ Complément ; Carences Zn/VitA/Fe spécifiques. \textbf{VIH :} Infection et destruction sélective des \textbf{LT4/CD4+} (récepteur CD4) $\rightarrow$ Effondrement aide lymphocytaire $\rightarrow$ Défaillance immunité cellulaire (LT8, Macrophages) et humorale (LB) $\rightarrow$ Infections opportunistes. \textbf{Citer les marqueurs CD4, les stades (OMS), les IO typiques.}
        \end{itemize}
    \end{enumerate}
\end{attention}

\newpage

\coursec{Exercices}

\courssub{Je vérifie mes connaissances}

\begin{exercice}[QCM : les bonnes définitions]
Pour chaque question, entoure la (ou les) bonne(s) réponse(s).
\begin{enumerate}
\item Un \cle{antigène} est :
\begin{enumerate}
\item une cellule du système immunitaire ;
\item une substance étrangère reconnue par l'organisme et capable de déclencher une réaction immunitaire ;
\item une protéine fabriquée par les lymphocytes B ;
\item un microbe toujours pathogène.
\end{enumerate}
\item Un \cle{anticorps} est :
\begin{enumerate}
\item une protéine produite par les lymphocytes B en réponse à un antigène ;
\item une substance capable de neutraliser spécifiquement un antigène ;
\item un organe lymphoïde ;
\item une toxine bactérienne.
\end{enumerate}
\item La \cle{vaccination} consiste à :
\begin{enumerate}
\item injecter des anticorps déjà fabriqués ;
\item introduire dans l'organisme un antigène inoffensif pour provoquer la fabrication d'anticorps et de cellules mémoires ;
\item tuer directement les microbes dans le sang ;
\item soigner une maladie déjà déclarée.
\end{enumerate}
\item L'immunité acquise après une vaccination est :
\begin{enumerate}
\item naturelle et passive ;
\item artificielle et active ;
\item naturelle et active ;
\item artificielle et passive.
\end{enumerate}
\end{enumerate}
\end{exercice}

\begin{exercice}[Vrai ou faux justifié]
Indique si chaque affirmation est vraie ou fausse, puis justifie ta réponse en une ou deux phrases.
\begin{enumerate}
\item Les anticorps produits contre le bacille de la tuberculose protègent aussi contre le virus de la rougeole.
\item Le thymus, la rate et les ganglions lymphatiques sont des organes du système immunitaire.
\item L'allaitement confère au nourrisson une immunité active.
\item Une personne guérie de la varicelle possède des lymphocytes mémoires dirigés contre le virus de la varicelle.
\item Le stress et la malnutrition n'ont aucune influence sur les défenses de l'organisme.
\end{enumerate}
\end{exercice}

\begin{exercice}[Définitions à compléter]
Recopie et complète chaque définition avec les mots ou expressions qui conviennent.
\begin{enumerate}
\item La \textbf{réaction immunitaire} est l'ensemble des moyens mis en \oe uvre par l'organisme pour \ldots\ les \ldots\ qui le pénètrent.
\item Les \textbf{lymphocytes} sont des \ldots\ sanguines qui interviennent dans les réactions immunitaires : les lymphocytes \ldots\ produisent les anticorps, les lymphocytes \ldots\ détruisent les cellules infectées.
\item On dit qu'une immunité est \textbf{passive} lorsque l'organisme reçoit des \ldots\ déjà fabriqués par un autre organisme.
\item Les \textbf{cellules mémoires} permettent une réponse immunitaire plus \ldots\ et plus \ldots\ lors d'un second contact avec le même antigène.
\end{enumerate}
\end{exercice}

\begin{exercice}[Tableau à compléter : les types d'immunité]
Recopie et complète le tableau suivant en indiquant, pour chaque situation, si l'immunité est \textbf{naturelle} ou \textbf{artificielle}, et si elle est \textbf{active} ou \textbf{passive}.
\begin{center}
\begin{tabularx}{\linewidth}{|X|c|c|}
\hline
\rowcolor{popblueL} \textbf{Situation} & \textbf{Naturelle ou artificielle ?} & \textbf{Active ou passive ?} \\
\hline
Allaitement d'un nourrisson (anticorps du lait maternel) & & \\
\hline
Vaccination par le BCG contre la tuberculose & & \\
\hline
Guérison après une varicelle & & \\
\hline
Injection d'une antitoxine (sérum antitétanique) & & \\
\hline
\end{tabularx}
\end{center}
\end{exercice}

\courssub{Je m'entraîne}

\begin{exercice}[Schéma à légender : la réaction immunitaire]
Reproduis et légende le schéma ci-dessous qui résume les quatre étapes d'une réaction immunitaire à médiation humorale, puis rédige une légende explicative de quatre lignes.
\begin{enumerate}
\item Étape 1 : pénétration d'un \ldots\ dans l'organisme ;
\item Étape 2 : reconnaissance de l'antigène par les \ldots\ ;
\item Étape 3 : multiplication des lymphocytes et production d'\ldots\ ;
\item Étape 4 : neutralisation et destruction de l'\ldots\ ; persistance des cellules \ldots .
\end{enumerate}
\end{exercice}

\begin{exercice}[Justifier un fait de la vie courante]
Le petit Moussa, élève au collège de Bafoussam, a été vacciné contre la rougeole dans sa première année de vie, avec un rappel. Un camarade de classe atteint de la rougeole tousse près de lui pendant toute une semaine : Moussa ne tombe pas malade.\\
Explique en cinq lignes maximum pourquoi Moussa n'a pas contracté la maladie, en utilisant obligatoirement les mots : \emph{antigène}, \emph{anticorps}, \emph{cellules mémoires}.
\end{exercice}

\begin{exercice}[Étude de document : primo-injection et rappel]
On injecte à un enfant un vaccin contre le tétanos (primo-injection au jour 0), puis une dose de rappel au jour 30. On mesure le taux d'anticorps antitétaniques dans son sang. Les résultats sont consignés dans le tableau ci-dessous.
\begin{center}
\begin{tabularx}{\linewidth}{|l|X|X|X|X|X|X|X|}
\hline
\rowcolor{popblueL} \textbf{Temps (jours)} & 0 & 10 & 20 & 30 & 40 & 50 & 60 \\
\hline
\textbf{Taux d'anticorps (u.a.)} & 0 & 5 & 15 & 12 & 60 & 45 & 35 \\
\hline
\end{tabularx}
\end{center}
\begin{enumerate}
\item Trace sur papier millimétré la courbe du taux d'anticorps en fonction du temps (1 cm pour 10 jours ; 1 cm pour 10 u.a.).
\item Compare la rapidité et l'amplitude de la réponse après la primo-injection et après le rappel.
\item Explique cette différence en faisant intervenir les cellules mémoires.
\item Conclus sur l'intérêt des doses de rappel dans le calendrier vaccinal.
\end{enumerate}
\end{exercice}

\begin{exercice}[Ordre logique d'une réaction immunitaire]
Voici, dans le désordre, les étapes de la réaction immunitaire qui suit une vaccination contre l'hépatite B :
\begin{enumerate}
\item[a.] production massive d'anticorps spécifiques ;
\item[b.] introduction d'un antigène inoffensif dans l'organisme ;
\item[c.] formation de lymphocytes mémoires ;
\item[d.] reconnaissance de l'antigène par les lymphocytes B ;
\item[e.] destruction rapide du virus lors d'un contact ultérieur réel.
\end{enumerate}
\begin{enumerate}
\item Remets ces étapes dans l'ordre chronologique.
\item Précise à quel type d'immunité (naturelle/artificielle, active/passive) appartient cette protection, en justifiant.
\end{enumerate}
\end{exercice}

\courssub{Je me prépare au Probatoire}

\begin{exercice}[Sujet type Probatoire : la vaccination au Cameroun]
Lors d'une campagne du Programme Élargi de Vaccination (PEV) dans un centre de santé de Yaoundé, des enfants reçoivent le vaccin antirougeoleux. Une mère de famille déclare : « Puisque mon enfant a reçu le vaccin, c'est le vaccin qui tuera les microbes toute sa vie, il n'a plus besoin de rien. »
\begin{enumerate}
\item Définis les termes : \emph{antigène}, \emph{anticorps}, \emph{vaccination}.
\item Cite deux organes lymphoïdes et précise le rôle général du système immunitaire.
\item Explique, en quatre étapes ordonnées, le mécanisme par lequel le vaccin antirougeoleux protège l'enfant.
\item La réaction immunitaire suite à la vaccination est dite \emph{spécifique} et \emph{acquise}. Justifie chacun de ces deux qualificatifs.
\item En t'appuyant sur tes connaissances, corrige en quelques lignes l'erreur contenue dans la déclaration de la mère de famille.
\end{enumerate}
\end{exercice}

\begin{exercice}[Sujet type Probatoire : étude expérimentale de la réponse immunitaire]
Dans un laboratoire, on réalise les expériences suivantes sur trois lots de souris :
\begin{itemize}
\item \textbf{Lot 1} : on injecte une toxine tétanique atténuée (anatoxine) ; 15 jours plus tard, on injecte la toxine active : les souris survivent.
\item \textbf{Lot 2} : on injecte directement la toxine active à des souris n'ayant rien reçu : les souris meurent.
\item \textbf{Lot 3} : on injecte à des souris saines du sérum prélevé sur les survivantes du lot 1, puis la toxine active : les souris survivent.
\end{itemize}
\begin{enumerate}
\item Identifie le rôle joué par l'anatoxine dans le lot 1.
\item Quel est le rôle du lot 2 dans cette expérience ?
\item Interprète les résultats du lot 1 en décrivant les étapes de la réaction immunitaire mise en jeu.
\item Que démontre le lot 3 sur la nature des substances protectrices présentes dans le sérum ?
\item Pour chacun des lots 1 et 3, précise le type d'immunité obtenue (active ou passive, naturelle ou artificielle) et indique laquelle de ces deux immunités est durable, en justifiant.
\end{enumerate}
\end{exercice}

\begin{exercice}[Problème ouvert : renforcer l'immunité d'un élève]
Ariane, élève de Première A à Douala, présente le profil suivant : elle saute régulièrement le petit-déjeuner et mange peu de fruits et de légumes ; elle dort en moyenne 5 heures par nuit ; elle n'est pas à jour de ses vaccinations ; elle vit un stress important à l'approche du Probatoire. Depuis le début de l'année, elle a eu trois angines et un paludisme.
\begin{enumerate}
\item Releve dans ce profil quatre facteurs susceptibles d'affaiblir le système immunitaire d'Ariane.
\item Pour chacun de ces facteurs, explique en une ou deux phrases par quel mécanisme il fragilise les défenses de l'organisme.
\item Propose trois comportements correctifs concrets qu'Ariane pourrait adopter, en justifiant chacun d'eux par son effet bénéfique sur l'immunité.
\item Rédige un court message de sensibilisation (cinq à huit lignes) qu'Ariane, une fois convaincue, pourrait afficher dans sa classe pour encourager ses camarades à protéger leur système immunitaire.
\end{enumerate}
\end{exercice}

\newpage

\coursec{Corrigés des exercices}

\begin{corrige}[Exercice 1 — QCM : les bonnes définitions]
\begin{enumerate}
\item \textbf{Bonne réponse : b.} Un antigène est une substance étrangère (protéine de surface d'un microbe, toxine, molécule d'un vaccin...) reconnue comme \cle{non-soi} par l'organisme et capable de déclencher une \cle{réaction immunitaire}. La réponse a est fausse (un antigène n'est pas une cellule immunitaire), la c est fausse (la protéine fabriquée par les lymphocytes B est l'\cle{anticorps}), la d est fausse (un antigène peut être inoffensif, comme l'\cle{anatoxine} d'un vaccin).
\item \textbf{Bonnes réponses : a et b.} L'anticorps est une protéine (\cle{immunoglobuline}) produite par les \cle{lymphocytes B} en réponse à un antigène, capable de se fixer \emph{spécifiquement} sur cet antigène pour le \cle{neutraliser}. Les réponses c et d sont fausses : un anticorps n'est ni un organe ni une toxine.
\item \textbf{Bonne réponse : b.} La \cle{vaccination} introduit dans l'organisme un antigène inoffensif (microbe tué ou atténué, anatoxine) qui provoque la fabrication d'anticorps et de \cle{cellules mémoires}. La réponse a décrit la \cle{sérothérapie} (immunité passive) ; la vaccination ne tue pas directement les microbes (c) et ne soigne pas une maladie déclarée (d) : elle \emph{prévient}.
\item \textbf{Bonne réponse : b.} L'immunité vaccinale est \textbf{artificielle} (provoquée par l'intervention médicale) et \textbf{active} (c'est l'organisme lui-même qui fabrique ses anticorps et ses cellules mémoires).
\end{enumerate}
\textbf{Point de vigilance :} ne pas confondre vaccination (antigènes injectés, immunité active) et injection de sérum (anticorps injectés, immunité passive).
\end{corrige}

\begin{corrige}[Exercice 2 — Vrai ou faux justifié]
\begin{enumerate}
\item \textbf{Faux.} La réaction immunitaire est \cle{spécifique} : un anticorps ne se fixe que sur l'antigène qui a provoqué sa fabrication. Les anticorps anti-bacille de la tuberculose sont donc inefficaces contre le virus de la rougeole.
\item \textbf{Vrai.} Le \cle{thymus}, la \cle{rate}, les \cle{ganglions lymphatiques} (ainsi que la \cle{moelle osseuse} et les \cle{amygdales}) sont des \cle{organes lymphoïdes}, lieux de formation, de maturation ou de stockage des \cle{lymphocytes}.
\item \textbf{Faux.} L'allaitement transmet au nourrisson des anticorps \emph{déjà fabriqués} par la mère : c'est une immunité \textbf{passive} (et naturelle), temporaire, car le nourrisson ne fabrique pas ses propres anticorps.
\item \textbf{Vrai.} Lors du premier contact avec le virus de la varicelle, l'organisme a produit des \cle{lymphocytes mémoires} qui persistent après la guérison : c'est l'immunité \cle{naturelle active}, durable.
\item \textbf{Faux.} Le \cle{stress} prolongé et la \cle{malnutrition} (carences en protéines, vitamines) diminuent la production et l'efficacité des lymphocytes et des anticorps : les défenses de l'organisme sont affaiblies et les infections deviennent plus fréquentes.
\end{enumerate}
\end{corrige}

\begin{corrige}[Exercice 3 — Définitions à compléter]
\begin{enumerate}
\item La \textbf{réaction immunitaire} est l'ensemble des moyens mis en \oe uvre par l'organisme pour \textbf{éliminer (détruire)} les \textbf{antigènes (corps étrangers)} qui le pénètrent.
\item Les \textbf{lymphocytes} sont des \textbf{cellules} sanguines (globules blancs / leucocytes) qui interviennent dans les réactions immunitaires : les lymphocytes \textbf{B} produisent les anticorps, les lymphocytes \textbf{T} détruisent les cellules infectées.
\item On dit qu'une immunité est \textbf{passive} lorsque l'organisme reçoit des \textbf{anticorps} déjà fabriqués par un autre organisme.
\item Les \textbf{cellules mémoires} permettent une réponse immunitaire plus \textbf{rapide} et plus \textbf{intense (efficace/ample)} lors d'un second contact avec le même antigène.
\end{enumerate}
\end{corrige}

\begin{corrige}[Exercice 4 — Tableau à compléter : les types d'immunité]
\begin{center}
\begin{tabularx}{\linewidth}{|X|c|c|}
\hline
\rowcolor{popblueL} \textbf{Situation} & \textbf{Naturelle ou artificielle ?} & \textbf{Active ou passive ?} \\
\hline
Allaitement d'un nourrisson (anticorps du lait maternel) & Naturelle & Passive \\
\hline
Vaccination par le BCG contre la tuberculose & Artificielle & Active \\
\hline
Guérison après une varicelle & Naturelle & Active \\
\hline
Injection d'une antitoxine (sérum antitétanique) & Artificielle & Passive \\
\hline
\end{tabularx}
\end{center}
\textbf{Justifications :} l'immunité est \emph{naturelle} lorsqu'elle s'acquiert sans intervention médicale (contact réel avec le microbe ou transmission mère-enfant), \emph{artificielle} lorsqu'elle résulte d'un acte médical (vaccin, sérum). Elle est \emph{active} lorsque l'organisme fabrique lui-même ses anticorps (maladie, vaccin), \emph{passive} lorsqu'il reçoit des anticorps tout faits (lait maternel, sérum).
\end{corrige}

\begin{corrige}[Exercice 5 — Schéma à légender : la réaction immunitaire]
Mots à placer :
\begin{enumerate}
\item Étape 1 : pénétration d'un \textbf{antigène} dans l'organisme ;
\item Étape 2 : reconnaissance de l'antigène par les \textbf{lymphocytes B} ;
\item Étape 3 : multiplication des lymphocytes et production d'\textbf{anticorps} ;
\item Étape 4 : neutralisation et destruction de l'\textbf{antigène} ; persistance des cellules \textbf{mémoires}.
\end{enumerate}

\begin{popfigure}
\begin{tikzpicture}[
    node distance=1.5cm and 2cm,
    box/.style={draw=popblue, thick, rounded corners, fill=popblueL, text width=2.8cm, align=center, minimum height=1.3cm, font=\small},
    arrow/.style={-Stealth, thick, popdark}
]
\node[box, align=center] (ag){1. Pénétration\\ de l'antigène};
\node[box, right=of ag, align=center] (rec){2. Reconnaissance\\ par les lymphocytes B};
\node[box, right=of rec, align=center] (prolif){3. Multiplication clonale\\ et différenciation};
\node[box, below=1.2cm of prolif, align=center] (plasma){Lymphocytes B plasmocytes\\ (Usines à anticorps)};
\node[box, above=1.2cm of prolif, align=center] (mem){Lymphocytes B mémoires\\ (Stockage de l'information)};
\node[box, right=of prolif, align=center] (elim){4. Neutralisation /\ Élimination\\ de l'antigène};

\draw[arrow] (ag) -- (rec);
\draw[arrow] (rec) -- (prolif);
\draw[arrow] (prolif) -- (plasma);
\draw[arrow] (prolif) -- (mem);
\draw[arrow] (plasma) -- (elim);
\draw[arrow, dashed, poppurple] (mem) -- ++(1.5,0) |- (elim) node[pos=0.25, above, font=\small, poppurple] {Second contact : réponse rapide};
\end{tikzpicture}
\legende{Schéma simplifié de la réaction immunitaire humorale (à médiation anticorps) lors d'un premier contact. Les cellules mémoires assurent une réponse secondaire plus rapide et plus intense.}
\end{popfigure}

\textbf{Légende explicative (modèle) :} Lorsqu'un antigène pénètre dans l'organisme, il est reconnu comme étranger par des lymphocytes B spécifiques (étape 2). Ceux-ci se multiplient (multiplication clonale) et se différencient en \textbf{plasmocytes} (usines à anticorps) et en \textbf{cellules mémoires} (étape 3). Les anticorps produits se fixent sur l'antigène, le neutralisent et facilitent son élimination par phagocytose (étape 4). Les cellules mémoires persistent et garantissent une réponse plus rapide, plus intense et plus efficace en cas de nouveau contact avec le même antigène.
\end{corrige}

\begin{corrige}[Exercice 6 — Justifier un fait de la vie courante]
\textbf{Modèle de réponse :} La vaccination a introduit dans l'organisme de Moussa un \emph{antigène} inoffensif du virus de la rougeole, qui a provoqué la fabrication d'\emph{anticorps} spécifiques et de \emph{cellules mémoires}. Lorsque le vrai virus, toussé par son camarade, pénètre dans son organisme, les cellules mémoires le reconnaissent immédiatement et déclenchent une production rapide et massive d'anticorps qui détruisent le virus avant l'apparition des symptômes : Moussa ne tombe pas malade.

\textbf{Points de vigilance :} les trois mots imposés doivent apparaître et être correctement employés ; il faut bien dire que c'est le \emph{second contact} (rapide grâce aux cellules mémoires) qui protège, et non le vaccin qui « tue » directement le virus.
\end{corrige}

\begin{corrige}[Exercice 7 — Étude de document : primo-injection et rappel]
\begin{enumerate}
\item \textbf{Courbe représentative :}
\begin{popfigure}
\begin{tikzpicture}
\begin{axis}[
    width=\linewidth, height=8cm,
    xlabel={Temps (jours)}, ylabel={Taux d'anticorps (u.a.)},
    xmin=0, xmax=65, ymin=0, ymax=70,
    xtick={0,10,20,30,40,50,60}, ytick={0,15,30,45,60},
    grid=major, grid style={popline},
    axis lines=middle, enlargelimits=0.05,
    clip=false,
    tick label style={font=\small},
    label style={font=\small}
]
% Réponse primaire
\addplot[popblue, thick, mark=*, mark size=2] coordinates {
    (0,0) (10,5) (20,15) (30,12)
};
% Injection de rappel (marqueur)
\addplot[poporange, only marks, mark=triangle*, mark size=3] coordinates {(30,12)};
% Réponse secondaire
\addplot[poppurple, thick, mark=*, mark size=2] coordinates {
    (30,12) (40,60) (50,45) (60,35)
};
% Annotations
\node[popblue, anchor=south, font=\small] at (axis cs:20,17) {Réponse primaire};
\node[poppurple, anchor=south, font=\small] at (axis cs:40,62) {Réponse secondaire (rappel)};
\node[poporange, anchor=north, rotate=90, font=\small] at (axis cs:30,10) {Injection de rappel (J30)};
\end{axis}
\end{tikzpicture}
\legende{Kinétique du taux d'anticorps sériques après une primo-injection (jour 0) et une injection de rappel (jour 30). La réponse secondaire est plus précoce, plus intense et plus durable.}
\end{popfigure}

\item \textbf{Comparaison :} après la primo-injection, la réponse est \emph{lente} (pic au jour 20, soit 20 jours de latence) et \emph{faible} (maximum 15 u.a.). Après le rappel, la réponse est \emph{rapide} (pic au jour 40, soit 10 jours après l'injection) et \emph{ample} (maximum 60 u.a., soit 4 fois plus : $60/15 = 4$).
\item \textbf{Explication :} la primo-injection a engendré des \cle{lymphocytes mémoires}. Lors du rappel, ces cellules reconnaissent immédiatement l'antigène et déclenchent une \cle{réponse secondaire} plus rapide et plus intense, avec une production massive d'anticorps.
\item \textbf{Conclusion :} les doses de rappel renforcent et prolongent l'immunité en augmentant le taux d'anticorps circulants et le nombre de cellules mémoires ; elles sont donc indispensables dans le calendrier vaccinal (ex. : PEV au Cameroun) pour assurer une protection durable.
\end{enumerate}
\end{corrige}

\begin{corrige}[Exercice 8 — Ordre logique d'une réaction immunitaire]
\begin{enumerate}
\item \textbf{Ordre chronologique :} b $\rightarrow$ d $\rightarrow$ a $\rightarrow$ c $\rightarrow$ e.\\
Introduction de l'antigène inoffensif (b) ; reconnaissance par les lymphocytes B (d) ; production massive d'anticorps spécifiques (a) ; formation de lymphocytes mémoires (c) ; destruction rapide du virus lors d'un contact ultérieur réel (e).
\item \textbf{Type d'immunité :} il s'agit d'une immunité \textbf{artificielle} (provoquée par la vaccination, acte médical) et \textbf{active} (c'est l'organisme vacciné qui fabrique lui-même ses anticorps et ses cellules mémoires). Elle est durable grâce aux cellules mémoires.
\end{enumerate}
\end{corrige}

\begin{corrige}[Exercice 9 — Sujet type Probatoire : la vaccination au Cameroun]
\begin{enumerate}
\item \textbf{Définitions :}
\begin{itemize}
\item \emph{Antigène} : substance étrangère à l'organisme, reconnue comme \cle{non-soi}, capable de déclencher une réaction immunitaire.
\item \emph{Anticorps} : protéine (\cle{immunoglobuline}) produite par les \cle{lymphocytes B} en réponse à un antigène, capable de se fixer \emph{spécifiquement} sur lui pour le \cle{neutraliser}.
\item \emph{Vaccination} : introduction dans l'organisme d'un antigène inoffensif (microbe tué, atténué ou \cle{anatoxine}) afin de provoquer une \cle{immunité active} protectrice.
\end{itemize}
\item \textbf{Organes lymphoïdes} (deux attendus parmi) : le \cle{thymus}, la \cle{rate}, les \cle{ganglions lymphatiques}, la \cle{moelle osseuse}, les \cle{amygdales}. \textbf{Rôle général du système immunitaire :} défendre l'organisme en reconnaissant et en éliminant les agents pathogènes et les substances étrangères (\cle{non-soi}).
\item \textbf{Mécanisme en quatre étapes :} (1) le vaccin introduit un antigène inoffensif du virus de la rougeole ; (2) les lymphocytes B reconnaissent cet antigène (reconnaissance spécifique) ; (3) ils se multiplient (multiplication clonale) et produisent des anticorps spécifiques, tandis que des lymphocytes mémoires se forment ; (4) lors d'un contact réel ultérieur avec le virus, les cellules mémoires déclenchent une réponse secondaire rapide et intense qui détruit le virus avant la déclaration de la maladie.
\item \textbf{Justification des qualificatifs :} la réaction est \emph{spécifique} car les anticorps fabriqués ne reconnaissent que l'antigène de la rougeole et aucun autre (relation clé-serrure anticorps-antigène) ; elle est \emph{acquise} car elle n'existe pas à la naissance (immunité innée) : elle s'est établie après le contact avec l'antigène vaccinal (adaptative).
\item \textbf{Correction de l'erreur :} le vaccin ne tue pas directement les microbes : il « éduque » l'organisme à fabriquer ses propres anticorps et cellules mémoires. De plus, chaque vaccin ne protège que contre une maladie précise (spécificité antigénique) et l'immunité peut s'atténuer avec le temps : des rappels, une bonne hygiène de vie et une alimentation équilibrée restent nécessaires.
\end{enumerate}
\textbf{Barème indicatif (sur 10) :} définitions 3 pts (1 pt chacune) ; organes et rôle 2 pts ; mécanisme en 4 étapes ordonnées 2,5 pts ; justification spécifique/acquise 1,5 pt ; correction de l'erreur 1 pt.\\
\textbf{Points de vigilance Probatoire :} exiger les quatre étapes \emph{ordonnées et numérotées} ; ne pas accepter « le vaccin tue les microbes » ; la spécificité doit être explicitée par la relation anticorps-antigène ; citer au moins deux organes lymphoïdes distincts.
\end{corrige}

\begin{corrige}[Exercice 10 — Sujet type Probatoire : étude expérimentale de la réponse immunitaire (modèle classique Ramon/Glenny)]
\begin{enumerate}
\item \textbf{Rôle de l'anatoxine (lot 1) :} elle joue le rôle d'\textbf{antigène vaccinal} (vaccin) : toxine tétanique rendue inoffensive (par le formol) mais conservant ses propriétés antigéniques, elle déclenche une réaction immunitaire sans provoquer la maladie.
\item \textbf{Rôle du lot 2 :} c'est le \textbf{lot témoin} (témoin de virulence) : il montre que la toxine active est bien mortelle pour des souris non immunisées, ce qui valide l'expérience (sans lui, on ne pourrait pas attribuer la survie du lot 1 à l'immunisation).
\item \textbf{Interprétation du lot 1 :} l'anatoxine est reconnue par les lymphocytes B, qui se multiplient et produisent des anticorps antitétaniques ; des \cle{cellules mémoires} se forment. Quinze jours plus tard (délai nécessaire à la réponse primaire), l'injection de toxine active provoque une \cle{réponse secondaire} rapide et intense : les anticorps neutralisent la toxine avant qu'elle n'atteigne le système nerveux, les souris survivent.
\item \textbf{Ce que démontre le lot 3 :} le sérum des survivantes (lot 1) contient des substances protectrices capables de protéger des souris saines naïves : ce sont des \textbf{anticorps} (substances solubles, transférables par le sérum, spécifiques de la toxine tétanique). C'est la preuve que l'immunité est humorale (à médiation anticorps).
\item \textbf{Types d'immunité :}
\begin{itemize}
\item Lot 1 : immunité \textbf{artificielle active} (l'organisme fabrique lui-même ses anticorps après injection d'antigène). Elle est \emph{durable} car elle s'accompagne de la formation de cellules mémoires.
\item Lot 3 : immunité \textbf{artificielle passive} (injection d'anticorps déjà fabriqués : \cle{sérothérapie}). Elle est \emph{temporaire} (quelques semaines) car les anticorps reçus sont éliminés sans laisser de mémoire immunologique.
\end{itemize}
\end{enumerate}
\textbf{Barème indicatif (sur 10) :} 1 pt par question (questions 3 et 5 pouvant valoir 2 pts chacune selon le barème retenu).\\
\textbf{Points de vigilance Probatoire :} identifier explicitement le lot témoin et son rôle ; dans le lot 3, conclure sur la \emph{nature chimique} (anticorps/protéines) et non seulement sur l'existence d'une protection ; bien distinguer durable (active, mémoire) et temporaire (passive, pas de mémoire) ; utiliser le terme « anatoxine » ou « toxoïde ».
\end{corrige}

\begin{corrige}[Exercice 11 — Problème ouvert : renforcer l'immunité d'un élève]
\begin{enumerate}
\item \textbf{Quatre facteurs affaiblissants identifiés :}
\begin{enumerate}
\item L'alimentation déséquilibrée et carencée (saut du petit-déjeuner, peu de fruits et légumes, apports protéiques insuffisants).
\item Le manque de sommeil chronique (5 h par nuit au lieu de 7--8 h recommandées).
\item Les vaccinations non à jour (absence de rappels : diphtérie, tétanos, coqueluche, hépatite B, fièvre jaune...).
\item Le stress important et prolongé (pression scolaire, anxiété).
\end{enumerate}
\item \textbf{Mécanismes biologiques d'affaiblissement :}
\begin{itemize}
\item \emph{Malnutrition :} les carences en \cle{protéines}, \cle{vitamines} (A, C, D, E) et \cle{sels minéraux} (fer, zinc, sélénium) réduisent la synthèse des anticorps (protéines) et la prolifération/différenciation des lymphocytes (division cellulaire), affaiblissant la réponse humorale et cellulaire.
\item \emph{Manque de sommeil :} le sommeil insuffisant perturbe la sécrétion de cytokines (interleukines) et l'activité des cellules NK (\emph{Natural Killers}) et des lymphocytes T, diminuant l'efficacité de la surveillance immunitaire et de la réponse inflammatoire.
\item \emph{Vaccinations non à jour :} l'absence de rappels laisse diminuer le taux d'anticorps circulants (catabolisme naturel des IgG) et le nombre de cellules mémoires spécifiques : la protection contre les maladies ciblées (tétanos, hépatite B, méningite...) s'affaiblit, risque de réapparition de la maladie.
\item \emph{Stress prolongé :} il provoque la libération chronique de \cle{cortisol} (hormone surrénalienne) qui a un effet \cle{immunosuppresseur} : inhibition de la prolifération lymphocytaire, diminution de la production d'anticorps et de l'activité phagocytaire, augmentant la vulnérabilité aux infections (herpès, rhinites, etc.).
\end{itemize}
\item \textbf{Trois comportements correctifs justifiés :}
\begin{itemize}
\item \textbf{Alimentation équilibrée et variée :} consommer des fruits et légumes à chaque repas (vitamines, antioxydants), des sources de protéines de qualité (poisson, viande, œufs, \textbf{haricots}, \textbf{arachides}, soja -- sources locales accessibles) pour fournir les acides aminés nécessaires à la synthèse des immunoglobulines et au renouvellement cellulaire.
\item \textbf{Sommeil réparateur et gestion du stress :} dormir 7 à 8 heures par nuit ; pratiquer des techniques de relaxation (respiration, sport modéré, organisation du travail) pour limiter le cortisol et restaurer l'efficacité des défenses immunitaires.
\item \textbf{Mise à jour vaccinale :} se rendre au centre de santé (CMA, hôpital de district) pour recevoir les rappels dus (dTcoq, Hépatite B, Fièvre jaune, HPV pour les filles...) afin de reconstituer une immunité active durable via la stimulation des cellules mémoires.
\end{itemize}
\item \textbf{Message de sensibilisation (modèle pour affichage classe / radio scolaire) :}
\begin{quote}
\textbf{« Chers camarades, notre système immunitaire est le soldat qui défend notre corps contre les microbes. Pour le garder fort et prêt au combat :}
\begin{itemize}
\item \textbf{Mangeons équilibré :} des fruits et des légumes à chaque repas, des protéines (poisson, haricots, arachides, œufs) pour fabriquer nos anticorps.
\item \textbf{Dormons suffisamment :} 7 à 8 heures par nuit, le sommeil recharge nos défenses.
\item \textbf{Faisons nos vaccins et rappels :} au centre de santé le plus proche, c'est gratuit et ça protège pour longtemps.
\item \textbf{Gérons notre stress :} organisons nos révisions, prenons l'air, respirons. Un esprit serein dans un corps reposé résiste mieux.
\end{itemize}
\textbf{Protégeons nos défenses, c'est protéger notre avenir et notre réussite au Probatoire ! »}
\end{quote}
\end{enumerate}
\textbf{Points de vigilance Probatoire :} chaque facteur cité doit être relié à un \emph{mécanisme biologique précis} (et pas seulement énuméré) ; les comportements proposés doivent être justifiés par leur effet physiologique sur l'immunité ; le message doit rester scientifiquement correct tout en étant mobilisateur et adapté au contexte camerounais (aliments locaux, structures de santé).
\end{corrige}

\newpage

\coursec{Fiche bilan}

\begin{popfigure}
\centering
\begin{tikzpicture}[mindmap, grow cyclic, every node/.style={concept, text=white, font=\small\bfseries},
  root concept/.append style={concept color=popblue, font=\small\bfseries, minimum size=3.2cm},
  level 1/.append style={level distance=3.6cm, sibling angle=60, minimum size=2.6cm},
  level 2/.append style={level distance=2.4cm, minimum size=1.9cm, font=\scriptsize}]
\node[root concept]{Système immunitaire}
  child[concept color=popgreen]{ node {Organes et cellules}
    child { node {Moelle, thymus, rate, ganglions} }
    child { node {Lymphocytes B et T, phagocytes} }
  }
  child[concept color=poporange]{ node {Antigène et anticorps}
    child { node {AG : marqueur étranger} }
    child { node {AC : protéine spécifique} }
  }
  child[concept color=poppurple]{ node {Réaction immunitaire}
    child { node {Phagocytose} }
    child { node {Réactions B et T} }
  }
  child[concept color=poppink]{ node {Types d'immunité}
    child { node {Innée / acquise} }
    child { node {Naturelle / artificielle} }
  }
  child[concept color=popgold]{ node {Vaccination}
    child { node {Immunité acquise artificielle} }
    child { node {Mémoire immunitaire} }
  }
  child[concept color=popdark]{ node {Facteurs d'affaiblissement}
    child { node {Malnutrition, stress, VIH, paludisme} }
  };
\end{tikzpicture}
\legende{Carte mentale de la leçon : le système immunitaire, ses acteurs, ses réactions, la vaccination et les facteurs qui l'affaiblissent.}
\end{popfigure}

\begin{bilan}
\textbf{Définitions clés}
\begin{itemize}
  \item \cle{Système immunitaire} : ensemble d'organes (moelle osseuse, thymus, rate, ganglions lymphatiques) et de cellules (lymphocytes, phagocytes) qui défendent l'organisme contre les agents pathogènes.
  \item \cle{Antigène (AG)} : substance étrangère à l'organisme, reconnue comme telle, capable de déclencher une réaction immunitaire.
  \item \cle{Anticorps (AC)} : protéine (immunoglobuline) fabriquée par les lymphocytes B, spécifique d'un antigène, capable de le neutraliser.
  \item \cle{Réaction immunitaire} : ensemble des réactions de défense déclenchées par la présence d'un antigène : phagocytose, production d'anticorps, destruction des cellules infectées.
  \item \cle{Vaccination} : introduction dans l'organisme d'un antigène inoffensif (microbe tué ou atténué) pour provoquer une immunité durable sans provoquer la maladie.
\end{itemize}

\textbf{À connaître par cœur}
\begin{center}
\begin{tabularx}{\linewidth}{|l|X|}\hline
\rowcolor{popblueL} \textbf{Notion} & \textbf{À retenir} \\ \hline
Organes lymphoïdes & Moelle osseuse et thymus (primaires) ; rate, ganglions, amygdales (secondaires) \\ \hline
Cellules & Phagocytes (phagocytose) ; lymphocytes B (anticorps) ; lymphocytes T (destruction des cellules infectées) \\ \hline
Spécificité & Un anticorps ne fixe qu'un seul type d'antigène (relation clé--serrure) \\ \hline
Types d'immunité & Innée (à la naissance) ou acquise ; naturelle (après maladie) ou artificielle (après vaccination) \\ \hline
Vaccination & Immunité acquise artificielle : mémoire immunitaire, réponse plus rapide et plus forte au 2\textsuperscript{e} contact \\ \hline
Facteurs d'affaiblissement & Malnutrition, stress, fatigue, maladies (VIH/Sida, paludisme, diabète) \\ \hline
\end{tabularx}
\end{center}

\textbf{Méthodes express}
\begin{itemize}
  \item Expliquer une réaction immunitaire : reconnaissance de l'AG $\rightarrow$ activation des lymphocytes $\rightarrow$ production d'AC ou destruction des cellules infectées $\rightarrow$ élimination de l'AG.
  \item Justifier la vaccination : elle crée des cellules mémoires $\rightarrow$ protection rapide et durable, sans risque de la maladie ; elle protège aussi l'entourage (immunité collective).
  \item Proposer des comportements favorables : alimentation équilibrée, sommeil suffisant, activité physique, gestion du stress, hygiène, respect du calendrier vaccinal, dépistage et traitement des maladies.
\end{itemize}
\end{bilan}

\begin{aretenir}[Checklist avant le Probatoire]
\begin{itemize}
  \item $\square$ Je sais citer les organes du système immunitaire et leur rôle.
  \item $\square$ Je sais nommer les cellules immunitaires (phagocytes, lymphocytes B et T) et leur fonction.
  \item $\square$ Je sais définir un antigène et un anticorps sans les confondre.
  \item $\square$ Je sais expliquer la spécificité de la réaction antigène--anticorps.
  \item $\square$ Je sais décrire les étapes d'une réaction immunitaire (phagocytose, réactions B et T).
  \item $\square$ Je sais distinguer immunité innée et immunité acquise, naturelle et artificielle.
  \item $\square$ Je sais expliquer le principe de la vaccination et le rôle de la mémoire immunitaire.
  \item $\square$ Je sais justifier l'intérêt de la vaccination pour l'individu et pour la communauté.
  \item $\square$ Je sais citer au moins trois facteurs qui affaiblissent le système immunitaire.
  \item $\square$ Je sais proposer des comportements concrets pour renforcer mon immunité.
  \item $\square$ Je sais légender un schéma simple (organe lymphoïde, rencontre AG--AC).
  \item $\square$ Je rédige mes réponses avec le vocabulaire scientifique exact (antigène, anticorps, lymphocyte, phagocytose).
\end{itemize}
\end{aretenir}

\findecours

\end{document}
